%------------------------------------------------------------------------------%
% DIFFERENTIAL TOPOLOGY                                                        %
%------------------------------------------------------------------------------%
% File  : Differential_Topology.tex                                            %
% Author: Klaus Hoeflinger, Beethovenstrasse 11, D-73117 Wangen                %
% Email : khoeflinger@t-online.de                                              %
%------------------------------------------------------------------------------%

%------------------------------------------------------------------------------%
%                       DOCUMENT CLASS and INCLUDE FILES                       %
%------------------------------------------------------------------------------%
\documentclass[11pt,twoside]{article}
\usepackage{geometry}
\geometry{paperheight=241mm,
          paperwidth=152mm,
          top=22mm,
          bottom=17mm,
          left=20mm,
          right=20mm,
          textwidth=112mm,
          textheight=202mm,
          headheight=10mm,
          headsep=3mm,
          footskip=9mm,
          includehead,
          includefoot,
          marginparsep=0mm,
          marginparwidth=0mm}

\renewcommand{\baselinestretch}{0.945}\normalsize

\AtBeginDocument{\setlength\abovedisplayskip{2mm}}
\AtBeginDocument{\setlength\belowdisplayskip{2mm}}

%------------------------------------------------------------------------------%
% define title formats                                                         %
%------------------------------------------------------------------------------%
\usepackage{titlesec}

\renewcommand{\thesection}{\arabic{section}}
\renewcommand{\sfdefault}{FiraSans}
\renewcommand{\ttdefault}{pcr}
\renewcommand{\rmdefault}{antiqua}

\titleformat{\part}[display]
{\normalfont \large \rmfamily \bfseries \centering}
{\small{\thepart.}}{2pt}{}

\titleformat{\section}[runin]
{\normalfont \rmfamily \bfseries}
{\thesection}{1pt}{.\;}[.]

%\titleformat{\section}
%{\normalfont \bfseries \small \filcenter}
%{\thesection.\;}{0mm}{}

\titleformat{\subsection}[runin]
{\normalfont \itshape}
{}{0mm}{}

\titleformat{\subsubsection}[runin]
{\normalfont \itshape}
{}{}{}{}

\titleformat{\paragraph}[runin]
{\normalfont}
{}{}{}{}

\titleformat{\subparagraph}[runin]
{\normalfont}
{}{}{}{}

\titlespacing*{\part}          { 0pt}{ 0pt}{ 8pt}
\titlespacing*{\section}       { 0pt}{ 7pt}{ 4pt}
\titlespacing*{\subsection}    { 0pt}{14pt}{ 6pt}
\titlespacing*{\subsubsection} { 0pt}{ 6pt}{12pt}
\titlespacing*{\paragraph}     { 0pt}{ 6pt}{12pt}
\titlespacing*{\subparagraph}  { 0pt}{ 0pt}{12pt}

%------------------------------------------------------------------------------%
% define enumerate parameters                                                  %
%------------------------------------------------------------------------------%
\usepackage{enumitem}
\setlength{\parindent}{3pt}
\setlength{\parskip}  {0pt}

%\setlist{noitemsep}

\setlist[enumerate]{
         leftmargin=12pt,
         rightmargin=0pt,
         itemsep=1pt,
         itemindent=3pt,
         topsep=3pt,
         labelsep=4pt,
         partopsep=0pt,
         labelwidth=0pt,
         listparindent=0pt,
         parsep=0pt}

\setlist[itemize]{
         leftmargin=12pt,
         rightmargin=0pt,
         itemsep=0pt,
         itemindent=1pt,
         topsep=1pt,
         labelsep=4pt,
         partopsep=1pt,
         labelwidth=0pt,
         listparindent=0pt,
         parsep=0pt}

%------------------------------------------------------------------------------%
% define packages                                                              %
%------------------------------------------------------------------------------%
\usepackage{upref}
\usepackage{amssymb}
\usepackage{slantsc}
\usepackage{amsmath}
\usepackage{amsthm}
\usepackage{thmtools}
\usepackage{amscd}
\usepackage{verbatim}
\usepackage{latexsym}
\usepackage{multicol}
\usepackage{graphicx}
\usepackage{setspace}
\usepackage[ngerman,english]{babel}
\usepackage[protrusion=true,expansion=true]{microtype}
\usepackage{tikz}
\usetikzlibrary{arrows,arrows.meta,decorations.markings}
\usepackage{mathrsfs}
\usepackage{amsfonts}
\usepackage{datetime2}
\usepackage{textgreek}
\usepackage{hyperref}
\usepackage{pifont}
\usepackage{relsize}
%\usepackage{biblatex}
\relsize{-0.05}
\usepackage[skip=0mm]{parskip}
\usetikzlibrary{decorations.pathmorphing}

\usepackage{mlmodern}
\usepackage[T5,T1]{fontenc}
\usepackage{ragged2e}

%\usepackage{mathabx}
%\usepackage{times}

%\usepackage{fontspec}
%\setmainfont{Nimbus Roman No9 L}
%\usepackage[T1]{fontenc}

%------------------------------------------------------------------------------%
% define page layout                                                           %
%------------------------------------------------------------------------------%
\usepackage{fancyhdr}
\pagestyle{fancy}
\setlength{\headheight}{30.0pt}
\renewcommand{\headrulewidth}{0pt}

\AtBeginDocument{
  \addtolength\abovedisplayskip{-0.0\baselineskip}
  \addtolength\belowdisplayskip{-0.0\baselineskip}
}

\fancyhead{}
\fancyfoot{}
\addtolength{\headwidth}{0.02\marginparwidth}

\makeatletter
\let\ps@plain\ps@fancy
\makeatother

\renewcommand\sectionmark[1]
{
 \def\sectionname{#1}
 \markright{\thesection #1}
}

\makeatletter
\@addtoreset{section}{part}
\makeatother

%------------------------------------------------------------------------------%
% define footnote without number                                               %
%------------------------------------------------------------------------------%
\newcommand{\myfootnote}[1]{%
\renewcommand{\thefootnote}{}%
\footnotetext{#1}%
\renewcommand{\thefootnote}{\arabic{footnote}}%
}

%------------------------------------------------------------------------------%
% define numbering in equations                                                %
%------------------------------------------------------------------------------%
\numberwithin{equation}{section}

%------------------------------------------------------------------------------%
% define newtheorem                                                            %
%------------------------------------------------------------------------------%
\declaretheoremstyle[
headfont=\scshape, 
headindent = 0mm, %\parindent,
%postheadhook = {\hspace{0mm}\newline},
%postheadspace = \newline, 
spaceabove = 3mm, 
spacebelow = 3mm,
numberwithin=section,
name=Theorem]{khMATH}
\declaretheorem[style=khMATH]{theorem}

\declaretheoremstyle[
headfont=\bfseries, 
headindent = 0mm, %\parindent,
%postheadhook = {\hspace{0mm}\newline},
%postheadspace = \newline, 
spaceabove = 3mm, 
spacebelow = 3mm,
numberwithin=part,
name=Theorem]{khMATH1}
\declaretheorem[style=khMATH1]{theorem1}

\declaretheoremstyle[
headfont=\scshape, 
headindent = 0mm, %\parindent,
%postheadhook = {\hspace{0mm}\newline},
%postheadspace = \newline, 
spaceabove = 3mm, 
spacebelow = 3mm,
numberwithin=section,
name=Corollary]{khMATH1}
\declaretheorem[style=khMATH1]{corollary}

%            name         counter     label              numeration-mode
\theoremstyle{plain}
\newtheorem*{introduction}            {\sc Introduction}
\newtheorem {standard}                {}                 [section]
\newtheorem*{definition}              {\sc Definition}
\newtheorem{satz1}                    {\sc Satz}
\newtheorem{lemma}                    {\sc Lemma}
\newtheorem*{proposition}             {\sc Proposition}
%\newtheorem{corollary}               {\sc Corollary}
\newtheorem*{corollar}                {\sc Korollar}
\newtheorem*{remark}                  {\sc Remark}
\newtheorem*{remarks}                 {\sc Remarks}
\newtheorem*{example}                 {Example}
\newtheorem*{examples}                {Examples}
\newtheorem*{theo}                    {\sc Theorem}

%------------------------------------------------------------------------------%
% define tabular                                                               %
%------------------------------------------------------------------------------%
\usepackage{tabularx}
\newcolumntype{L}[1]{>{\raggedright\arraybackslash}p{#1}}
\newcolumntype{C}[1]{>{\centering\arraybackslash}  p{#1}} 
\newcolumntype{R}[1]{>{\raggedleft\arraybackslash} p{#1}} 

%------------------------------------------------------------------------------%
% define \sh, \zB                                                              %
%------------------------------------------------------------------------------%
\def\dsh{{d.\,h.}}
\def\ie{{i.\,e.}}
\def\zB{z.\,B.}
\renewcommand{\abstractname}{ }

%------------------------------------------------------------------------------%
% change default spacing for cases                                             %
%------------------------------------------------------------------------------%
\makeatletter
\renewcommand*\env@cases[1][1.6]{%
  \let\@ifnextchar\new@ifnextchar
  \left\lbrace
  \def\arraystretch{#1}%
  \array{@{}l@{\quad}l@{}}%
}
\makeatother

\newenvironment{rcases}
  {\left.\begin{aligned}}
  {\end{aligned}\right\rbrace}

%------------------------------------------------------------------------------%
% change default for \predisplaypenalty                                        %
%------------------------------------------------------------------------------%
\predisplaypenalty=990
\allowdisplaybreaks \numberwithin{equation}{section}

%------------------------------------------------------------------------------%
% general notations                                                            %
%------------------------------------------------------------------------------%
\def\*{\star}
\def\={\,=\,}
\def\+{\,+\,}
\def\xsubset{{\,\subset\,}}
\def\xtimes{{\,\times\,}}
\def\xeof{{\,\in\,}}
\def\eq{{\,=\,}}
\def\xto{{\,\to\,}}
\def\eqdef{\,:=\,}
\def\:{{\,:\,}}
\def\ele{{\,\in\,}}
\def\exists{\mbox{ exists }}
\def\and{\mbox{ and }}
\def\for{\mbox{ for }}
\def\fa{\mbox{ for all }}
\def\fe{\mbox{ for each }}
\def\qfa{\quad \mbox{ for all }}
\def\df{\,\dotfill}
\def\iff{if and only if }
\def\wrt{with respect to}

\makeatletter
\newcommand \Df {\leavevmode \cleaders \hb@xt@ .70em{\hss .\hss }\hfill \kern \z@}
\makeatother

\def\sql{\mathfrak{a}}               % sesquilinear form a
\def\DF{B}                           % Dirichlet form a
%------------------------------------------------------------------------------%
% number sets and special elements                                             %
%------------------------------------------------------------------------------%
\def\P{\mathbb{H}}                   % half plane of $\C$
\def\J{I}                            % interval I
\def\N{{\mathbb{N}}}                 % unsigned integers
\def\Q{{\mathbb{Q}}}                 % rational integers
\def\Z{{\mathbb{Z}}}                 % integers
\def\R{{\mathbb{R}}}                 % real numbers
\def\Rn{{\mathbb{R}}^n}              % real numbers in Rn
\def\Rd{{\mathbb{R}}^d}              % real numbers in Rn
\def\C{{\mathbb{C}}}                 % complex numbers
\def\F{{\mathbb{K}}}                 % arbitrary field
\def\Torus{{\mathbb{T}}}             % complex circle line
\def\T{\tau}                         % upper bound of interval (0,t)
\def\sign#1{\operatorname{sign}(#1)} % sign of a number

%------------------------------------------------------------------------------%
% set theory                                                                   %
%------------------------------------------------------------------------------%
\def\G{G}                            % domain $G$
\def\clG{\overline{\G}}              % closure of $G$
\def\bndG{\partial \G}               % boundary of $G$
\def\setA{\mathcal{A}}               % set A
\def\setB{\mathcal{B}}               % set B
\def\setC{\mathcal{C}}               % set C
\def\setM{M}                         % set M
\def\setN{N}                         % set N
\def\setS{\mathcal{S}}               % set S
\def\famF{\mathfrak{F}}              % family of sets
\def\indI{\mathcal{I}}               % index set I
\def\pot#1{\mathfrak{P}(#1)}         % power set

\def\relR{\mathbf{R}}                % relation R
\def\mapf{f}                         % mapping f
\def\mapg{g}                         % mapping g

\def\set#1#2{\{ \, #1 \,: \, #2 \, \}}
\def\tensor#1#2{#1 \otimes #2}
\def\Tens{\oplus}

\def\cal#1{{\mathcal{#1}}}
\def\aut#1{\operatorname{Aut}(#1)}

\def\cross#1#2{#1 {\,\times\,} #2}
\def\multicross#1#2{#1 \times  \ldots \times  #2}
\def\multitens#1#2 {#1 \otimes \ldots \otimes #2}

%------------------------------------------------------------------------------%
% group theory                                                                 %
%------------------------------------------------------------------------------%
\def\1{{\mathsf{1}}}                % unit element of a monoid
\def\0{{\mathsf{0}}}                % unit element of an Abelian monoid
\def\kernel#1{\mathfrak{N}(#1)}     % kernel of a homomorphism
\def\cokernel#1{\mathfrak{C}{(#1)}} % cokernel of a homomorphism
\def\Hom#1{\operatorname{Hom}(#1)}

%------------------------------------------------------------------------------%
% linear algebra                                                               %
%------------------------------------------------------------------------------%
\def\vecV{\mathit{V}}               % vector space V
\def\vecH{\mathit{H}}               % vector space H
\def\vecW{\mathit{W}}               % vector space W
\def\vecU{\mathit{U}}               % subspace U
\def\convC{\mathcal{C}}             % convex set C
\def\basH{\mathfrak{B}}             % Hamel base B

\def\LO#1{\mathscr{L}(#1) }         % algebra of linear transformations

\def\scalar#1#2{ \langle #1,#2 \rangle }
\def\trace#1{\operatorname{tr}(#1)}
\def\dim#1{{\operatorname{dim} \, #1 }}
\def\codim#1{{\operatorname{codim} \, #1 }}
\def\dim{\operatorname{dim}}
\def\codim{\operatorname{codim}}
\def\ind{\operatorname{ind}}

\def\genG{\cal{G}}

\def\algA{\cal{A}}
\def\algB{\cal{B}}
\def\algU{\cal{U}}

\def\idI{\cal{I}}

%------------------------------------------------------------------------------%
% topology                                                                     %
%------------------------------------------------------------------------------%
\def\compK{{K}}          % compact set C
\def\openO{\mathcal{O}}  % open set O
\def\openU{\mathcal{U}}  % open set U
\def\U{\mathcal{U}}      % neighbourhood U
\def\V{\mathcal{V}}      % neighbourhood V
\def\d{\mathit{d}}       % metric function d

\def\interior#1{\mathring{#1}}
\def\closure#1{{\overline{#1}}}

\def\topT{\mathcal{T}}   % topology T
\def\topX{X}             % topological space X
\def\topY{Y}             % topological space Y
\def\topZ{Z}             % topological space Z
\def\metrS{M}            % metric space M

\def\grpG{G}             % topological group G
\def\grpA{A}             % topological group A
\def\grpB{B}             % topological group B
\def\grpC{C}             % topological group C
\def\grpH{H}             % topological group H
\def\grpU{U}             % topological group U
\def\grpV{V}             % topological group V
\def\topG{G}             % topological group G
\def\topH{H}             % topological group H
\def\topU{U}             % topological subgroup U

\def\subS{\cal{S}}
\def\riR{R}              % ring R

%------------------------------------------------------------------------------%
% calculus                                                                     %
%------------------------------------------------------------------------------%
\def\supp#1{\operatorname{supp}(#1)}
\def\singsupp#1{\operatorname{sing \, supp}(#1)}

\def\re{\operatorname{Re}}
\def\im{\operatorname{Im}}
\def\grad{\operatorname{grad}}

%\def\abs#1 {\left|#1\right|}
\def\abs#1 {\lvert #1 \rvert}
%\def\abs#1 {\left|\,#1\,\right|}

%------------------------------------------------------------------------------%
% measure theory                                                               %
%------------------------------------------------------------------------------%
\def\salgA{\Sigma}               % sigma-algebra S
\def\salgB{\cal{B}}              % sigma-algebra B
\def\Borel#1{\cal{B}_{#1}}       % Borel-algebra 
\def\LB{\lambda}                 % Lebesgue-Borel measure

%------------------------------------------------------------------------------%
% function spaces                                                              %
%------------------------------------------------------------------------------%
\def\Lp{L^p(\G)}                 % spaces L_p
\def\Lq{L^q(\G)}                 % spaces L_q
\def\Lh{L^2(\G)}                 % spaces L_2
\def\Li{L^{\infty}(\G)}          % spaces L_infty
\def\Lc{L_{loc}^1(\G)}           % spaces L_1^{loc}
\def\Ck{C^k(\G)}                 % spaces C_k

\def\BV#1{BV(#1) }               % set of functions of bounded variation
\def\AC#1{AC(#1) }               % set of absolutely continuous functions

\def\MP#1{\mathscr{M}^+(#1) }    % set of positive measures
\def\MS#1{\mathscr{M}^-(#1) }    % vector space of finite signed measures
\def\MC#1{\mathscr{M}(#1) }      % vector space of complex measures
\def\MCR#1{\mathscr{M}_{r}(#1) } % vector space of regular complex measures

\def\SobW{W^{k,p}(\G)}                % Sobolev space W
\def\SobO{\mathring{W}^{k,p}(\G)}     % Sobolev space W
%\def\WN#1#2{\mathring{W}^{#1}(#2) }   % Sobolev space W
\def\WN#1#2{W_0^{#1}(#2) }   % Sobolev space W
%\def\HN#1#2{\mathring{H}^{#1}(#2) }   % Sobolev space H
\def\HN#1#2{H_0^{#1}(#2) }   % Sobolev space H
%------------------------------------------------------------------------------%
% functional analysis                                                          %
%------------------------------------------------------------------------------%
\def\snorm#1{\lVert #1 \rVert}
\newcommand{\norm}[1]{\left\lVert #1 \right\rVert}

\def\topV{V}                       % topological vector space V
\def\topW{W}                       % topological vector space W
\def\normS{X}                      % normed space

\def\banX{\mathit{X}}              % Banach space X
\def\banY{\mathit{Y}}              % Banach space Y
\def\banZ{\mathit{Z}}              % Banach space Z
\def\dbanX{\mathit{X^*}}           % dual space of Banach space X
\def\dbanY{\mathit{Y^*}}           % dual space of Banach space Y
\def\dbanZ{\mathit{Z^*}}           % dual space of Banach space Z

\def\banA{\mathit{A}}              % Banach algebra A
\def\banB{\mathit{B}}              % Banach algebra B
\def\banC{\mathit{C}}              % C*-algebra C


\def\domain#1{\mathfrak{D}(#1)}    % domain
\def\range#1{\mathfrak{R}(#1)}     % range

\def\isoJ{\mathfrak{J}}            % isometry J
\def\repA{{\cal{A}_{\sql}}}        % representation operator A
\def\linS{{S}}                     % linear operator S
\def\linG{{G}}                     % linear operator G
\def\linL{{L}}                     % linear operator L
\def\linT{{T}}                     % linear operator A
\def\linA{\mathit{A}}              % linear operator A
\def\linB{{B}}                     % linear operator B
\def\linM{{M}}                     % linear operator M
\def\linU{{U}}                     % linear operator U
\def\linI{{I}}                     % linear operator I
\def\A{\cal{A}}                    % linear differential operator A
\def\BDO{\cal{B}}                  % linear boundary differential operator B
\def\linU{{U}}                     % unitary operator U
\def\linP{{P}}                     % projection P
\def\compA{k}                      % compact operator k
\def\linFR{f}                      % finite rank operator f
\def\fredA{A}                      % Fredholm operator F

\def\HAM{\cal{H}}                  % Hamilton operator
\def\PA{\partial^{\alpha}}         % elementary differential operator
\def\DO{\cal{A}}                   % linear differential operator
\def\BC{\cal{B}}                   % boundary operator
\def\SLO{\cal{L}_{p,q}}            % Sturm-Liouville operator
\def\MO{\cal{L}_{M}}               % momentum operator

\def\HS#1{S(#1)           }        % algebra of Hilbert-Schmidt operators
\def\FR#1{F(#1)           }        % algebra of finite rank operators
\def\TC#1{N(#1)           }        % algebra of trace class operators
\def\BU#1{\mathscr{U}(#1) }        % algebra of unitary linear operators
\def\BO#1{\mathscr{L}(#1) }        % algebra of bounded linear operators
\def\CO#1{\mathscr{C}(#1) }        % algebra of compact linear operators
\def\CL#1{\mathscr{C}(#1) }        % algebra of compact linear operators
\def\FO#1{\Phi(#1) }               % set of Fredholm operators
\def\FK#1#2{\Phi_{#2}(#1) }        % set of Fredholm operators with index k

\def\hilH{\mathit{H}}              % Hilbert space H
\def\clU{\mathit{U}}               % closed subspace U

\def\orthO{\mathfrak{O}}           % orthonormal system O
\def\orthS{\mathfrak{S}}           % orthonormal system S
\def\basB{\mathfrak{B}}            % basis B of a Hilbert space

\def\GL#1 {\mathbf{GL}(#1)}        % linear group
\def\OG#1 {\mathbf{O}(#1)}         % linear group
\def\SOG#1 {\mathbf{SO}(#1)}       % linear group

%------------------------------------------------------------------------------%
% distribution theory                                                          %
%------------------------------------------------------------------------------%
\def\disF{F}                  % distribution F
\def\disG{G}                  % distribution G
\def\disH{H}                  % distribution H
\def\disU{U}                  % distribution U
\def\disT{T}                  % distribution T
\def\SF{C^{\infty}(\G)}       % space of smooth functions
\def\TF{C_c^{\infty}(\G)}     % space of compactly supported smooth functions
\def\TFRd{C_c^{\infty}(\R^d)} % space of test functions
\def\SRd{\mathscr{S}(\R^d)}   % Schwartz space
\def\SR1{\mathscr{S}(\R)}     % Schwartz space
\def\B1#1{C^1(#1) }           % space of continuously differentiable functions
\def\Bm#1{C^m(#1) }           % space of continuously differentiable functions
\def\Ck{C^k(\G)}
\def\TD{\mathscr{S}'(\R^d)}   % space of temperated distributions
\def\E{\cal{E}}               % space of distributions with compact support
\def\D{\mathscr{D}}           % D
\def\FT{\mathcal{F}}          % Fourier transformation F
\def\FT{\mathscr{F}}          % Fourier transformation F

%------------------------------------------------------------------------------%
% other definitions                                                            %
%------------------------------------------------------------------------------%
\def\Proof{\textsc{Proof}}
\def\FRECHET{Fr\'{e}chet}
\def\ARZELA{Arzel\'{a}}
\def\GARDING{G\r{a}rding}
\def\HOELDER{H\"older}
\def\SCHROEDINGER{Schr\"odinger}
\def\JOERGENS{J\"orgens}
\def\WUEST{W\"uest}
\def\KAEHLER{K\"aehler}
\def\HOEFLINGER{H\"oflinger}
\def\POINCARE{Poincar\'{e}}
\def\FA{Functional Analysis}

\def\Author   {K.\,{\HOEFLINGER}}
\def\AuthorUC {K.\,HOEFLINGER}
\def\Version  {1.0}
\def\Email    {khoeflinger@t-online.de}


%\renewcommand\thesection{\Roman{section}}

\def\Title{Differential Topology}

\titleformat{\part}%[display]
{\normalfont \rmfamily \bfseries \centering}
{{\thepart.}}{3pt}{}

\titleformat{\section}[runin]
{\normalfont \bfseries \filcenter}
{\thesection.\,}{0mm}{}[.]

\titlespacing*{\part}   {5pt}{12pt}{ 4pt}
\titlespacing*{\section}{0pt}{ 3pt}{ 3pt}

%------------------------------------------------------------------------------%
%                                BEGIN OF DOCUMENT                             %
%------------------------------------------------------------------------------%
\begin{document}
\thispagestyle{empty}

\centerline{\Large{\textbf{\Title}}}
\vspace{4pt}
\centerline{\small{{\textsc{\Author}}}}
\vspace{10pt}

\fancyhead[C] {\small{\textsc{\Title}}}
\fancyhead[OL]{\small{\thepage}}
\fancyhead[ER]{\small{\thepage}}

\myfootnote
{
  Version {\Version} from \today;
  Email:  \textit{khoeflinger@\,t-online.de}
}

%------------------------------------------------------------------------------%
%                                DEFINITIONS                                   %
%------------------------------------------------------------------------------%
\def\Off{\mathcal{O}}          % open Umgebung O
\def\OffO{\mathcal{O}}         % open Umgebung O
\def\matrA{A}                  % matrix A

\def\lieG{\mathfrak{G}}        % Lie group G
\def\lieH{\mathfrak{H}}        % Lie group H
\def\lieU{\mathcal{U}}         % Lie subgroup U

\def\GL#1 {\mathbf{GL}(#1)}    % linear group
\def\SL#1 {\mathbf{SL}(#1)}    % special linear group
\def\OG#1 {\mathbf{O}(#1)}     % orthogonal group
\def\SOG#1{\mathbf{SO}(#1)}    % orthonormal group
\def\UG#1 {\mathbf{U}(#1)}     % unitary group U
\def\SUG#1{\mathbf{SU}(#1)}    % special unitary group SU

\def\gl#1{\mathfrak{gl}(#1)}   % Lie algebra of the linear group
\def\sl#1{\mathfrak{sl}(#1)}   % Lie algebra of the special linear group
\def\so#1{\mathfrak{so}(#1)}   % Lie algebra of the orthonormal group
\def\o#1 {\mathfrak{o }(#1)}   % Lie algebra of the orthogonal group
\def\su#1{\mathfrak{su}(#1)}   % Lie algebra of the unitary group SU
\def\u#1 {\mathfrak{u }(#1)}   % Lie algebra of the unitary group U

\def\topM{{M}}        % topological manifold M
\def\topN{{N}}        % topological manifold M

\def\difM{{M}}        % differential manifold M
\def\difN{{N}}        % differential manifold N
\def\difS{{S}}        % differential submanifold S

\def\vecX{X}                   % vector field X
\def\vecY{Y}                   % vector field Y
\def\vecZ{Z}                   % vector field Z

\def\g{{\mathfrak{g}}}         % Lie algebra g
\def\h{{\mathfrak{h}}}         % Lie algebra h
\def\gebG{\mathcal{G}}         % domain G
\def\UmgU{\mathcal{U}}         % neighbourhood U
\def\UmgV{\mathcal{V}}         % neighbourhood V

\def\bracket#1{\left<#1\right>}

\newpage
\thispagestyle{empty}
\fancyhead[C] {\small{\textsc{I. Manifolds}}}
%------------------------------------------------------------------------------%
\part{Manifolds}
%------------------------------------------------------------------------------%
Manifolds form an important class of topological spaces,
which look locally like the Euclidean spaces $\Rd$.
This property enables to transwith many concepts from differential calculus
of several variables to manifolds.

\section{Topological Manifolds}
%------------------------------------------------------------------------------%
A topological space $\topM$ is called \textit{locally Euclidean},
if for each point $p \xeof \topM$ there is an open neighbourhood $\OffO_p$
and a \textit{local} homeomorphism
$\kappa \: \OffO_p \xto \Off' \subseteq \Rd$,
where the dimension $n \eq n(p)$ may depend from the point $p$.

\subparagraph{} %- topological manifolds
One is led to the notation of \textit{topological manifold}
by requiring additional topological features for the space $\topM$.
More precisely,
a topological space $\topM$ is said to be a \textit{topological manifold}
if 
\begin{itemize}[leftmargin=45pt]
\item[(M-1)]
$\topM$ is locally Euclidean;
\item[(M-2)]
$\topM$ is Hausdorff;
\item[(M-3)]
$\topM$ satisfies the second countability axiom.
\end{itemize}

\subparagraph{}
Notice that the Hausdorff property is not a consequence of local Euclidicity,
but must be explicitly required.

\subparagraph{}
It is often the case in monographs that a manifold does not have to fulfill
the second countability axiom.
All manifolds of interest however own this property,
hence it makes sense to regard it in the definition of this concept.

\paragraph{} %- charts
Each local homeomorphism $\kappa \: \OffO \xto \Rd$
is called a $d$-dimensional \textit{chart} for $\topM$,
and the numbers $\kappa(p) \eq \{x_1,\ldots,x_n\}$ are denoted
the \textit{coordinates} of a point $p$ {\wrt} this chart.
The open set $\OffO$ is said to be the \textit{domain of the chart}.
Any chart $(\Off,\kappa)$ is called \textit{centered} {\wrt} $p \xeof \topM$,
if $\kappa(p) \eq (0,\ldots,0) \xeof \Rd$.

\subparagraph{}
If $\kappa_1 \: \Off_1 \xto \R^{d_1}$ and $\kappa_2 \: \Off_2 \xto \R^{d_2}$
are two charts in $\topM$ satisfying
$\Off_1 \cap \Off_2 \neq \emptyset$,
then $n_1 \eq n_2$ holds.
This is a consequence of a theorem quoting 
that the dimensions of two finite-dimensional linear spaces
being homeomorphic to each other are identical.
Hence the mapping $p \mapsto n(p)$ is constant on each connected component
of a topological manifold.
If all connected components of $\topM$ own one and the same dimension $d$,
then $d$ is called the \textit{dimension} of $\topM$.

\subparagraph{}
If the whole manifold $\topM$ is mapped by \textit{one} chart $\kappa$
to the space $\Rd$,
then $(M,\kappa)$ is called a \textit{global chart} for $\topM$.
A collection of $d$-dimensional charts,
whose domains are covering the manifold $\topM$,
is denoted a $d$-dimensional \textit{atlas} for $\topM$.
The validness of the second countability axiom implies
that every topological manifold owns mostly one countable atlas.

\subparagraph{}
Topological manifolds are locally path-connected
due to the local euclidicity.
Hence the concepts of path-connected and connected are equivalent
and the path-connection components of a manifold coincide with their connected components.
Each connected component forms an open set,
an its complementary set is open, too.
Moreover,
the second countability axiom implies
that a topological manifold owns mostly countable connected components.

\begin{comment}
\subparagraph{}
A topological manifold is \textit{semi-locally simple-connected}.
Aus der \"Uber\-lagerungs\-theorie topological spaces is bekannt,
dass sie damit a \textit{universelle \"Uberlagerung} besitzt.
\end{comment}

\subparagraph{}
Every topological manifold $\topM$ is \textit{locally compact}.
Since a locally compact Hausdorff space forms the union
of mostly countable many compact subspaces, 
the set $\topM$ itself proves to be \textit{paracompact}.

\section{Differentiable Manifolds}
%------------------------------------------------------------------------------%
A topological manifold should behave during map changes in a way
that the transition from one map $(\openO_1,\kappa_1)$
to the next one $(\openO_2,\kappa_2)$
within the common domain $\openO_1 \cap \openO_2$ occurs
in a \textit{differentiable} manner.
This leads us to the concept of \textit{differentiable structure}.

\subparagraph{}
Let $\topM$ be a topological manifold and $\cal{A}$ an atlas for $\topM$.
If $(\Off_1,\kappa_1)$ and $(\Off_2,\kappa_2)$ are two charts in $\cal{A}$
with non-empty common domain of the charts,
then on the range $\kappa_1(\Off_1 \cap \Off_2)$ there is defined
a \textit{transformation of coordinates} rsp. a \textit{change of charts}
\[
\tau \eqdef \kappa_2 \circ \kappa_1^{-1}: \kappa_1(\Off_1 \cap \Off_2)
\subseteq \Rd \xto \Rd.
\]
Since $\kappa_1$ and $\kappa_2$ are homeomorphisms,
the mapping $\tau$ is a homeomorphism,
too.
The charts $\kappa_1$ and $\kappa_2$ are called
\textit{differentiable compatible},
if $\tau$ is a diffeomorphism (of class $C^k, k \xeof \N$)
between $\kappa_1(\Off_1 \cap \Off_2)$ and $\kappa_2(\Off_1 \cap \Off_2)$.
Moreover, 
two charts are said to be compatible,
if their domains are disjoint.
The totality of all differentiable compatible charts
is distinguished as follows:
an atlas $\cal{A}$ for a topological manifold $\topM$
is called \textit{differentiable},
if each two charts in $\cal{A}$ are differentiable compatible.

\subparagraph{}
Two differentiable atlases ${\cal{A}}_1$, ${\cal{A}}_2$ for $\topM$
are called \textit{equivalent},
if the union ${\cal{A}}_1 \cup {\cal{A}}_2$,
that is,
the totality of all charts of ${\cal{A}}_1$ and ${\cal{A}}_2$
is a differentiable atlas,
too.
This is the case
iff each chart of the first atlas is differentiable compatible
with all charts of the second atlas and vice versa.
Obviously we have an equivalence relation on the set
of all differentiable atlases of the manifold,
and the union of two atlases in one and the same equivalent class
forms also an atlas in this class.

\subparagraph{}
Thus we obtain for each equivalence class $[{\cal{A}}]$ a \textit{maximal}
rsp. \textit{complete} differentiable atlas ${\cal{A}}_{max}$
by putting together all atlases of the equivalence class,
\[
{\cal{A}}_{max} \eqdef
\bigcup_{{\cal{A}}_i \in [{\cal{A}}]} {\cal{A}}_i.
\]
A complete differentiable atlas $\cal{A}_{max}$ has the property,
that each chart for $\topM$
that is differentiable compatible with all other charts of $\cal{A}_{max}$,
is contained in the atlas itself.
If a differentiable atlas is complete,
then it is denoted a \textit{differentiable structure} on $\topM$.

\paragraph{}
Now we are able to introduce the notion of \textit{differentiable manifold}:

\begin{definition}
Let $\topM$ be a $d$-dimensional topological manifold
and $\cal{A}$ be a differentiable structure on $\topM$.
Then the pair $(\topM,\cal{A})$ is called
a ($d$-dimensional) differentiable manifold.
\end{definition}

\subparagraph{}
Strictly speaking we should call $\topM$ to be a differentiable manifold
of class $C^k$,
that is,
each transformation of coordinates $\tau$ is a diffeomorphism
in $\Rd$ of the class $C^k(\Omega,\tau(\Omega))$ for $\Omega \xsubset \Rd$.
A manifold of the class $C^{\infty}$ is called \textit{analytical}.
When speaking of a differentiable manifold in what follows,
we shall always assume that it is of the class $C^{\infty}$.

\paragraph{}
Examples of differentiable manifolds are:

\begin{itemize}[leftmargin=15pt]
\item[1.]
$\difM := \Rd$ is a differentiable manifold.
The identity $Id \: \Rd \xto \Rd$ is a global chart for $\difM$.
In general,
each finite-dimensional real vector space forms a differentiable manifold.
\item[2.]
$S^d := \{x \xeof \R^{d+1}\: \norm{x} \, \eq 1\}$.
The unit sphere in $\R^{d+1}$ is a standard example for
a $d$-dimensional differentiable manifold.
Eeach differentiable atlas for $S^d$ consists of at least two charts.
\item[3.]
Let $\topM$ be a topological space, endowed with the discrete topology.
Then $\topM$ is a $0$-dimensional differentiable manifold,
and each mapping of the singleton $\{p\}$
to $\{0\}$ for $p \xeof \difM$ forms a chart for $\topM$.
\item[4.]
Let $\openO \xsubset \R^2$ be open and $f \: \openO \xto \R$ be differentiable.
Then the surface
$F := \{(x,y,f(x,y)) \: (x,y) \xeof \openO\} \xsubset \R^3$
is a $2$-dimensional differentiable manifold.
\end{itemize}

\subparagraph{}
An \textit{open} subset $\difM' \xsubset \difM$ of a $d$-dimensional
differentiable manifold $\difM$ is
a differentiable manifold of dimension $d$.
Its topology is the trace topology of $\difM$ in $\difM'$,
and the differentiable structure on $\difM'$ is obtained
by restricition of the charts of $\difM$ on $\difM'$.

\subparagraph{}
Performing the \textit{topological sum} $\difM_1 + \difM_2$
of manifolds $\difM_1$ and $\difM_2$,
the result is again a differentiable manifold,
supposed their dimensions concide.
In this case a differentiable atlas is given by the union of charts
of the differentiable structures of $\difM_1$ and $\difM_2$.
This construction is also possible
for countable many topological sums of manifolds.

\subparagraph{}
Furthermore,
the \textit{cartesian product} $\cross{M_1}{M_2}$ of two
manifolds is a differentiable manifold.
Its topology is the product topology of both manifolds,
and a differentiable atlas is obtained 
by the union of charts of the form
\[
\cross{\kappa_1}{\kappa_2}: \cross{O_1}{O_2} \xto \cross{O'_1}{O'_2},
\quad \kappa_1 \xeof \cal{A}(M_1), \, \kappa_2 \xeof \cal{A}(M_2).
\]
The product-manifold $\cross{M_1}{M_2}$ owns the Dimension
$n_1 + n_2$.

\paragraph{}
Examples of differentiable manifolds are:

\begin{itemize}[leftmargin=15pt]
\item[1.]
$D^n := \{x \xeof \Rd \: \norm{x} \, < 1\}$,
the open ball in $\Rd$,
is a $d$-dimensional differentiable manifold,
when considered as an open set of the Euclidean space $\Rd$.

\item[2.]
Each connected component of a differentiable manifod $\topM$ is,
considered as an open subset of $\topM$,
a differentiable manifold again.
The manifold $\topM$ itself is the topological sum of all connected components.

\item[3.]
The open cylinder $S^1 \xtimes \R$ and the torus $\T^2 := S^1 \xtimes S^1$
are examples of $2$-dimensional product manifolds.
\end{itemize}

\subparagraph{}
The quotient $\difM/\sim$ of a differentiable manifold $\difM$
by a given equivalence relation $\sim$,
however,
does not own a differentiable structure any more.
In fact,
in most cases the quotient topology is not Hausdorff.

\section{Differentiable Mappings}
%------------------------------------------------------------------------------%
Countinuous functions between topological spaces are mappings,
which are compatible with the topological structures of the spaces.
Rather similar,
we are looking for mappings between differentiable manifolds,
that are compatible with their differentiable structures.
The concept of differentiability of mappings between such manifolds
is obtained by pulling back the images and preimages by means of charts
to subsets of the Euclidean space $\Rd$.

\begin{definition}
(differentiable mapping).
Let $\difM,\difN$ be differentiable manifolds.

\begin{itemize}[leftmargin=15pt]
\item[1.]
A continuous mapping $f\: \difM \xto \difN$ is called
\textit{differentiable} in a point $p \xeof \difM$,
if for each chart $(\Off_1,\kappa_1)$ around $p$
and each chart $(\Off_2,\kappa_2)$ around $f(p)$ the mapping
\[
f'(\kappa_1,\kappa_2) := \kappa_2 \circ f \circ k_1^{-1}:
\kappa_1(\Off_1) \xto \kappa_2(\Off_2) \subseteq \Rd
\]
is differentiable.
\item[2.]
The mapping $f$ is called differentiable on $\difM$,
if $f$ is differentiable in each point $p \xeof \difM$.
\end{itemize}
\end{definition}

\paragraph{}
Examples for differentiable mappings are:

\begin{itemize}[leftmargin=15pt]
\item[1.]
The inclusion $Id \: \difM' \xto \difM$ of an open subset
$\difM' \xsubset \difM$ into the manifold $\difM$ is differentiable.
\item[2.]
Especially for $\difN \eq \R$,
a differentiable function $f\: \difM \xto \R$ is called a \textit{scalar field}.
The set of arbitrarily many differentiable scalar fields is usually
denoted by $C^{\infty}(\difM,\R)$.
\end{itemize}

\subparagraph{}
If $\difM_1,\difM_2,\difM_3$ are differentiable manifolds and
$f\: \difM_1 \xto \difM_2$ and $g \: \difM_2 \xto \difM_3$ are differentiable,
then $g \circ f \: \difM_1 \xto \difM_3$ is differentiable,
too.
Differentiable manifolds together with the differentiable mappings between them
form a category.

\subparagraph{}
The isomorphisms between differentiable manifolds are defined as follows:

\begin{definition}
Let $\difM,\difN$ be differentiable manifolds.

\begin{itemize}[leftmargin=15pt]
\item[1.]
A differentiable mapping $f \: \difM \xto \difN$ is called
a \textit{diffeomorphism} between $\difM$ and $\difN$,
if $f$ is bijective and $f^{-1} \: \difN \xto \difM$ is differentiable as well.

\item[2.]
If there is a diffeomorphism between $\difM$ and $\difN$,
then $\difM,\difN$ are said to be \textit{diffeomorphic} to each other.
\end{itemize}
\end{definition}

\subparagraph{}
But notice
that differentiability and bijectivity of $f$ do not imply
that $f^{-1}$ is differentiable.
Consider the example $f \: \R \xto \R$ with $f(x) := x^3$.
The mapping $f$ is merely bijective,
but $f^{-1}$ is not differentiable in $x \eq 0$.

\begin{example}
Each differentiable manifold is diffeomorphic
to the topological sum of its connected components.
\end{example}

\subparagraph{}
Let $\cal{A}_1$ and $\cal{A}_2$ be two differentiable structures on $\difM$.
Then the identity $Id \: \difM \xto  \difM$ is differentiable
iff $\cal{A}_1 \eq \cal{A}_2$. 
Hence the property of differentiability of the identity mapping depends
on the underlying differentiable structures.
For a topological manifold there is in general more than one
differentiable structures (or no such structure).
If $\cal{A}_1$ rsp. $\cal{A}_2$ are two structures,
then the identity mapping $Id:(\difM,\cal{A}_1) \xto (\difM,\cal{A}_2)$
is not necessarily diffeomorphic.
For example,
the sphere $S^7$ has in total 28 different differentiable
(exotic) structures (\textsc{J.\,Milnor} 1963).

\paragraph{}
The notion of \textit{rank} of a differentiable mapping $f$ is based
on the Jacobi-matrix $J_p f$ within an admissable chart:


\begin{definition}
(Rank of a differentiable mapping).

Let $\difM,\difN$ be differentiable manifolds
and $f \: \difM \xto \difN$ be differentiable.
Moreover,
let $(U,\kappa_1)$ and $(V,\kappa_2)$ be charts around $p \xeof \difM$ rsp.
$f(p) \xeof \difN$.
Then the \textit{rank} of $f$ in $p$ is the rank
of the Jacobi-matrix of $\kappa_2 \circ f \circ \kappa_1^{-1}$ in $\kappa_1(p)$.
\end{definition}

\subparagraph{}
The rank of $f$ in $p \xeof \difM$ is independent of the charts
$\kappa_1$, $\kappa_2$ and thus well-defined.
If $f$ has maximal rank in a point $p$,
then half-continuity of the mapping $p \mapsto \operatorname{Rg}_p f$ implies
that $f$ owns maximal rank in a neighbourhood of $p$ and consequently
the inverse functions exists.
This is the quotation of the following theorem,
generalizing the well-known rank theorem:

\begin{theorem} (Implicit function theorem for manifolds).
Let $\difM,\difN$ be two $d$-dimensional differentiable manifolds,
$p \xeof \difM$ and $f \: \difM \xto \difN$ a differentiable mapping
such that
$\operatorname{Rg}_p f \eq n$.
Then is $f$ in $p$ a local diffeomorphism,
that is,
there are open set $\Off_p \xsubset \difM$
and $\Off_{f(p)} \xsubset \difN$
such that $\Off_p$ and $\Off_{f(p)}$ are diffeomorphic by the mapping $f$.
\end{theorem}

\begin{theorem} (rank theorem).
Let $\difM,\difN$ be differentiable manifolds
such that
$\operatorname{dim} \difM \eq r+s$ and
$\operatorname{dim} \difN \eq r+n$.
If a differentiable mapping $f\: \difM \xto \difN$ has constant rank $r$
in a neighbourhood $U$ of a point $p \xeof \difM$,
then there is a chart $\kappa_1$ around $p$
and a chart $\kappa_2$ around $f(p)$
such that
\[
\kappa_2 \circ f \circ \kappa_1^{-1} \:
\R^r \times \R^m \xto \R^r \times \Rd
\]
with $(x,y) \mapsto (x,\0)$ for $x \xeof \R^r$ and $y \xeof \R^m$.
\end{theorem}

\section{Regular and Singular Points}
%------------------------------------------------------------------------------%
\textit{Regular} and \textit{singular} points
{\wrt} a differentiable mapping $f\: \difM \xto \difN$
occur during the classification of the elements in $\difM$ and $\difN$,
which are determined by the values of the rank of $f$ in these points.

\begin{definition}
Let $\difM,\difN$ be differentiable manifolds
and $f\: \difM \xto \difN$ be a differentiable mapping.
\begin{itemize}[leftmargin=15pt]
\item[1.]
A point $p \xeof \difM$ is called a \textit{regular point} of $f$,
if $\operatorname{Rg}_p f \eq dim \, \difN$.

\item[2.]
The non-regular points of $f$ in $\difM$ are called \textit{singular}
or \textit{critical}.

\item[3.]
A point $q \xeof \difN$ is called a \textit{regular value} of $f$,
if the set $f^{-1}(q)$ is empty or consists of regular points only.

\item[4.]
A point $q \xeof \difN$ is called a \textit{singular value} of $f$,
if there is a singular point $p \xeof \difM$ such that $f(p) \eq q$.
\end{itemize}
\end{definition}

\paragraph{}
Example:
let $f\: \R \xto \R$, $x \mapsto x^2{+}1$.
Then $x \eq 0$ is a singular point and each point $x \neq 0$
is a regular point of $f$.
Moreover,
$y \eq 1$ is a singular value and $y > 1$ are regular values of $f$.

\begin{theorem}
(regular point theorem).
Let $\difM,\difN$ be differentiable manifolds with dimensions
$d+s$ rsp. $d$ and $f\: \difM \xto \difN$ be differentiable.
If $p \xeof \difM$ is a regular point of $f$,
then there is a chart $\kappa_1$ around $p$
and a chart $\kappa_2$ around $f(p)$
such that
$\kappa_2 \circ f \circ \kappa_1^{-1}$ is the canonical projection
$\pi:\R^{n+s} \xto \R^s$ ist,
that is,
\[
(x_1,\ldots,x_{n+s}) \mapsto (x_{s+1},\ldots,x_{s+n}).
\]
\end{theorem}

\paragraph{}
Sets of \textit{measure zero} on a manifold can be defined 
by means of those in the manifold $\Rd$.

\begin{definition}
Let $\difM$ be a $d$-dimensional differentiable manifold
and $\A \xsubset \difM$.
Then $\A$ has measure zero,
if for each chart $(\openO,\kappa)$ the range
$\kappa(\A \cap \openO) \xsubset \Rd$ is a set of measure zero.
\end{definition}

\subparagraph{}
This definition does not depend on the choice of $\kappa$,
since all changes of charts are diffeomorphic.

\begin{theorem} (\textsc{Sard}'s theorem).
Let $\difM,\difN$ be differentiable manifolds
and $f \: \difM \xto \difN$ be differentiable.
Then set of critical values in $\difN$ has measure zero.
\end{theorem}

\paragraph{}
As a consequence of this theorem we have:

\begin{theorem} (\textsc{Brown}'s theorem).
The set of regular values of a differentiable mapping $f \: \difM \xto \difN$
is dense in $\difN$.
\end{theorem}

\section{Immersions and Submersions}
%------------------------------------------------------------------------------%
An \textit{immersion} rsp. \textit{submersion} is a differentiable mapping
between manifolds and distinguishes by the fact
that its rank is maximal in each point.

\begin{definition}
Let $\difM,\difN$ be differentiable manifolds.
A differentiable mapping $f \: \difM \xto \difN$ is called

\begin{itemize}[leftmargin=15pt]
\item[1.]
an \textit{immersion},
if $\operatorname{dim} \difM \leq \operatorname{dim} \difN$
and the rank of $f$ in each point $p \xeof \difM$ maximal ist,
that is,
$\operatorname{Rg}_p f \eq \operatorname{dim} \difM$.

\item[2.]
a \textit{submersion},
if $\operatorname{dim} \difM \geq \operatorname{dim} \difN$ and
$\operatorname{Rg}_p f \eq \operatorname{dim} \difN$.
\end{itemize}
\end{definition}

\paragraph{}
Examples for immersions and submersions are:

\begin{itemize}[leftmargin=15pt]
\item[1.]
Each injection $\difM' \xto \difM$ is an immersion.

\item[2.]
The mapping $f\: \R^2 \setminus \{0\} \xto S^1$ with
$f(x,y) := \sqrt{x^2+y^2}^{-1} (x,y)$ is a submersion.
\end{itemize}

\subparagraph{}
A differentiable mapping $f$ is a submersion
iff
each point $p \xeof \difM$ is regular
or each point $q \xeof \difN$ is a regular value.

\subparagraph{}
By the rank theorem,
an immersion has a local representation 
{\wrt} a suitable chart of the form
\[
(x_1,\ldots,x_m) \mapsto (x_1,\ldots,x_m,0,\ldots,0).
\]
However,
a submersion looks locally like a projection
\[
(x_1,\ldots,x_n,0,\ldots,0) \mapsto (x_1,\ldots,x_n).
\]

\section{Submanifolds}
%------------------------------------------------------------------------------%
Given a topological space $\topX$,
each subset $\topX' \xsubset \topX$ is a topological subspace in a natural way,
the injection $\iota \: \topX' \xto \topX$ is continuous,
and $\topX'$ is homeomorphic to the space $\iota(\topX')$.
The mapping $\iota$ is called an \textit{embedding}.

\subparagraph{}
In the differentiable category the definition of the concepts
\textit{submanifold} and \textit{embedding} is a little bit more complicated.

\subparagraph{}
If $\difM,\difN$ are differentiable manifolds with dimensions $d+s$ and $d$,
$f\: \difM \xto \difN$ is differentiable
and $q \xeof \difN$ is a regular value of $f$,
then to each point $p \xeof f^{-1}(q)$ there is a chart $(\openO,\kappa)$
such that
$\kappa(\difM \cap \openO) \eq \R^s \cap \kappa(\openO)$.
Hence the subset $f^{-1}(q) \xsubset \difM$ lies {\wrt} a suitable chart
in $\difM$ like the subset $\R^s$ in the manifold $\R^{d+s}$.
Dies gibt Anla\"ss to folgender definition.

\paragraph{}
The concept of \textit{submanifold} is defined as follows:

\begin{definition}
Let $\difM$ be a $d$-dimensional differentiable manifold.
A subset $\difM' \xsubset \difM$ is called
a \textit{$k$-dimensional submanifold} of $\difM$,
if there is a chart $(\openO,\kappa)$ around each point $p$ of $\difM'$
such that
\[
\kappa(\difM' \cap \openO) \eq \R^k \cap \kappa(\openO).
\]
The number $d-k$ is called the \textit{codimension} of $\difM'$ in $\difM$.
\end{definition}

\paragraph{}
Examples of submanifolds are:

\begin{itemize}[leftmargin=15pt]
\item[1.]
The $0$-dimensional submanifolds are exactly the discrete subsets of $\difM$.

\item[2.]
The $d$-dimensional submanifolds of a $d$-dimensional
manifold $\difM$ are exactly the open subsets of $\difM$.
\end{itemize}

\subparagraph{}
The \textit{regular value theorem} enables to construct submanifolds
in a simple way:

\begin{theorem} (Regular value theorem).
Let $\difM,\difN$ be differentiable manifolds
and $f\: \difM \xto \difN$ a differentiable mapping.
If $q \xeof \difN$ is a regular value of $f$,
then $f^{-1}(q)$ is a submanifold of $\difM$,
whose codimension is equal to the dimension of $\difN$.
\end{theorem}

\begin{itemize}[leftmargin=15pt]
\item[1.]
If $\difM$ is a $d$-dimensional differentiable manifold
and $f \: \difM \xto \R$ is differentiable,
then for each $c \xeof \R$ the set $f^{-1}(c)$
is a $(d-1)$-dimensional submanifold of $\difM$.

\subparagraph{}
Especially for $\difM \eq \R^2$ and $f \: \R^2 \xto \R$
with $f(x,y) := x^2 {+} y^2{-}1$
we have $f^{-1}(0) \eq S^1$.
Hence the circle $S^1$ is a $1$-dimensional submanifold of $\R^2$.

\item[2.]
Let $\difM$ be the set of $(n,n)$-matrices with real coefficients
and $\difN$ be the subset of \textit{symmetrc} $(n,n)$-matrices.
Then the mapping
$f\: \difM \xto \difN$ with $\matrA \mapsto \matrA \circ \matrA^T$
is differentiable.
The unit matrix $I \xeof \difN$ is a regular value of $f$,
and the preimage $f^{-1}(I)$ is the set $\OG{n}$ of orthogonal $(n,n)$-matrices.
Due to the last theorem,
$\OG{n}$ is a submanifold of $\difM$,
whose dimension is $n(n-1)/2$,
and the set $\SOG{n} \xsubset \OG{n}$ is a open submanifold of $\OG{n}$.
\end{itemize}

\section{Embeddings}
%------------------------------------------------------------------------------%
In the topological category,
an embedding of a topological space
$\topX$ into a further topological space $\topY$ is a homeomorphism
between $\topX$ and a subset $\topY' \xsubset \topY$.

\subparagraph{}
In the differentiable category,
however,
one defines embeddings as follows:

\begin{definition}
Let $\difM,\difN$ be differentiable manifolds and
$f\: \difM \xto \difN$ a differentiable mapping.
Then $f$ is called an \textit{embedding} of $\difM$ in $\difN$ if

\begin{itemize}[leftmargin=15pt]
\item[1.]
$f(\difM)$ is a submanifold of $\difN$;
\item[2.]
$f$ is a diffeomorphism between $\difM$ and $f(\difM)$.
\end{itemize}
\end{definition}

\subparagraph{}
Hence $\difM$ is embedded into $\difN$,
if $\difM$ is diffeomorphic to a submanifold of $\difN$.
For example,
the mapping $f \: S^1 \xto \R^2$ with $t \mapsto (\cos(t), \sin(t))$
is an embedding of $S^1$ into $\R^2$.

\paragraph{}
\textsc{Whitney}\textit{'s embedding theorem} claims
that each differentiable manifold can be embedded
into the Euclidean space $\Rd$,
supppsed the dimension is chosen high enough:

\begin{theorem}
Each $d$-dimensional differentiable manifold $\difM$ can be embedded
into the manifold $\R^{2d+1}$.
If $f$ is such an embedding,
then $f(\difM) \subseteq \R^{2d+1}$ is even closed.
\end{theorem}

\subparagraph{}
But notice that
there is no canonical embedding of a manifold into one of the spaces $\Rd$.

\section{The Tangential Space}
%------------------------------------------------------------------------------%
When speaking of a direction in a point of a differentiable manifold,
we have to provide vectors,
which are attached to this point.
This means
that to each point there is to assign a vector space,
whose elements represent all directions of the manifold.

\subparagraph{}
This leads to the concept of the \textit{tangent vector}
in a differentiable manifold.
Its definition should be beared from the manifold itself,
that is,
without using any larger manifold, in which the first one is embedded.

\subparagraph{}
Let $\difM$ be a differentiable manifold,
$p \xeof \difM$ a point and $F_p(M)$ be the set of
all differentiable scalar-valued functions
in a neighbourhood von $p$.
Then on $F_p(M)$ there is defined an equivalence relation according to
\[
f_1 \sim f_2, \quad \mbox{if} \quad
f_1(q) \eq f_2(q) \quad \mbox{ for all } q \xeof U,
\]
where $U$ is a sufficiently small neighbourhood of $p$,
on which $f_1$ and $f_2$ are well-defined.

\subparagraph{}
Let $F_p(M)/\sim$ be the set of these equivalence classes.
They are called the \textit{Funktionskeime} in $p \xeof \difM$.
The set $F_p(M)/\sim$ owns the algebraic structure of a $\R$-algebra.

\begin{definition}
A \textit{tangent vector} in a point $p \xeof \difM$
is an algebra derivative $v \: F_p(M)/\sim \, \xto \R$.
\end{definition}

\subparagraph{}
A \textit{geometric} definition of the tangent vector can be done
by means of differentiable curves.
A \textit{differentiable curve} $\gamma:I \subseteq \R \xto \difM$ is a mapping
such that
for all charts $(U,\kappa) \xeof \cal{A}$
the mappings $\kappa \circ \gamma: I \xto \R$ are differentiable.
A \textit{tangent vector} in $p \xeof \difM$ is an equivalence class
of differentiable curves $\gamma:(-\epsilon,\epsilon) \xto \difM$ fulfilling
$\gamma(0) \eq p$,
where
$\gamma_1 \sim \gamma_2 \Leftrightarrow (\kappa \circ \gamma_1)'(0) =
(\kappa \circ \gamma_2)'(0)$
for each chart $(\Off,\kappa)$ with $p \xeof \Off$.
This definition is equivalent to the algebraic definition.

\subparagraph{}
The derivation $v$ is also linear on $F_p(M)/\sim$
and satisfies the product rule
$v(f \cdot g) \eq v(f) \cdot g(p) + f(p) \cdot v(g)$
for $[f], [g] \xeof F_p(M)/\sim$.

\subparagraph{}
For the tangent vectors of a point $p$
we can define in a usual way the addition and multiplication with numbers
such that
the set of all tangent vectors in $p$ forms a $\R$-vector space.

\begin{definition}
(tangential space, cotangential space).
Let $\difM$ be a differentiable manifold and $p \xeof \difM$.
The vector space $T_p \difM$ of tangent vectors in $p$ is called
the \textit{tangential space} of $\difM$ in $p$.

\subparagraph{}
The dual space $T_p^{\*} \difM$ of the tangential space in $p$ is called
the \textit{cotangential space} of $\difM$ in $p$.
\end{definition}

\subparagraph{}
Let $\difM$ be a differentiable manifold and $p \xeof \difM$.
Each chart $(\Off,\kappa)$ with $p \xeof \Off$ induces
a base of the tangential space $T_p \difM$.

\subparagraph{}
For any $i \xeof \{1,\ldots,n\}$
and any function $f \xeof F_p(M)$ the mapping
$f \circ \kappa^{-1}:\Off \xsubset \Rd \xto \R$ and the mapping
\[
\frac{\partial}{\partial x_i}(f) :=
\frac{\partial (f\circ \kappa^{-1})}{\partial x^i} \mid_{\kappa(p)}
\]
is real-valued,
which assigns to $f$ the derivative {\wrt} the $i$-ten coordinate
of $f \circ \kappa^{-1}$.
The mapping $\frac{\partial}{\partial x_i}$ owns all properties
of a tangent vector.
Let $B_p$ be the set of tangent vectors defined on this way.
Then enth\"alt $B_p$ genau $d$ linearly  independent vectors and
and thus forms a base of the tangential space $T_p \difM$.
It is called the \textit{induced base} in $p$ of the chart $(\Off,\kappa)$
The vectors of these bases are denoted as $\partial_i$.

\subparagraph{}
If a differentiable manifold $\difM$ is $d$-dimensional,
then the related tangential space $T_p \difM$ in each point $p \xeof \difM$
is a $d$-dimensional vector space.

\begin{definition}
A base $B$ of the tangential space $T_p \difM$ in a point $p \xeof \difM$
is called \textit{natural},
if there is a chart $(\Off,\kappa)$ with $p \xeof \Off$
such that
$B$ is the induced base of $(\Off,\kappa)$.
\end{definition}

\subparagraph{}
If a base $B$ of $T_p \difM$ is natural,
let $dx^i, \ldots, dx^n$ be the covectors of the dual base of $B$
in the cotangential space.
They assign to each tangent vector $v \xeof T_p \difM$ its $i$-te coordinate
{\wrt} the base vector $\partial_i$,
\[
dx^i(v) \eq dx^i(v^j \partial_j) \eq v^j \, dx^i(\partial_j) \eq v^i,
\qquad (i=1,\ldots,n).
\]
The covectors $dx^i$ are also called \textit{coordinate differentials}.

\subparagraph{}
Putting together all tangent vectors,
which appear in the points of a differentiable manifold,
we obtain a \textit{bundle} of vectors:

\begin{definition}
Let $\difM$ be a differentiable manifold
and $T\,\difM$ be the disjoint union of all tangent vectors,
\[
T\,\difM  \eq \bigsqcup_{p \xeof \difM} \cross{\{p\}}{T_p \difM}.
\]
Moreover, let
$\pi: T \, \difM \xto \difM$ with $v \xeof T_p \difM \xsubset T \, \difM \mapsto p$
be the mapping,
which assigns to each tangent vector its base point $p$.
Then the triple $(T\,\difM, \pi, \difM)$ is called
the \textit{tangent bundle} of $\difM$.
\end{definition}

\subparagraph{}
A topology $\topT$ on the tangent bundle is constructed
by means of the initial topology.
This is the coarsest topology on $T\,\difM$
such that the projection $\pi$ is still continuous.
Endowed with this topology,
$T\,\difM$ is locally Euclidean to the space $\R^{2d}$,
fullfills \textsc{Hausdorff}'s separation axiom
and the second countability axiom.
Eeach tangent bundle is in a natural way
a $2d$-dimensional topological manifold.

\subparagraph{}
Moreover,
there is a canonical differentiable structure on $T\,\difM$
inherited from the base manifold $\difM$.
The bundle $T\,\difM$ is even a differentiable manifold.

\section{The Differential of a Differentiable Map}
%------------------------------------------------------------------------------%
Given a point $p \xeof \difM$,
a differentiable mapping $f$ induces
in a natural way a vector space homomorphism
between den tangential spaces $T_p$ and $T_{f(p)}$.

\subparagraph{}
Let $\difM,\difN$ be differentiable manifolds,
$f \: \difM \xto \difN$ differentiable and $p \xeof \difM$.
If $v \xeof T_p \difM$,
then for each scalar field $\omega \xeof C^{\infty}(N,\R)$
the mapping $v'(\omega) \eq v(\omega \circ f)$
forms a tangent vector in $T_{f(p)}(N)$.
This leads to the following definition:

\begin{definition}
The \textit{differential} of a differentiable mapping
$f \: \difM \xto \difN$ in a point $p \xeof \difM$
is the {linear mapping}
\[
\operatorname{d}f_p : T_p \difM \to T_{f(p)} N \quad \mbox{with}
\quad \operatorname{d}f(v) \eq v', \quad v \xeof T_p \difM.
\]
\end{definition}

\subparagraph{}
Taking the differentials is a functor
from the category of punctured differentiable manifolds
to the category of $\R$-vector spaces.

\paragraph{}
Consider the following example:
for $\difM$ arbitrary and $\difN \eq \R$
we obtain as the differential of a scalar field $f$
in $p \xeof \difM$ the covector $\operatorname{d}f \xeof T_p^\* \difM$,
that is,
a linear functional on the space of tangent vectors in $p$,
which assigns to each $v \xeof T_p \difM$ the directional derivative of $f$
according to
\[
\operatorname{d}f(v) \eqdef v(f) \xeof \R \qquad (v \xeof T_p \difM).
\]

\section{Transversality}
%------------------------------------------------------------------------------%
The concept of \textit{transversality} describes the situation
that two submanifolds of a given manifold are as orthogonal to each other
as possible.

\begin{definition}
Let $\difM$ be a differentiable manifold.
Two submanifolds $\difM_1,\difM_2$ of $\difM$ are called \textit{transversal}
to each another,
if for $p \xeof \difM_1 \cap \difM_2$
$T_p\,\difM_1 + T_p\,\difM_2 \, \eq \, T_p\,\difM$ holds.
\end{definition}

\subparagraph{}
Hence in each point $p$ the tangential spaces of the submanifolds span
the tangential space in $p$.

\begin{theorem}
Let $\difM$ be a differentiable manifold.
If two submanifolds $\difM_1, \difM_2 \subseteq \difM$ are transversal
to each another,
then $\difM_1 \cap \difM_2$ is a submanifold of $\difM$.
\end{theorem}

\subparagraph{}
The concept of transversal submanifolds is defined as follows:

\begin{definition}
Let $\difM,\difN$ be differentiable manifolds
and $f \: \difM \xto \difN$ be a differentiable mapping.
Then $f$ is called \textit{transversal}
to a submanifold $\difS \subseteq \difM$,
if for all $p \xeof f^{-1}(S)$
\[
T_{f(p)}\,\difN \= T_{f(p)}\,\difS \, + \, df(T_p\,\difM).
\]
\end{definition}

\subparagraph{}
Let us consider the following examples:

\begin{itemize}[leftmargin=15pt]
\item[1.]
A mapping $f$ is transversal to $S$,
if $f(\difM) \cap S \eq \emptyset$.
\item[2.]
Each mapping $f\: \difM \xto \difN$ is transveral to $\difN$.
\item[3.]
If $S \eq \{q\}$ is a singleton,
then $f$ is transversal to $\{q\}$ iff $q$ is a regular point of $f$.
\end{itemize}

\begin{theorem}
Let $f\: \difM \xto \difN$ be a differentiable mapping,
which is transveral to the submanifold $\difS \subseteq \difN$.
Then $f^{-1}(\difS)$ is a submanifold of $\difM$ and
\[
\operatorname{dim} f^{-1}(S) \= \operatorname{dim} \difM
                              - \operatorname{dim} \difN
                              + \operatorname{dim} \difS.
\]
\end{theorem}

\newpage
\thispagestyle{empty}
\fancyhead[C] {\small{\textsc{II. Tensors and Lie-Algebras}}}
%------------------------------------------------------------------------------%
\part{Tensors and Lie-Algebras}                                                %
%------------------------------------------------------------------------------%
Tensors are the main objects in multilinear algebra.
The concept of \textit{tensor} unifies notations like \textit{vector},
\textit{linear form}, \textit{linear mapping}, \textit{bilinear form},
\textit{multilinear form} and many others by a general formalism.
Tensors are important objects in differential geometry
and theoretical physics (relativity theory, classical field theory).

\section{Generalities on Tensors}
%------------------------------------------------------------------------------%
In the following we consider the multiple tensor product of a $d$-dimensional
real vector space $V$ and its dual space $V^\*$.
Let $r$ factors $V$ and $s$ factors $V^\*$ be given,
then we obtain the $\R$-vector space $T_{r,s}(V)$ by repeated performing
the tensor product according to
\[
T_{r,s}(V) \eqdef
\underbrace{{\multitens{V}   {V}   }}_r
\otimes
\underbrace{{\multitens{V^\*} {V^\*} }}_s.
\]
Especially for $r \eq s \eq 0$ let $T_{0,0}(V) := \R$.

\paragraph{}
The sets $T_{r,s}(V)$ constitute finite-dimensional vector spaces,
too,
having the dimensions $n^{r+s}$.
The subsequent definitions are fundamental:

\begin{definition}
Let $r,s \xeof \N$ and $V$ be a $n$-dimensional real vector space.

\begin{itemize}[leftmargin=15pt]
\item[1.]
A vector $T \xeof T_{r,s}(V)$ is called a ($r$,$s$)-tensor on $V$ 
of \textit{level} $r+s$.
\item[2.]
The tensor $T \xeof T_{r,s}(V)$ is called $s$\textit{-fold covariant} and 
$r$\textit{-fold contravariant}.
Especially for $r \eq 0$, 
$T$ is a \textit{pure covariant} tensor of level $s$, 
and for $s \eq 0$
the vector $T$ is a \textit{pure contravariant} tensor of level $r$.
\end{itemize}
\end{definition}

\paragraph{}
For finite tensor products of vector spaces $V_1,\ldots,V_n$
there is an isomorphism
\[
(V_1 \otimes \ldots \otimes V_n)^\* \cong 
V_1^\* \otimes \ldots V_n^\*.
\]
This isomorphism yields for the spaces $T_{r,s}(V)$
by usage of $V \cong V^{\*\*}$ and 
$Hom_{\R}(\tensor{V_1}{V_2},\R) \cong Hom_{\R}(V_1 \times V_2,\R)$
the isomorphism
\[
T_{r,s}(V) \cong
(\underbrace{V^\* \times \ldots \times V^\*}_r
 \times
 \underbrace{V    \times \ldots \times V   }_s)^\*,
\]
that is,
the elements of $T_{r,s}(V)$ can be considered as real-valued multilinear forms
on finite cartesian products of $V^\*$ and $V$.

\subparagraph{}
for a finite-dimensional vector space and smal pairs of values $(r,s)$
the $(r,s)$-tensors can be described explicitely.
Consider the following examples:

\begin{itemize}[leftmargin=15pt]
\item[1.]
$(r=1, s=0)$.\\
A simple contravariant tensor of stage $(1,0)$
is a linear mapping $T \: V^\* \xto \R$
and hence an element of the bidual space $V^{\*\*}$. 
Since $V$ may be identified with $V^{\*\*}$, 
the $(1,0)$-tensors are the elements of the vector space $V$ itself.

\item[2.]
$(r=0, s=1)$.\\
A covariant tensor $T$ of stage $(0,1)$ is a linear mapping $T\:V \xto \R$
and hence a covector rsp. a linear functional.

\item[3.]
$(r=1, s=1)$.\\
Let $T$ be a (1,1)-tensor with $T(f,v) \xeof \R$ for $f \xeof V^\*$
and $v \xeof V$. 
Varying for a fixed vector $v$ the argument $f$ 
along $V^\*$, the mapping
\[
T_v(f) \: V^\* \xto \R, \qquad T_v(f) := T(f,v)
\]
defines a contravariant tensor $T_v$. 
By identification of $T_v$ with a vector $w \xeof V$, 
we obtain from $T$ a linear mapping
\[
A_T \: V \xto V, \qquad A_T(v) \eq w \quad \for v \xeof V.
\]
Hence each (1,1)-tensor$T$ is associated with a linear mapping
$A_T \: V \xto V$ this way.

\subparagraph{}
The (1,1)-tensor $\bracket{f,v}$ is of special interest.
It assigns to each pair $(f,v) \xeof \cross{V^\*}{V}$ the real number $f(v)$.
The mapping $\bracket{,}$ is called the \textit{inner product}
in the space $\cross{V^\*}{V}$.
The associated operator $A_{\bracket{,}}$ is the identity mapping on $V$.

\item[4.]
$(r=0, s=2)$.\\
A two-fold covariant tensor on $V$ is nothing but a bilinear form on $V$. 
\end{itemize}

\paragraph{}
In the following the tensors $T \xeof T_{r,s}(V)$ are regarded as real-valued
multilinear mappings.
Let $T$ a ($r_1$,$s_1$)-tensor and $U$ a ($r_2$,$s_2$)-tensor. 
Then the \textit{product} $T \otimes U$ is the multilinear mapping given by
\[
(T \otimes U)(
   f_1,\ldots,f_{r1},f_{r1+1},\ldots,f_{r1+r2},
   v_1,\ldots,v_{s1},v_{s1+1},\ldots,v_{s1+s2}) :=
\]
\[
  T(f_1,\ldots,f_{r1},v_1,\ldots,v_{s1}) 
\cdot
  U(f_1,\ldots,f_{r2},v_1,\ldots,v_{s2}),
\]
hence $T \otimes U$ is a $(r_1+r_2,s_1+s_2)$-tensor.

\subparagraph{}
For the product of two tensors, which are sums and products of tensors,
the distributive laws and the associative law hold:
\[
( T_1 + T_2 ) \otimes U \eq T_1 \otimes U + T_2 \otimes U,
\]
\[
U \otimes ( T_1 + T_2 ) \eq U \otimes T_1 + U \otimes T_2,
\]
\[
( \lambda \cdot T ) \otimes U \eq \lambda \cdot ( T \otimes U ), \qquad
\lambda \xeof \R,
\]
\[
T \otimes ( U \otimes V ) \eq ( T \otimes U ) \otimes V.
\]

\paragraph{}
Let $V$ be a $n$-dimensional vector space.
A base of the vector space $T_{r,s}(V)$ can be setup from a base of the
vector space $V$ and its cobase in $V^\*$. 
Let $\{e_1,\ldots,e_n\}$ be base vectors in $V$ and 
$\{\epsilon^1,\ldots,\epsilon^n\}$ be base covectors in $V^\*$. 
Then for a ($r$,$s$)-tensor$T$ the following representation holds:
\[
T \eq \sum_{1 \leq j,k \leq n} 
T_{j_1,\ldots,j_r}^{k_1,\ldots,k_s} \cdot 
\epsilon^{j_1} \otimes \ldots \otimes \epsilon^{j_r} \otimes
e_{k_1} \otimes \ldots \otimes e_{k_s},
\]

\begin{definition}
The coefficients $T_{j_1,\ldots,j_r}^{k_1,\ldots,k_s} \xeof \R$ are called
the \textit{coordinates} of the tensor {\wrt} the base vectors 
$e_i \xeof V$ and $\epsilon^j \xeof V^\*$.
\end{definition}

\subparagraph{}
This representation implies
that one may obtain base tensors by repeated performing the product
of base vectors $e_i$ and base covectors $\epsilon^j$.

\paragraph{}
In multilinear algebra it is convenient
to formulate equations in a way
that for a product of tensors we have to sum over those indices,
which appear in any factor as uppercase as well as
in a further factor as lowercase \textit{(Einstein's sum convention)}.
Consider the following examples:

\begin{itemize}[leftmargin=15pt]
\item[1.]
Let $V$ be a $n$-dimensional vector space, $e_i$ a base of $V$ and 
$\epsilon^j$ the dual base in $V^\*$.
Then we briefly write $v \eq v^i \, e_i$ for $v \xeof V$ and
$f \eq f_i \, \epsilon^i$ for $f \xeof V^\*$.

\item[2.]
Let $b$ a bilinear form in $V$.
We briefly write 
$b(v,w) \eq v^i \, w^j \cdot b_{ij}$ with $b_{ij} := b(e_i,e_j)$.
\end{itemize}

\paragraph{}
Performing the \textit{trace} of a given tensor is an operation,
which assigns to a $(r,s)$-tensor with $r,s \geq 1$
a $(r-1,s-1)$-tensor.

\begin{definition}
Let $V$ be a finite-dimensional vector space and
$T \xeof T_{r,s}(V)$ and $r,s > 0$.
Then the trace $\operatorname{Spur}(T) \xeof T_{r-1,s-1}(V)$
is defined according to
\[
\operatorname{Spur}(T)(f_1,\ldots,f_{i-1},f_{i+1},\ldots,f_r,
                       v_1,\ldots,v_{j-1},v_{j+1},\ldots,v_s) \, := 
\]
\[
f_i(v_j) \, \cdot \, T(f_1,\ldots,f_r,v_1,\ldots,v_s).
\]
\end{definition}

\subparagraph{}
Especially for a tensor $T \xeof T_{1,1}(V)$ we define
$\operatorname{Spur}(T) := f(v) \xeof \R$.

\section{Tensors in Euclidean Vector Spaces}
%-------------------------------------------
In the following
let $(V,g)$ a $n$-dimensional Euclidean vector space 
and $\iota \: V \xto V^\*$ the isomorphism
between the spaces $V$ and $V^\*$
induced by the inner product.
The pairs of vectors $(v, \iota v) \xeof \cross{V}{V^\*}$ 
are called \textit{associated to each other} {\wrt} the inner product.

\subparagraph{}
Chosing a base $\{e_i\}$ in $V$ and
the corresponding dual base $\{\epsilon^j\}$ in the dual space ${V^\*}$ ,
for the coordinates of $v$ and $f := \iota v$
\[
f(w) \eq f_i w^i \eq g(v^i e_i, w^k e_k) \eq g_{ik} \,
v^i w^k, \qquad w \xeof V
\]
holds,
where $g_{ik} := g(e_i,e_k)$ for $1 \leq i,k \leq n$.
By comparisons we obtain the relations
\[
f_i \eq g_{ik} \, v^k, \qquad (i=1,\ldots,n).
\]
Since an inner product $g$ is regular,
the matrix $(g_{ik})$ is invertible with inverse matrix
$g^{ik} := (g_{ik})^{-1}$,
and reversely
\[
v^k \eq g^{kj} f_j, \qquad (k=1,\ldots,n)
\]
holds.
These are the relationships of the coordinates for an associated pair
of vectors {\wrt} given by dual bases.

\subparagraph{}
By means of the mapping $\iota$
it is possible to assign to each ($r$,$s$)-tensor on $V$ a ($r+1$,$s-1$)-tensor
and a ($r-1$,$s+1$)-tensor,
supposed $s \geq 1$ and $r \geq 1$.
In the first case a covariant argument $v_i$ is replaced by
the contravariant argument $\iota v_i$, 
whereas in the second case the contravariant argument $f_i$ replaced
by covariant argument $\iota^{-1} f_i$. 
In both cases the stage and the function value of the tensor does not change.
Tensors,
which are characterized by such substitutions,
are called \textit{associated with each another}.

\subparagraph{}
Consider the following examples:

\begin{itemize}[leftmargin=15pt]
\item [1.]
to each simple covariant tensor $f \xeof V^\*$
there is a simple contravariant tensor $v \xeof V^{\*\*}$
such that
\[
f(w) \eq v(\iota w), \qquad w \xeof V.
\]
Reversely,
to each contravariant tensor $v \xeof V^{\*\*}$ 
there is a covariant tensor $f \xeof V^\*$
such that
\[
v(g) \eq f(\iota^{-1} g), \qquad g \xeof V^\*.
\]
The tensors $v$ and $f$ are associated to each another.

\item [2.]
The (0,2)-tensor $T_1$, the (1,1)-tensor $T_2$ and the (2,0)-tensor $T_3$
are associated with each another,
if
\[
T_1(v,w) \eq T_2(\iota v, w) \eq T_3(\iota v, \iota w) \qquad
\fa v,w \xeof V.
\]
\end{itemize}

\section{Anti-Symmetric Tensors}
%------------------------------------------------------------------------------%
In differential geometry the class of \textit{anti-symmetric}
tensors on a finite-dimensional $\R$-vector space is of special interest.

\begin{definition}
Let $V$ be a finite-dimensional $\R$-vector space and
$T \xeof T_{r,s}(v)$ with $r \geq 2$ or $s \geq 2$. 
The tensor $T$ is called \textit{anti-symmetric} rsp. \textit{alternating},
if in case of $s \geq 2$ for any vectors $v_1,v_2 \xeof V$
\[
T(\ldots,v_1,\ldots,v_2,\ldots) \eq - \, T(\ldots,v_2,\ldots,v_1,\ldots)
\]
and in case of Falle $r \geq 2$ for any covectors $f_1, f_2 \xeof V^\*$
\[
T(\ldots,f_1,\ldots,f_2,\ldots) \eq - \, T(\ldots,f_2,\ldots,f_1,\ldots).
\]
\end{definition}

\subparagraph{}
This implies
that an anti-symmetric tensor adpots the value $0$
iff
the set of covariant rsp. contravariant arguments is linearly dependent.

\subparagraph{}
The concept of anti-symmetry is thus defined for arbitrary $(r,s)$-tensors.
For applications,
however,
tensors are important,
if they are either purely covariant ($r=0$) or purely contravariant ($s=0$).

\subparagraph{}
The anti-symmetric and $s$-fold covariant tensors with $s \geq 2$ form
a \textit{linear subspace} of the vector space $T_{0,s}(V)$, 
denoted $\Lambda^s \, V$.
Especially for $s=0$ and $s=1$ let $\Lambda^0 V := \R$ and $\Lambda^1 V := V$.

\subparagraph{}
Similarly,
the $r$-fold contravariant anti-symmetric tensors with
$r \geq 2$ form a \textit{linear subspace} $\Lambda^r \, V^\*$ of
the vector space $T_{r,0}(V)$. 
Speziell for $r=0$ and $r=1$ let
$\Lambda^0 V^\* := \R$ and $\Lambda^1 V^\* := V^{\*\*}$.

\paragraph{}
Given a $d$-dimensional real vector space $V$,
a \textit{determinant function} is defined to be as a real-valued
and alternating $d$-fold multilinear form on $V$. 
But the following definition is equivalent:

\begin{definition}
A \textit{determinant function} on a $d$-dimensional vector space $V$
is an anti-symmetric tensor of level $(0,d)$.
\end{definition}

\subparagraph{}
A linear mapping $\phi \: V \xto V$ can be represented for a given base 
$\{e_i\}$ by means of a $(d,d)$-matrix $A$. 
The columns of $A$ are the coordinates of the vectors $A(e_i)$ {\wrt} this base.
The problem,
whether $\phi$ is an injective mapping,
is usually solved by computing the \textit{determinant} of $A$.
In fact, it owns the value $0$
iff
the column vectors of $A$ are linear dependent.

\subparagraph{}
A determinant function is called \textit{trivial}
if it adopts the value $0$ for arguments $e_1,\ldots,e_n$
forming a base of $V$.
Two non-trivial determinant functions $\delta_1, \delta_2$
differ in their values {\wrt} the arguments 
$v_1,\ldots,v_n$ by a \textit{constant} factor only, 
which depends on $\delta_1$ and $\delta_2$ only.
In this case there is a number $\lambda(\delta_1,\delta_2) \xeof \R$
such that
\[
\delta_1(v_1,v_2,\ldots,v_d) \= \lambda(d_1,d_2) \cdot
\delta_2(v_1,v_2,\ldots,v_d)
\]
for all $v_1,v_2,\ldots,v_d \xeof V$.

\begin{definition}
Two non-trivial determinant functions $\delta_1,\delta_2 \Lambda^d$
are called \textit{equivalent},
if $\lambda(\delta_1,\delta_2) > 0$.
\end{definition}

\subparagraph{}
Hence the set of all non-trivial determinant functions 
splits into two equivalence classes,
which both are denoted an \textit{orientation} of the underlying vector space.

\subparagraph{}
If an orientation of a vector space $V$ is given,
then the totality of bases of $V$ split into two classes: 
let $d$ be a determinant function of the given orientation,
then two bases $B$ and $B'$ are called \textit{equivalent to each other},
if
$d(b_1,\ldots,b_d) \cdot d(b'_1,\ldots,b'_d) > 0$.
This implies an equivalence relation on the set of all bases of $V$.
The bases of the first equivalence class
are called \textit{positively oriented},
whereas those  of the second one are called \textit{negatively oriented}.

\subparagraph{}
The tensor product of two anti-symmetric tensors is in general 
not an anti-symmetric tensor again.
Indeed,
the product of pure covariant tensors $f_1,f_2 \xeof V^\*$ is
\[
(f_1 \otimes f_2) (u,v) \eq f_1(u) \cdot f_2(v), \qquad u,v \xeof V,
\]
so it fulfills the symmetry condition
$(f_1 \otimes f_2) (u,v) \eq (f_1 \otimes f_2) (v,u)$
and is consequently not anti-symmetric.
Thus for purely covariant anti-symmetric tensors
one is seeking for a modified tensor product,
called the \textit{(exterior product)} $\wedge$, 
which for two anti-symmetric tensors $T_1$, $T_2$ provides 
an anti-symmetric tensor $T_1 \wedge T_2$ again.
For example,
if one defines for $f_1, f_2 \xeof V^\*$
\[
(f_1 \wedge f_2) (u,v) := f_1(u) \cdot f_2(v) - f_1(v) \cdot f_2(u),
\]
then $f_1 \wedge f_2$ is a two-fold covariant tensor fulfilling
\[
(f_1 \otimes f_2)(u,v) \eq - (f_1 \otimes f_2)(v,u).
\]

\subparagraph{}
Generalizing this approach,
the \textit{exterior product} of two purely covariant tensors
is introduced as follows:

\begin{definition}
Let $T_1$ be a (0,$s_1$)-tensor,
$T_2$ be a (0,$s_2$)-tensor and
$\pi$ be a permutation of the set $\{1,2,\ldots,s_1+s_2\}$.
Then the tensor $T_1 \wedge T_2\: \cross{V^{s_1}}{V^{s_2}} \xto \R$ with
\[
T_1 \wedge T_2(v_1,\ldots,v_{s1+s2}) := 
\frac{1}{s_1! s_2!} \sum_{\pi} sign(\pi) \,
(T_1 \otimes T_2) (v_{\pi(1)},\ldots,v_{\pi(s_1+s_2)})
\]
is called the \textit{exterior product} of $T_1$ and $T_2$.
\end{definition}

\subparagraph{}
The exterior product of purely contravariant tensors is defined
in a similar way.

\subparagraph{}
Taking the exterior product of two covariant tensors is consequently a mapping 
$\wedge \: \cross{\Lambda^{s_1} V}{\Lambda^{s_2} V} \xto \Lambda^{s_1+s_2} V$. 
Since in case of $s_1 + s_2 > d$ the $d$ arguments of the exterior product
are linearly dependent,
we have $T_1 \wedge T_2 \eq 0$
for each $s_1$-fold covariant tensor $T_1$
and $s_2$-fold covariant tensor $T_2$ with $s_1 + s_2 > d$.

\subparagraph{}
As an example,
let us compute the exterior product of two covectors 
$f,g \xeof \Lambda^1 \, V$ of a two-dimensional vector spaces $V$.
Let $\{\epsilon^1,\epsilon^2\}$ be a base of $V^\*$,
then with $f \eq f_1 \epsilon^1 + f_2 \epsilon^2$ and
$g \eq g_1 \epsilon^1 + g_2 \epsilon^2$ for the exterior product $f \wedge g$,
applied to two vectors $u,v \xeof V$ we obtain
\[
(f \wedge g)(u,v) \eq f(u)g(v) - f(v)g(u) \eq 
(f_1 \, g_2 - f_2 \, g_1) \, (v^1 w^2 - v^2 w^1).
\]
On the other hand,
$(\epsilon^1 \wedge \epsilon^2)(u,v) \eq v^1 w^2 - v^2 w^1$ holds
such that a comparison yields
\[
(f \wedge g)(u,v) \eq (f_1 \, g_2 - f_2 \, g_1) \cdot 
(\epsilon^1 \wedge \epsilon^2)(u,v)
\]
and hence
$f \wedge g \eq (f_1 g_2 - f_2 g_1) \, \cdot \, \epsilon^1 \wedge \epsilon^2$.

\section{The Vector Spaces $\Lambda^s \, V$}
%------------------------------------------------------------------------------%
Tensors,
which are built up by $s$-fold performing the exterior product
of dual base vectors $\epsilon^i$,
form a base of the vector spaces $\Lambda^s \, V$, 
where the products with ascending indication of the factors
are already sufficient to form this base. 
The following representation theorem for anti-symmetric tensors is nearby: 

\begin{theorem}
Let $T \xeof \Lambda^s V$ be an anti-symmetric covariant tensor and
$\{\epsilon^i\}$ be a base of $V^\*$.
Then there are real coefficients $t_{l_1 l_2 \ldots l_n}$
such that
\[
T \eq \sum_{l_1 < l_2 < \ldots < l_s} t_{l_1 l_2 \ldots l_s} \cdot
\epsilon^{l_1} \wedge \epsilon^{l_2} \ldots \wedge \epsilon^{l_s}.
\]
\end{theorem}

\subparagraph{}
This identity is called the \textit{canonic representation}
of an anti-symmetric tensor {\wrt} the base $\{\epsilon^1,\ldots,\epsilon^d\}$.

\subparagraph{}
A similar representation for $r$-fold contravariant tensors in the space
$\Lambda^r \, V^\*$ is valid:
\[
T \eq \sum_{l_1 < l_2 < \ldots < l_r} t^{l_1 l_2 \ldots l_r} \cdot
e_{l_1} \wedge e_{l_2} \ldots \wedge e_{l_r},
\]
where the vectors $e_1,\ldots,e_d$ form a base of $V$.
Since it is possible to elect from a set of $n$ elements exactly $\binom{n}{s}$ 
ordered $s$-tuples,
the spaces $\Lambda^s\,V$ and $\Lambda^r\,V^\*$ own the dimensions
\[
dim \, \Lambda^s\,V    \, \eq \, \binom{n}{s}, \qquad
dim \, \Lambda^r\,V^\* \, \eq \, \binom{n}{r}.
\]

\subparagraph{}
Let the vector space $\Lambda \, V$
be the direct sum of the spaces $\Lambda^sV (s=1,\ldots,n)$,
\[
\Lambda \, V := 
\Lambda^0 \, V \oplus \Lambda^1 \, 
             V \oplus \Lambda^2 \, 
             V \, \oplus \, \ldots \, \oplus \Lambda^n \, V.
\]
Then $\Lambda\,V$ together with the exterior product forms
a graduated anti-commutative $\R$-Algebra,
called the \textit{exterior algebra} rsp. \textit{Gra\"ssmann algebra} on $V$.
For the dimension of the exterior algebra $\Lambda \, V$ we have
\[
dim \, \Lambda\,V \, \eq \, \sum_{k=0}^d \binom{n}{k} \, \eq \,
\sum_{k=0}^d dim \Lambda^k V \, \eq \, 2^d.
\]

\paragraph{}
The vector spaces $\Lambda^s\,V$ and $\Lambda^s\,V^{\*}$ are related
to each other by their dual spaces:

\begin{theorem}
Let $V$ be a $d$-dimensional vector space.
Then $(\Lambda^s \, V)^\* \eq \Lambda^s \, V^\*$ for $s \eq 0,\ldots,d$ holds,
that is,
the spaces $\Lambda^s \, V$ and $\Lambda^s \, V^\*$ are dual to each other.
\end{theorem}

\subparagraph{}
The isomorphism $\iota$ between these spaces is defined in a natural way.
If $\phi \xeof (\Lambda^s\,V)^\*$, then
\[
(\iota \phi)(f_1,\ldots,f_s) := \phi(f_1 \wedge \ldots \wedge f_s),
\qquad f_i \xeof V^{\*}.
\]
The mapping $\iota$ is injective and linear.
Since both spaces own one and the same dimension,
$\iota$ actually proves to be surjective.

\paragraph{}
In the following
let $(V,b)$ be a $d$-dimensional \textit{Euclidean} vector space. 
The presence of an inner product enables the construction of inner products
in the spaces $\Lambda^s\, V$.

\subparagraph{}
The natural isomorphisms between $V$ and $V^\*$ on the one hand side and
$\Lambda^s\,V^\*$ and $(\Lambda^s\,V)^\*$ on the other side are crucial.
Its combination yields a natural isomorphism $\iota$
between the spaces $(\Lambda^s\,V)^\*$ and $\Lambda^s\,V$.
The inner product in the space $\Lambda^s\,V$ can be defined for
$f_1,f_2 \xeof \Lambda^s\,V$ according to
\[
\bracket{f_1,f_2} := (\iota^{-1}(f_1))(f_2) \eq (\iota^{-1}(f_2))(f_1).
\]
To the end,
the bilinear form $\bracket{f_1,f_2}$ is symmetric and regular and constitutes
an inner product on the space $\Lambda^s\,V$.

\subparagraph{}
It is our final aim to establish a \textit{duality} between the spaces 
$\Lambda^s\,V$ and $\Lambda^{n-s}\,V$:
performing for two tensors $f \xeof \Lambda^s\,V$ and
$g \xeof \Lambda^{n-s}\,V$ the exterior product $f \wedge g$,
then $f \wedge g \xeof \Lambda^n\,V$.
The vector space $\Lambda^n\,V$ of the determinant functions on $V$
is one-dimensional.
Selecting the \textit{canonical volume form} $\epsilon$ as a base
of this space and
assigning to each positive oriented orthonormal base $B \subset V$
the value $+1$,
then
\[
f \wedge g \= \Delta(f,g) \cdot \epsilon, \quad \Delta(f,g) \xeof \R.
\]

\begin{theorem}
Let $(V,b)$ be a $d$-dimensional, oriented and Euclidean vector space.
Then to each linear form $h \xeof \Lambda^s\,V$
there is exactly one linear form $h^\* \xeof \Lambda^{n-s}\,V$
such that
for arbitrary linear forms $f$
\[
f \wedge h^\* \eq g(f,h) \cdot \epsilon.
\]
\end{theorem}

\paragraph{}
This theorem gives rise to the following definition:

\begin{definition}
Let $(V,b)$ be a $d$-dimensional, oriented and Euclidean vector space.
The mapping 
$\* \: \Lambda^s\,V \xto \Lambda^{n-s}\,V, \, \*(h) \eq h^\*$
is called the \textit{Hodge operator} or the \textit{star operator}
on the space of linear forms of degree $s$.
\end{definition}

\subparagraph{}
The Hodge operator is involutional (up to the sign):

\begin{theorem}
For a linear form $f \xeof \Lambda^s\,V$ on a $n$-dimensional
vector space $V$ we have
$\*\* f \eq (-1)^{s(n-s)} \cdot f$.
\end{theorem}

\section{Induced Mappings}
%------------------------------------------------------------------------------%
Let $V,W$ be $n$-dimensional $\R$-vector spaces and
$\phi \: V \xto W$ be a linear mapping.
Then for $i=1,\ldots,n$ the mappings
\[
\phi^\*_i \: \Lambda^i \, W \xto \Lambda^i \, V
\]
are defined between the spaces $\Lambda^i\,W$ and $\Lambda^i\,V$,
where the values of $\phi^\*_i$ are given element-wise
according to
\[
\phi^\*_i(f) := f^\* \fa f \in\Lambda^i \, W,
\]
\[
\phi^\*(v_1,\ldots,v_n) := f(\phi(v_1),\ldots,\phi(v_n)). 
\]
The operator $\phi$ need not be injective or surjective.

\begin{definition}
The mappings $\phi^\*_i$ are called the \textit{induced} ones
by $\phi \xeof Hom_{\R}(V,V)$
between the vector spaces $\Lambda^i\,W$ and $\Lambda^i\,V$.
\end{definition}

\section{Lie-Rings}
%------------------------------------------------------------------------------%
In applications a special class of rings is of interest,
whose 
multiplication is in general not associative.

\begin{definition}
Let $(L,+)$ be an Abelian group and $[\,] \: L \xtimes L \to L$
be a further but not necessarily associative operation.
The triple $(L,+,[\,])$ is called a \textit{Lie ring} if

\begin{itemize}[leftmargin=35pt]
\item[(L-1)]
$[x,x] := [\,](x,x) \eq 0$ for all $x \xeof L$
\item[(L-2)]
$[x,[y,z]] + [y,[z,x]] + [z,[x,y]] \eq 0$ for all $x,y,z \xeof L$
(\textit{Jacobi identity}).
\item[(L-3)]
$[x,a+b] \eq [x,a] + [x,b]$ and $[a+b,y] \eq [a,y] + [b,y]$ (distribution laws).
\end{itemize}
\end{definition}

The operation $[\,]$ is often denoted the \textit{commutator map}.
Let us mention the following examples of Lie rings:

\begin{itemize}[leftmargin=15pt]
\item[1.]
the Abelian group $(\R^3,+)$ with the vector product $\times$ 
as commutator map,
that is,
$[a,b] := a \times b$.

\item[2.]
the Abelian group of real ($n$,$n$)-matrices with the matrix operation
$[A,B] := A \circ B - B \circ A$.
\end{itemize}

\paragraph{}
A textit{Lie ring homomorphism} between the Lie rings $L$ and $M$ is
a group homomorphism $\phi \: L \xto M$, 
which respects the commutator map in the sense of
\[
\phi([x,y]) \= [ \phi(x) \, ,  \, \phi(y) ]
\]
for all $x,y \in L$.
Moreover,
if $\phi$ is bijective,
then $\phi$ is said to be a \textit{Lie Ring isomorphism}.
Further concepts in ring theory,
like subrings, ideals and so on can be defined for Lie rings as well.

\section{Lie Algebras}
%------------------------------------------------------------------------------%
Lie-algebras are non-associative algebraic structures, 
which (considered as a ring) form Lie rings.
Linear operators, differential operators ans vector fields on
manifolds are mathematical objects
that can be endowed with the structure of a Lie algebra.

\begin{definition}
Let $R$ be a unital ring, $L$ a $R$-module and
$[\,] \: L \xtimes L \xto L$ a not necessarily associative operation.
The triple $(L,+,[\,])$ is called a \textit{$R$-Lie algebra}
if $(L,[\,])$ is a Lie ring.
\end{definition}

\paragraph{}
Consider the following examples of Lie algebras:

\begin{itemize}[leftmargin=15pt]
\item[1.]
The $\R$-module $\R^3$ with the vector product $[\,]$ is a $R$-Lie algebra.
\item[2.]
The ring of endomorphisms $E(V,V)$ of a $\F$-vector space with the
commutator map $[\phi,\psi] := \phi \circ \psi - \psi \circ \phi$.
\item[3.]
The $\R$-module of ($n$,$n$)-matrices with matrix operation
$[A,B] := A \circ B - B \circ A$
is a $R$-Lie algebra.
\item[4.]
The set of derivations of a $\F$-algebra with $[D_1,D_2] := D_1D_2-D_2D_1$.
\end{itemize}

\paragraph{}
If $L_1$ and $L_2$ are Lie algebras,
then its cartesian product
$\cross{L_1}{L_2}$ is a Lie algebra again.
The mapping $[\,]$ is defined pointwise,
that is,
\[
[(x_1,y_1),(x_2,y_2)] \eqdef ([x_1,x_2], [y_1,y_2])
\]
for $x_1,x_2 \xeof L_1$ and $y_1,y_2 \xeof L_2$.

\subparagraph{}
An element $x \xeof L$ is called \textit{central},
if it commutes with all other elements of $L$,
that is,
$[a,y] \eq 0$ holds for all $y \xeof L$.
The \textit{center} $Z(L)$ of a Lie algebra $L$ is the set of central elements in $L$.
A Lie-algebra $L$ is said to be \textit{commutative},
if all elements are central,
that is,
$Z(L) \eq L$.

\subparagraph{}
A subset $U$ of a Lie algebra $L$ is called a \textit{Lie subalgebra},
if $U$ is an submodule of $L$ and $[x,y] \xeof U$ for all $x,y \xeof U$.
Moreover,
$U$ is denoted an \textit{ideal} in $L$,
if $[x,y] \xeof L$ for all $x \xeof U$ and all $y \xeof U$.

\paragraph{}
A \textit{Lie-algebra homomorphism} between two Lie-algebras is
a module homomorphism and a Lie-ring homomorphism at the same time.

\subparagraph{}
If $L$ is a Lie-algebra and $I$ is an ideal in $L$,
the $L/I$ is in a natural way a $R$-module as well as a Lie ring.
Lie-algebras can be factorized by their ideals,
and the quotient $L/I$ is a Lie-algebra again.
The projection $\pi \: L \xto L/I$ constitutes a Lie-epimorphism.

\subparagraph{}
Let $\phi \:L_1 \xto L_2$ be a Lie-algebra homomorphism between
Lie-algebras $L_1$ and $L_2$.
Then $\operatorname{Kern} \phi$ is an ideal in $L_1$,
and $L_1 / \operatorname{Kern} \phi$ is isomorphic to $\phi(L_1)$.

\subparagraph{}
Each element $x \xeof L$ of a Lie-algebra $L$ generates
a Lie-algebra endomorphism $\phi_x \: L \xto L$
according to $\phi_x(y) := [x,y]$ for all $y \xeof L$.

\paragraph{}
For finite-dimensional Lie-algebras over fields of characteristics
$\chi(\F) \neq 2$ the following structure theorem is valid (\textsc{Ado}'s theorem):

\begin{theorem}
Let $\F$ be a field with $\chi(\F) \neq 2$.
Then each finite-dimensional Lie-algebr over $\F$ is isomorphic
to a Lie-subalgebra of $\operatorname{GL}(n,\F)$ for some $n \xeof \N$.
\end{theorem}

\subparagraph{}
The Lie-algebras over a ring $R$ form a category denoted by $Lalg_R$.

\paragraph{}
Let us switch to representations of Lie algebras:

\begin{definition}
Let $L$ be a Lie algebra.
A \textit{representation} of $L$ on the $\F$-vector space $V$
is a
group homomorphism $\phi \: L \xto Hom_{\F}(V,V)$,
which is compatible with the commutator map in $L$ in the sense of
\[
\phi([x,y]) \= \phi(x) \circ \phi(y) - \phi(y) \circ \phi(x) \quad (x,y \in L).
\]
\end{definition}

\subparagraph{}
A representation $\phi$ is called \textit{faithful}, if $\phi$ is injective.
If the vector space $V$ owns dimension $n \xeof \N$,
then $n$ is denoted the \textit{dimension} of the representation.

\subparagraph{}
Let $\phi \: L \xto Hom_{\F}(V,V)$ be a representation of $L$ on $V$.
Then each linear operator $\phi(x) \in Hom_{\F}(V,V)$
gives rise to a representation
$\phi^T:L \to Hom_{\F}(V^{\*},V^{\*})$ von $L$ on the dual space $V^{\*}$
by means of the dual operator  $\phi^{\*}(x) \in Hom_{\F}(V^{\*},V^{\*})$,
where $\phi^T$ is defined by $x \mapsto -\phi^{\*}(x)$.

\begin{definition}
The mapping $\phi^T$ is called the \textit{dual representation} of $L$ on $V$.
\end{definition}

\paragraph{}
Furthermore,
for each fixed element $x \in L$
a linear operator
$\operatorname{ad}_x \in Hom_{\F}(L,L)$ is defined by the operation $[\,]$.
The assignment $x \mapsto \operatorname{ad}_x$ is a group homomorphism,
which respects the operation $[\,]$ respektiert,
and consequently a representation of $L$ on itself.

\begin{definition}
$\operatorname{ad} \: L \xto Hom_{\F}(L,L)$ with $x \mapsto \operatorname{ad}_x$
is said to be the \textit{adjoint} representation of $L$.
\end{definition}

\subparagraph{}
If $L$ is finite-dimensional (considered as a $\F$-vector space),
then for each linear operator $\phi \in Hom_{\F}(L,L)$
the \textit{trace} is well-defined as an element in $\F$.
Let for $x,y \in L$ the mapping $B \: \cross{L}{L} \xto \F$ be defined according to
\[
B(x,y) \eqdef Trace (\operatorname{ad}_x \circ \operatorname{ad}_y) \in \F.
\]
Then $B$ is a symmetric bilinear form on the Lie-algebra $L$,
called the \textit{Killing form} of $L$.

\newpage
\thispagestyle{empty}
\fancyhead[C] {\small{\textsc{III. Tensor Fields and Differential Forms}}}
%------------------------------------------------------------------------------%
\part{Tensor Fields and Differential Forms}                                    %
%------------------------------------------------------------------------------%
Continuous mappings between topological spaces are mappings,
which are compatible with the topological structures of the spaces.
In the same way one is searching for differentiable manifolds and mappings,
which are compatible with their differentiable structures.
To the end we obtain the concept of differentiability
for mappings between manifolds.

\section{Vector Fields}
%------------------------------------------------------------------------------%
Vector and tensor fields are the base objects of \textit{tensor analysis}
on manifolds.

\subparagraph{}
A \textit{scalar field} on a differentiable manifold $\difM$ is
a differentiable mapping $f\:\difM \xto \R$.
The set $C^{\infty}(\difM,\R)$ of all scalar fields on $\difM$ forms
with pointwise defined addition rsp. multiplication and
with scalar multiplication a commutative unital $\R$-algebra.

\subparagraph{}
If $f\:\difM \xto \difN$ is a differentiable mapping
between manifolds and $\phi$ is a scalar field on $\difN$,
then the mapping $\phi \circ f$ is a scalar field on $\difM$,
called the \textit{pullback} of $\phi$ under $f$.
The assignment
$\phi \xeof C^{\infty}(\difN,\R) \mapsto
 \phi \circ f \xeof C^{\infty}(\difM,\R)$
is compatible with the inner operations of these algebras
and thus forms an algebra homomorphism.
This construction provides a contravariant functor of the
category of differentiable manifolds to the category of algebras.

\paragraph{}
A \textit{vector field} assigns to each point $p$ of a differentiable
manifold a tangent vector of the jeweiligen tangential space.
Considering the tangential space itself as a differentiable manifold,
the concept of \textit{differentiable vector field} is meaningful.

\begin{definition}
Let $\difM$ be a $d$-dimensional differentiable manifold
and $T \difM$ be the related tangential space.
Then a (differentiable) \textit{vector field} on $\difM$
is a (differentiable) mapping $\vecX:\difM \xto T \difM$
with $\vecX(p) \xeof T_p \difM$.
\end{definition}

A vector field is often called
the \textit{contravariant tensor field of first stage}.

\subparagraph{}
If a vector field $\vecX$ on $\difM$ is differentiable,
then there is a representation of $\vecX$
for each chart $(\openO,\kappa) \xeof \cal{A}(\difM)$ according to
\[
\vecX(p) \= \sum_{i=1}^n x^i \partial_i(p), \quad p \xeof \openO,
\]
where the functions $x^i \: \kappa(\openO) \xto \R$ are differentiable.

\subparagraph{}
For a scalar field $f$ and a (differentiable) vector field $\vecX$
the product $f \cdot \vecX$ is a (differentiable) vector field again.
Thus the set $\cal{V}(\difM,\R)$ of differentiable vector fields on $\difM$
forms with pointwise defined addition and multiplication
with scalar fields a module over the ring $C^{\infty}(\difM,\R)$.

\subparagraph{}
A differentiable \textit{Covector field} on $\difM$ is defined similarly:

\begin{definition}
A \textit{Covector field} rsp. a \textit{linear form} on a
differentiable manifold $\difM$ is a differentiable mapping
$\vecX^\*:\difM \xto T^\*(\difM)$,
which assigns to each point $p \xeof \difM$ a covector
in the dual space $T_p^{\*}(\difM)$ of the tangential space.
\end{definition}
% xxx

\subparagraph{}
A linear form wird oft auch as \textit{covariant tensor field erster Stufe}
bezeichnet.
The set $\cal{V}^\*(\difM;\R)$ of linear forms on $\difM$ forms with
pointwise defined addition and multiplication with scalar fields
also a module over the ring $C^{\infty}(\difM,\R)$.

\section{Lie Products}
%------------------------------------------------------------------------------%
Let $\vecX,\vecY$ be two vector fields on a differentiable
manifold $\difM$.
The \textit{Lie-product} $[\vecX,\vecY]$ is a mapping,
which in each point $p \xeof \difM$ den beiden tangent vectors
$\vecX(p),\vecY(p)$ wieder a tangent vector aus $T_p \difM$ zuordnet,
and zwar in folgender Weise:

\subparagraph{}
For a scalar field $f$ on $\difM$
the mapping
$f_X: p \xeof \difM \mapsto \vecX(f) \xeof \R$
is a scalar field on $\difM$ again.
The erneute mapping $p \xeof \difM \mapsto Y(f_X) \xeof \R$ is demnach the
Verkettung $(\vecY \circ \vecX)$
and weist the scalar field $f$ a further scalar field zu.
$(\vecY \circ \vecX)$ is jedoch im allgemeinen kein vector field,
da the mapping the geforderte product rule for tangent vectors
nicht erf\"ullt (and hence no derivative any more).
The difference after changing the order
$(\vecY \circ \vecX)(f) - (\vecX \circ \vecY)(f)$
fulfills all properties of a tangent vector.

\begin{definition}
The \textit{Lie product} of two vector fields $\vecX,\vecY$
auf of a differentiable manifold $\difM$ is the vector field
\[
[\vecX,\vecY](f) \, := \, (\vecY \circ \vecX)(f) - (\vecX \circ \vecY)(f),
\quad f \xeof C^{\infty}(\difM,\R).
\]
\end{definition}

\subparagraph{}
The Lie product owns the following properties:
\begin{itemize}[leftmargin=15pt]
\item[1.]
$[\vecX,\vecY]$ is linear in both arguments and anti-symmetric.
\item[2.]
For each scalar field $f$ holds
$[f\vecX,\vecY] = f [\vecX,\vecY] - \vecY(f) \vecX$.
\item[3.]
Je drei vector fields $\vecX,\vecY,\vecZ$ erf\"ullen the \textit{Jacobi-Identit\"at}
\[
[\vecX,[\vecY,\vecZ]] + [\vecY,[\vecZ,\vecX]] + [\vecZ,[\vecX,\vecY]] = 0.
\]
\end{itemize}

\subparagraph{}
The set $\cal{V}(\difM,\R)$ of vector fields
forms with pointwise defined addition
and the Lie product as multiplication a Lie ring.
Nimmt man the pointwise defined multiplication of vector fields
with scalar fields hinzu,
then $\cal{V}(\difM,\R)$ forms a Lie algebra
over the ring $C^{\infty}(\difM,\R)$.

\subparagraph{}
Let $\{\partial_1,\ldots,\partial_n\}$ be a natural base of $T_p \difM$
and $\vecX^i, \vecY^i$ be the coordinates of the vector fields $\vecX,\vecY$.
Then gelten for the coordinates of the Lie-product $[\vecX,\vecY]$
die Gleichungen
\[
[\vecX,\vecY]^i
\, = \,
X^j \partial_j \vecY^i - \vecY^j \partial_j X^i, \qquad i=1,\ldots,n.
\]

Especially for natural base vectors we have $[\partial_i,\partial_j] \eq 0$.

\subparagraph{}
Lie-products play an important role in the coordinate-free description
of torsion and curvature forms of affine connections.
wichtige Rolle.

\section{Tensor Fields}
%------------------------------------------------------------------------------%
In generalization of the concepts \textit{scalar field}
and \textit{vector field} tensor fields on a manifold are defined as follows:

\begin{definition}
Let $\difM$ be a differentiable manifold.
A $s$-fach covariant and $r$-fach contravariant \textit{tensor field} $T$
on $\difM$ is a differentiable mapping,
which assignes to each point $p \xeof \difM$ an $(r,s)$-tensor
over the tangential space $T_p \difM$.
\end{definition}

\subparagraph{}
Differentiability of the tensor field again means
that for each chart $(\Off,\kappa) \xeof \cal{A}(\difM)$
the coordinate functions
\[
T_{i_1,i_2,\ldots,i_s}^{j_1,j_2,\ldots,j_r} \:
\kappa(\Off) \subseteq \R^d \xto \R
\]
are differentiable.

\subparagraph{}
The set $\cal{T}_r^s(\difM,\R)$ of $s$-fach covariant and $r$-fach
contravariant tensor fields on $\difM$ forms a module
over the ring $C^{\infty}(\difM,\R)$, too.

The Lie-derivative is an extension of the Lie-bracket for tensor fields.

\subparagraph{}
Let $\vecX$ be a vector field on a differentiable manifold $\difM$.
For each further vector field $\vecY \xeof \cal{V}(\difM,\R)$ sei the mapping
$\operatorname{L}_X: \cal{V}(\difM,\R) \xto \cal{V}(\difM,\R)$ erkl\"art gem\"a\"ss
$\operatorname{L}_X(\vecY) := [\vecX,\vecY]$.

\begin{definition}
$\operatorname{L}_X$ is called the \textit{Lie-derivative} von
$\vecY \xeof \cal{V}(\difM,\R)$ {wrt} of the vector field $\vecX$.
\end{definition}

\subparagraph{}
For a scalar field $f \xeof C^{\infty}(\difM,\R)$ wird the Lie-derivative
gem\"a\"ss $\operatorname{L}_X f := \vecX(f) = d \, f(\vecX)$ definiert,
that is,
$\operatorname{L}_X f$ stellt wiederum a scalar field dar.
Damit is the Lie-derivative $\operatorname{L}_X$ jeweils as operator
auf den modules $\cal{V}(\difM,\R)$ and $C^{\infty}(\difM,\R)$ erkl\"art.

\subparagraph{}
The Lie-derivative l\"a\"sst sich nun to an operator on den modules der
$(r,s)$-tensor fields erweitern.
The Bild of a $(r,s)$-tensor field soll hierbei wiederum
ein $(r,s)$-tensor field sein.
Ferner soll $\operatorname{L}_X$ linear on thesen modules operieren.

\begin{theorem}
Let $\vecX$ be a vector field on a differentiable manifold $\difM$.
Then there is genau a mapping
\[
\operatorname{L}_X:T_s^r(\difM) \xto T_s^r(\difM),
\]
welche for $r = s = 0$ the Lie-derivative of a scalar field and f\"ur
$r=1, s=0$ the Lie-derivative of a vector field ergibt.
\end{theorem}

\subparagraph{}
The tensor field $\operatorname{L}_X T$ is called the \textit{Lie-derivative}
of the tnesor field $T$ {\wrt} the vector field $\vecX$.

\section{Differential Forms}
%------------------------------------------------------------------------------%
\textit{Differential forms} play an important role
especially in integration theory on manifolds.

\begin{definition}
Let $\difM$ be a $d$-dimensional differentiable manifold.
A \textit{differential form of degree $s$} or briefly a \textit{$s$-form}
is a $s$-fach covariant, antisymmetric tensor field on $\difM$.
\end{definition}

\subparagraph{}
The following examples of differential forms are known:

\begin{itemize}[leftmargin=15pt]
\item[1.]
The scalar fields on $\difM$ are the differential forms of degree $0$.
\item[2.]
The covariant differential forms of degree $1$ on $\difM$ are
exactly the covector fields.
They are called \textit{linear forms} or \textit{Pfaff forms}
\end{itemize}

\subparagraph{}
Let for $s \xeof \N$ the set $\Omega^s \difM$
be the submodule of all differential forms of degree $s$ on $\difM$.
Especially for $s=0$ let $\Omega^0 \difM := C^{\infty}(\difM,\R)$,
that is,
the module of scalar fields.
Since the tangential spaces of $d$-dimensional differentiable manifolds
are $d$-dimensional,
too,
we obviously have $\Omega^s \difM \eq 0$ for $s > n$.

\subparagraph{}
If a point $p \xeof \difM$ is in the domain of a chart $(O,\kappa)$,
then the forms
\[
dx^i \wedge dx^j \wedge \ldots \wedge dx^s \qquad (i < j < \ldots < s)
\]
constitute a base of the vector space of anti-symmetric tensors of degree $s$
over $T_p \difM$.
Hence a differential form $\omega$ in such a domain $O$
has the representation
\[
\omega = \sum_{i < j < \ldots < s} w_{ij\ldots s}
dx^i \wedge dx^j \wedge \ldots \wedge dx^s.
\]
The coefficients $w_{ij\ldots s}$ are differentiable functions
$w_{ij\ldots s}:\kappa(O) \xto \R$.

\subparagraph{}
Let $\difM,\difN$ be differentiable manifolds and
$f\:\difM \xto \difN$ a differentiable mapping.
Then $f$ induces in a natural way a mapping
\[
f^\* : \Omega^s \difN \xto \Omega^s \difM, \qquad s=1,\ldots,n,
\]
where $f^\*$ definiert is gem\"a\"ss
\[
f^\* \omega(v_1,\ldots,v_n) := \omega(\operatorname{d}\,f(v_1),\ldots,
                                      \operatorname{d}\,f(v_n))
\]
for $v_1,\ldots,v_n \xeof T_p(\difM)$ and each point $p \xeof \difM$.
The assignment $f^\* : \Omega^s \difN \xto \Omega^s \difM$,
$\omega \xeof \Omega^s \difN \mapsto f^\* \omega \xeof \Omega^s \difM$
is called the \textit{pullback} of $f$.

\subparagraph{}
On the modules of differential forms of a manifold
one may define a mapping $\operatorname{d}$
that assignes to each $s$-form a $(s+1)$-form
such that
the following properties hold:

\begin{itemize}[leftmargin=15pt]
\item[1.]
$\operatorname{d}^2 \, \omega
:=
\operatorname{d} \, (\operatorname{d}\omega ) = 0$
(Komplexbedingung)
\item[2.]
$\operatorname{d}(\lambda \, \omega_1 + \mu \, \omega_2) =
  \lambda \, \operatorname{d} \omega_1 + \mu \, \operatorname{d} \omega_2$
(Linearit\"at)
\item[3.]
$\operatorname{d} \, f(v) = v(f)$ for a 0-form $f$
\item[4.]
$\operatorname{d} (\omega_1 \wedge \omega_2) =
 \operatorname{d}  \omega_1 \wedge \omega_2  +
 (-1)^s
                   \omega_1 \wedge (\operatorname{d} \omega_2)$,
where $s$ is the degree of $\omega_1$ (product rule).
\end{itemize}

\subparagraph{}
The existence of such a mapping
$\operatorname{d}:\Omega^s \difM \xto \Omega^{s+1} \difM$
is ensured by the operator
\[
\operatorname{d} \, \omega := \sum_{k, \, i < j < \ldots < s}
\frac{\partial \, \omega_{ij\ldots s}}{\partial x_k} \quad
dx^i \wedge dx^j \wedge \ldots \wedge dx^s,
\quad \omega \xeof \Omega^s \difM,
\]
since if fulfills all the required properties for $\operatorname{d}$.
Moreover,
this is the only mapping fullfilling the just mentioned properties,
leading to the following definition:

\begin{definition}
Let $\omega \xeof \Omega^s \difM$ be a differential form of degree $s$
on $\difM$.
The durch obige Forderungen eindeutig bestimmte Differentialform
$\operatorname{d} \omega \xeof \Omega^{s+1} \difM$
is called the \textit{exterior derivative} rsp. das
\textit{exterior Differential} of $\omega$.
\end{definition}

\subparagraph{}
A representation of the exterior derivative of $\omega$
without using coordinates can be obtained by means of the Lie product:
\[
\operatorname{d} \omega(u,v) \= u(\omega(v)) - v(\omega(u)) - \omega([u,v]),
\qquad u,v \xeof T_p(\difM).
\]

\subparagraph{}
Ist um a point $p \xeof \difM$ a cobase
gem\"a\"ss $\{\omega^i\}$ gegeben,
then is notwendigerweise
$\operatorname{d} \, \omega^i = 0$ for $i=1,\ldots,n$.

\paragraph{}
Consider the following examples:

Let $\difM = \R^3$ and $\Omega \xsubset \R^3$.
\begin{itemize}[leftmargin=15pt]
\item[1.]
A $0$-form is a scalar-valued function $f\:\Omega \xto \R$.
The exterior differential $\operatorname{d}f$ of $f$ is defined by
\[
\operatorname{d}f((v_1,v_2,v_3)) = v(f) =
\frac{\partial f}{\partial x} v_1 +
\frac{\partial f}{\partial y} v_2 +
\frac{\partial f}{\partial z} v_3,
\]
hence
\[
\operatorname{d}f = \frac{\partial f}{\partial x} dx +
                    \frac{\partial f}{\partial y} dy +
                    \frac{\partial f}{\partial z} dz.
\]
$\operatorname{d} f$ is called the \textit{gradient} of $f$.

\item[2.]
The \textit{divergence} of a continuously differentiable vector field
$A \xeof C^1(\Omega,\R^3)$ is the scalar-valued function
\[
\operatorname{div} \, A := \frac{\partial A_1}{\partial x} +
                           \frac{\partial A_2}{\partial y} +
                           \frac{\partial A_3}{\partial z}.
\]
\item[3.]
The $1$-forms $\omega$ are covectors of the form
\[
\omega = a \, dx + b \, dy + c \, dz,
\]
where $a,b,c:\R^3 \xto \R$ are differentiable functions.

\item[4.]
For a continuously differentiable vector field $A \xeof C^1(\Omega,\R^3)$
the \textit{rotation} of $A$ is defined as the vector field
\[
\operatorname{rot} A :=
    (\frac{\partial A_3}{\partial y} - \frac{\partial A_2}{\partial z}, \,
     \frac{\partial A_1}{\partial z} - \frac{\partial A_3}{\partial x}, \,
     \frac{\partial A_2}{\partial x} - \frac{\partial A_1}{\partial y}).
\]
The exterior derivative of a covector is
\[
\operatorname{d} \omega =
(\frac{\partial c}{\partial y} - \frac{\partial b}{\partial z}) dy \wedge dz +
(\frac{\partial a}{\partial z} - \frac{\partial c}{\partial x}) dz \wedge dx +
(\frac{\partial b}{\partial x} - \frac{\partial a}{\partial y}) dx \wedge dy.
\]
\item[5.]
Hence the differential operator $\operatorname{grad} \, f$ corresponds to 
the exterior derivative $\operatorname{d}$,
applied to the scalar field $f$.
The operator $\operatorname{rot} \, A$ is the exterior derivative
applied to a linear form $A$,
and $\operatorname{div} \, A$ is the exterior derivative
einer aus $A$ konstruierten 2-form $A'$.
The complex-property $\operatorname{d} \circ \operatorname{d} = 0$
der exterior derivative f\"uhrt daher to den Gleichungen
$\operatorname{rot} \, \operatorname{grad} f = 0$
(that is, a gradient field is always wirbelfrei) and
$\operatorname{div} \operatorname{rot}  A = 0$
(that is, a Wirbel field is always divergence free).
\end{itemize}

\paragraph{}
The exterior derivative commutes with the Lie-derivative.

\begin{theorem}
Let $\difM$ be a differentiable manifold,
$\vecX$ a vector field and $\omega \xeof \Omega^s \difM$.
Then holds
$\operatorname{L}_X (\operatorname{d}   \omega) =
 \operatorname{d}   (\operatorname{L}_X \omega)$.
\end{theorem}

\subparagraph{}
Ferner is the exterior derivative with den Pullbacks
von differentiable mappings vertauschbar.

\begin{theorem}
Let $\difM,\difN$ be differentiable manifolds and
$\omega \xeof \Omega^s \difN$.
Then holds
$f^\* (\operatorname{d} \omega) = \operatorname{d} (f^\* \omega)$.
\end{theorem}

\section{Integration on Manifolds}
%------------------------------------------------------------------------------%
differential forms are the \textit{nat\"urlichen} mathematischen Objekte,
for the a Integrationstheorie on orientierten differentiable
manifolds entwickelt werden kann.

\subparagraph{}
Gewisse charts of a Atlanten are in folgender Weise miteinander
compatible.

\begin{definition}
Let $\difM$ be a differentiable manifold.
\begin{itemize}[leftmargin=15pt]
\item[1.]
For two charts $(\Off_1,\kappa_1),(\Off_2,\kappa_2) \xeof \cal{A}(\difM)$
is called der Kartenwechsel $\tau = \kappa_2 \circ \kappa_1^{-1}$
\textit{orientierungserhaltend},
if for all points $p \xeof \kappa_1(\Off_1) \cap \kappa_2(\Off_2)$ the
Jacobi-Matrizen $J_p \tau$ den Bedingungen
$\operatorname{det} J_p \tau > 0$ gen\"ugen.
\item[2.]
An atlas $\cal{A}$ on $\difM$ is called a \textit{orientierender atlas},
if er nur charts enth\"alt,
deren Kartenwechsel orientierungserhaltend sind.
\end{itemize}
\end{definition}

\subparagraph{}
manifolds, the a orientierenden Atlanten besitzen,
werden wie folgt ausgezeichnet.

\begin{definition}
A differentiable manifold is called \textit{orientiert}
rsp. \textit{orientable},
if ihre differentiable Struktur a orientierenden Atlanten enth\"alt.
\end{definition}

Alternativ hierzu l\"a\"sst sich the concept of oriented manifold
my means of differential forms definieren.

\subparagraph{}
A \textit{volume form} on a $d$-dimensional manifold $\difM$
is a differential form of degree $d$ $\omega \xeof \Omega^n \difM$,
die in keinem point $p \xeof \difM$ verschwindet.
Falls on $\difM$ a volume form existiert, so is called $\difM$ orientable.

\begin{definition}
Let $\difM$ be a $d$-dimensional differentiable manifold.
\begin{itemize}[leftmargin=15pt]
\item[1.]
A subset $\A \xsubset \difM$ is called \textit{measurable},
if for each covering of $\A$ by charts $(\Off_i,\kappa_i)$
the sets $\kappa(\Off \cap \A) \xsubset \R^d$ are measurable.
\item[2.]
A set $\A \xsubset \difM$ is called \textit{small},
if it is contained in the domain $\Off$ a chart $(\Off,\kappa)$.
\end{itemize}
\end{definition}

\subparagraph{}
Each differentiable manifold can be decomposed into mostly countable
many small, measurable and pairwise disjoint subsets.
Dies erlaubt the definition of an integral for differential form
of degree $d$ as integrands.

\begin{definition}
A $d$-form $\omega$ on an \textit{oriented} $d$-dimensional
manifold $\difM$ is called \textit{integrable},
if for a Zerlegung $\{\A_i\}$ of $\difM$ in countable
many small and measurable subsets and a sequence
$\{(\Off_i,\kappa_i)\}_{i \xeof \N}$ of
orientation preserving charts with $\A_i \xsubset \Off_i$ holds:
for each $i \xeof \N$ we have
\[
a_i := \omega \circ \kappa_i^{-1} : \kappa_i(\Off_i) \xto \R
\]
is Lebesgue-integrable in the (open) set $\kappa_i(\A_i)$,
and
\[
\sum_{i=1}^{\infty} \int_{\kappa_i(\A_i)} \abs{a_i(x)} \, dx < \infty
\]
holds.
Then
\[
\int_{\difM} \omega \, := \,
\sum_{i=1}^{\infty} \int_{\kappa_i(\A_i)} \abs{a_i(x)} \, dx
\]
is denoted the Lebesgue integral of $\omega$.
\end{definition}

\section{Manifolds with Boundary}
%------------------------------------------------------------------------------%
Let the Euclidean half-space be defined according to
$\R^d_+ := \{x \xeof \R^d \: x_d \geq 0\}$.

\begin{definition}
A $d$-dimensional differentiable manifold is said to own a \textit{boundary},
if the domains of defintion of all charts in the differentiable structure
are homeomorphic to open subsets of $\R^d_+$.
\end{definition}

\subparagraph{}
Let $f\:\difM \xto \R$ be a scalar field.
Then folgt aus the Rangsatz,
dass for a $a \xeof \R$ the subset $f^{-1}((- \infty, a])$ in
kanonischer Weise a manifold with boundary is.

\paragraph{}
These are examples of manifolds with boundary:

\begin{itemize}[leftmargin=15pt]
\item[1.]
Each manifold is a manifold with boundary $\emptyset$.
\item[2.]
The closed ball $D^d := \{x \xeof \R^d \: \norm{x} \leq 1\}$ has a boundary.
\end{itemize}

\begin{definition}
Let $\difM$ be a $d$-dimensional differentiable manifold with boundary.
\begin{itemize}[leftmargin=15pt]
\item[1.]
A point $p \xeof \difM$ is called a \textit{boundary point} of $\difM$,
if er durch a chart (and damit each chart) $\kappa$
auf a point $\kappa(p)$ with $x_n \eq 0$ abgebildet wird.
\item[2.]
The boundary $\partial \, \difM$ is the set of all Randpunkte of $\difM$.
\item[3.]
The set $\difM \setminus \partial \, \difM$ is called the \textit{interior} of $\difM$.
\end{itemize}
\end{definition}

\begin{theorem}
Let $\difM$ be a $d$-dimensional differentiable manifold.
Then is der boundary a manifold der Dimension $n-1$,
and the Innere of $\difM$ is a $d$-dimensional manifold.
\end{theorem}

\begin{definition}
A differentiable manifold is called \textit{closed},
if it is compact and its boundary is empty.
\end{definition}

\section{The Stokes Theorem}
%------------------------------------------------------------------------------%
The \textit{Satz of Stokes} f\"uhrt the integration of a differential form
on a manifold with boundary
auf the integration ihrer exterior form
auf the boundary of the manifold zur\"uck.

\begin{theorem}
Let $\difM$ be a orientierte $d$-dimensional manifold with boundary and
$\omega \xeof \Omega^{n-1}(\difM)$ a $(n-1)$-form with compact support.
Then holds
\[
\int_{\partial \difM} \omega \, = \, \int_{\difM} \operatorname{d} \omega
\]
\end{theorem}

\subparagraph{}
Im Spezialfall a closed manifold $\difM$ holds f\"ur
each $(n-1)$-form $\omega \xeof \Omega^{n-1} \difM$
\[
\int_{\difM} \operatorname{d} \omega = 0.
\]

\begin{examples}
There are some classical variants
for Stoke's integral theorem being valid on manifolds:
\begin{itemize}[leftmargin=15pt]
\item[1.]
\textit{(Gau\"ss's integral theorem)}
Let $\Omega \subseteq \R^3$ be open
and $\vecX$ be a differentiable vector field auf $\Omega$.
Then for each compacte subset $K \subseteq \Omega$ with sufficiently
smooth boundary we have
\[
\int_K \operatorname{div} \, \vecX \, dV \= \int_{\partial K} \vecX \, dF.
\]
\item[2.]
\textit{(Stokes's integral theorem)}.
Let $\Omega \subseteq \R^3$ be open
and $\vecX$ be a differentiable vector field auf $\Omega$.
Then for each oriented compact subset $K \subseteq \Omega$ with
sufficiently smooth boundary we have
\[
\int_{K} \operatorname{rot} \vecX \, dF \= \int_{\partial K} \vecX \, ds.
\]
\end{itemize}
\end{examples}

\section{deRham Complex}
%------------------------------------------------------------------------------%
The exterior derivative $d$ of a differential form satisfies
the complex condition $d \circ d \eq 0$.
Hence the sequence of the modules $\{\Omega^i \difM\}$ with the
homomorphisms $d_i \:\Omega^i \difM \mapsto \Omega^{i+1} \difM$
forms an aufsteigenden chain complex,
whose homology and cohomology can be studied.

\begin{definition}
Let $\difM$ be a differentiable manifold.
The (aufsteigende) chain complex of the modules $\Omega^s \difM$
with the exterior derivatives as the module-homomorphisms
\[
\{0\}
\stackrel{d}{\longrightarrow} \Omega^0 \difM
\stackrel{d}{\longrightarrow} \Omega^1 \difM
\stackrel{d}{\longrightarrow} \ldots
\stackrel{d}{\longrightarrow} \Omega^n \difM
\stackrel{d}{\longrightarrow}
\{0\}
\]
is called the \textit{deRham complex} of $\difM$.
\end{definition}

\begin{definition}
A differential form $\omega \xeof \Omega^{s+1} \difM$ is called \textit{exact}
rsp.  \textit{integrable},
if it is the exterior derivative of a differential form
$\alpha \xeof \Omega^s \difM$
that is,
$\omega \eq d \, \alpha$.
\end{definition}

\subparagraph{}
Existiert a differential form $\alpha$ with $d \, \alpha \eq \omega$,
then holds wegen $d \circ d \eq 0$,
dass $0 = d (d \, \alpha) \eq d \, \omega$.
This leads to the following definition:

\begin{definition}
A differential form $\omega$ is called \textit{closed},
if $d \, \omega \eq 0$.
\end{definition}

Ist a differential form exakt, so is sie insbesondere closed.
Let $\partial \, \Omega^s \difM \subseteq \Omega^s \difM$
be the submodule of closed differential forms for $s \xeof \N$.

\subparagraph{}
The Tatsache, dass each closed differential form exakt is,
ist gleichbedeutend with exactness of the deRham complex.
Hierzu definiert man the cohomology modules a manifold wie folgt:

\begin{definition}
Let $\difM$ be a $d$-dimensional differentiable manifold.
For $k \eq 1,\ldots,n$ is called
\[
H^k(\difM,\R) := \frac
{\operatorname{Kern}\,(d(\Omega^{k} \difM))}
{\partial \, \Omega^{k-1} \difM}
\]
die $k$-th \textit{cohomology module} of the deRham-complex.
\end{definition}

\subparagraph{}
Let $H^{\star}(\difM,\R)$ be the total sum of the cohomology modules,
\[
H^{\star}(\difM,\R) := \bigoplus_{k=0}^{\infty} H^k(\difM,\R).
\]
Then the exterior product of differential forms induces a
product in $H^{\star}(\difM,\R)$
such that $H^{\star}(\difM,\R)$ sogar a $\R$-Algebra darstellt.
Diese is called the \textit{Kohomologie-Algebra} a manifold $\difM$.

\subparagraph{}
For $d$-dimensional manifolds holds
obviously $H^k(\difM,\R) = 0$ for $k > n$.

The elements aus $\partial \, \Omega^0 \difM$ are the scalar fields $f$ mit
$d \, f = 0$,
that is,
the constant scalar-valued functions on $\difM$.
If $\difM$ is even connected,
then $H^0(\difM,\R)$ is isomorphic to the real line $\R$.

\subparagraph{}
Es sei daran erinnert, dass each differentiable mapping $f\:\difM \xto \difN$
zweier $d$-dimensional manifolds a Folge of Pullbacks
$f^\*_k:\Omega^k \difN \xto \Omega^k \difM \, (k=1,\ldots,n)$
der related differential form modules induziert.
The Pullbacks $f^\*_k$ erzeugen selbst wieder mappings $f_k$
between den deRham-Kohomologie-modules:
\[
f_k: H^k(\difN,\R) \xto H^k(\difM,\R), \quad [x] \mapsto [f^\*(x)].
\]
For homotope mappings erh\"alt man identische Pullbacks:

\begin{theorem}
Let $\difM$ and $\difN$ be $d$-dimensional differentiable manifolds
and $f,g \:\difM \xto \difN$ be differentiable and homotopic mappings.
Then holds $f_k = g_k$ for all $k=1,\ldots,n$.
\end{theorem}

\subparagraph{}
Aus der Tatsache, that a kontrahierbare manifold
is homotopy-equivalent to a einpunktigen set and its
deRham cohomology modules are trivial, folgt:

\begin{theorem}
Let $\difM$ be a kontrahierbare differentiable manifold.
Then is each in $\difM$ closed $s$-form $\omega$ exakt.
\end{theorem}

\textit{Sternf\"ormige} sets are contradictable as well.

\begin{definition}
A subset $\Omega \subseteq \R^d$ is called \textit{star-shaped},
if es a point $s \xeof \Omega$ gibt
such that
for all further points in $\Omega$ the linear Verbindung
ebenfalls to $\Omega$ geh\"ort.\\
A subset $G \subseteq \difM$
of a differentiable manifold $\difM$ is called \textit{star-shaped},
if there is a chart $(\Off,\kappa)$ with $G \subseteq \Off$
such that
$\kappa(G) \xsubset \R^d$ is star-shaped.
\end{definition}

\subparagraph{}
Damit is each on a star-shaped set defined, closed
differential form exact.

\begin{example}
In a sufficiently small neighbourhood of a point $p \xeof \difM$,
in which the differential form $\omega$ is closed,
a star-shaped set can be found.
In this each closed differential form $\omega$ is exact
(the local variant of the Poincare's lemma).
\end{example}

\paragraph{}
The following and central theorem quotes
that for manifolds both cohomology theories are one and the same:

\begin{theorem}
Let $\difM$ be a differentiable manifold.
Then the deRham-cohomology is isomorphic to the singular cohomology
of $\difM$ (considered as a topological space).
\end{theorem}

\paragraph{}
The \textit{Poincare duality} is a basic result on
the structure of the homology and cohomology groups of a manifold.

\begin{theorem}
Let $\difM$ be a compact oriented manifold of dimension $d$.
Then there is an isomorphism $D_k$ between the groups $H_k(\difM)$
and $H^{n-k}(\difM)$ for $k=0,\ldots,n$.
\end{theorem}

\newpage
\thispagestyle{empty}
\fancyhead[C] {\small{\textsc{IV. Lie Groups}}}
%------------------------------------------------------------------------------%
\part{Lie Groups}                                                              %
%------------------------------------------------------------------------------%
An especially intensively studied and understood family of topological groups
are the \textit{Lie groups}.
Its theory is extremely well developed and plays
in many topics of modern mathematics a central role.
Hence the knowledge of their topological properties is of principle interest.

\section{Basic Notations}
%------------------------------------------------------------------------------%
\begin{definition}
A topological group $\lieG$ is called a \textit{Lie group},
if $\lieG$ is a differentiable manifold at the same time 
and the group operation $(g,h) \mapsto g \circ h^{-1}$ is differentiable.
\end{definition}

\subparagraph{}
This definition makes sense,
sinc the product of two manifolds is a manifold again.
It is sufficient to require in the 
definition of a Lie group only the continuity of the group operation.
In 1952 it was proved by \textsc{Montgomery} and \textsc{Zippin}
that continuity of the group operation within a Lie group already implies
differenciability.

\paragraph{}
Examples of Lie groups are:

\begin{itemize}[leftmargin=15pt]
\item[1.]
each finite group is a discrete, $0$-dimensional Lie group;
\item[2.]
$(\Rd,+)$ and $(\C^d,+)$ are $d$-dimensional Lie groups;
\item[3.]
$(\C^{\times} := \C \setminus \{0\}, \cdot)$ is a Lie group;
\item[4.]
$\GL{n,\F}$,
the group of invertible $\cross{n}{n}$-matrices with coefficients
in $\F=\R$ or $\F=\C$;
\item[5.]
finite products of Lie groups are Lie groups again.
Especially the torus $\T^n$,
considered as a multiple product of the group $(S^1,\cdot)$ for $n \xeof \N$,
is a Lie group.
\end{itemize}

\subparagraph{}
Let $\lieG$ be a Lie group.
A subgroup $\lieU \xsubset \lieG$ is called a \textit{Lie-subgroup}
rsp. a \textit{Liescher Normalteiler} of $\lieG$,
if $\lieU$ a (algebraische) rsp. normale subgroup of $\lieG$
and zugleich a submanifold of $\lieG$ is.

\paragraph{}
Lie-subgroups rsp. Lie normal subgroups can be characterized
by the following theorem, which however is difficult to prove:

\begin{theorem}
Let $\lieG$ be a Lie group.
A subgroup (rsp. a normal subgroup) $\lieU \xsubset \lieG$ is
a \textit{Lie} subgroup (rsp. a \textit{Lie} normal subgroup),
iff $\lieU$ is closed in $\lieG$.
\end{theorem}

\subparagraph{}
This theorem provides some more examples of Lie groups.
Especially the orthogonal and unitary subgroups of
classical groups are Lie groups.

\subparagraph{}
Rather similar to topological groups the following assertion holds:
If $\1 \xeof \lieG$ is the unit of a Lie group,
then its connected component $Z_{\1}$ is a Lie normal subgroup in $\lieG$.
Each other connected component of $\lieG$ is diffeomorphic to $Z_{\1}$.
Since $\lieG$ as a manifold satisfies the second countability axiom,
a Lie group owns mostly countable many connected components.

\begin{definition}
Let $\lieG,\lieH$ be Lie groups.
A \textit{Lie homomorphism} between $\lieG$ and $\lieH$
is a \textit{differentiable} group homomorphism $\phi:\lieG \xto \lieH$.
\end{definition}

Entsprechend is a \textit{Lie isomorphism} a group isomorphism
between den (algebraic) groups $\lieG$ and $\lieH$,
which at the same time is a diffeomorphism
between den manifolds $\lieG$ and $\lieH$ is.
Important examples of Lie isomorphisms are the left and right translation
and the conjugation.
The Lie groups and their homomorphisms constitute a category.

\subparagraph{}
The following concept weakens th concept of isomorphy between Lie groups:
two Lie groups $\lieG,\lieH$ are called \textit{locally isomorphic},
if there is a neighbourhood $U_{\1}$ of $\1 \xeof \lieG$,
a neighbourhood $V_{\1}$ of $\1 \xeof \lieH$
and a diffeomorphism $\phi:U_{\1} \xto V_{\1}$
such that
$\phi(g \circ h) \eq \phi(g) \circ \phi(h)$ for all $g,h \xeof U_{\1}$.

\section{Left-Invariant Vector Fields}
%------------------------------------------------------------------------------%
For a Lie group $\lieG$ and for a fest gew\"ahltes Element
$g \xeof \lieG$ sei $L_g:\lieG \xto \lieG$ the Linkstranslation on $\lieG$.
Since $L_g$ is a diffeomorphism,
the differential $d\,L_g$ induces
for each $h \xeof \lieG$ a vector space-isomorphism between den
tangential spaces $T_h\,\lieG$ and $T_{g \circ h}\,\lieG$.
Insbesondere for $h=\1$ erh\"alt man durch the Linkstranslation
einen isomorphism $\phi_{\1}$ between dem
tangential space $T_{\1}\,\lieG$ and the tangent space $T_g\,\lieG$.
Hence the tangent bundle of a Lie group owns the simple
product structure $T\,\lieG \eq \cross{\lieG}{T_{\1}\,\lieG}$.
By this reaseon Lie groups are always orientable.

\subparagraph{}
For a fixed given vector $v \xeof T_{\1}\,\lieG$ is ferner
in each point $g \xeof \lieG$ a vector
$v_g := d\,L_{\phi_g}(v) \xeof T_g\,\lieG$ definiert.
Hence certain vector fields distinguish from others by the following

\begin{definition}
A vector field $V:\lieG \xto \lieG$ is called \textit{left-invariant},
if for all $g \xeof \lieG$ holds: $dL_g(V(e)) = V(g)$.
\end{definition}

\subparagraph{}
If $V,W$ are left-invariant vector fields on a Lie group $\lieG$,
then each linear combination $\alpha V + \beta W$
and the Lie product $[V,W]$ are left-invariant,
too.
Thus the left-invariant vector fields form a Lie algebra $\g$,
which forms a subalgebra of the Lie algebra of all vector fields on $\lieG$.
For the dimension of the Lie algebra we have
\[
\dim \g \= \dim T_{\1} \lieG \= \dim \lieG.
\]

\begin{definition}
Let $\lieG$ be a Lie group.
Then the Lie-algebra $\g$ of left-invariant vector fields is denoted
the the Lie-algebra belonging to $\lieG$.
\end{definition}

\subparagraph{}
A left-invariant vector field $V$ on $\lieG$ is already uniquely determined
by a given vector in the tangential space of the unit element.
Thus there is a bijective mapping $\Phi: T_{\1}\,\lieG \xto \g$.
Defining the Lie-product of two vectors $v_1,v_2 \xeof T_{\1}\,\lieG$
according to
$[v_1,v_2] := [\Phi(v_1),\Phi(v_2)]$,
hence $T_{\1}(\lieG)$ can be seen as a Lie-algebra (being isomorphic to $\g$ ) as well.

\subparagraph{}
Sind $\lieG,\lieH$ Lie groups
and $\phi:\lieG \xto \lieH$ a Lie-homomorphism,
so wird for the related Lie-Algebren $\g,\h$ ein
Lie algebra homomorphism $\phi_L:\g \xto \h$ induziert.
Falls $\lieG,\lieH$ zueinander Lie-isomorphic sind,
so stellt $\phi_L$ an isomorphism of Lie-Algebren dar.

\subparagraph{}
The Zuordnung $\lieG \mapsto \g$ is daher a covariant functor of the
category der Lie groups into the category of Lie algebras and is called
the \textit{Lie functor}.

\paragraph{}
The following table shows an overview of the classical groups and their
Lie-algebras.
Let $M(n,\R)$ rsp. $M(n,\C)$ be the set of real rsp. complex $(n,n)$-matrices.
\vspace{2mm}

\begin{tabular}{|l|c|c|l|}
\hline
Lie group   & dimension    & $\g$        & elements of $\g$                                  \\
\hline
$\GL{n,\R}$ & $n^2$        & $\gl{n,\R}$ & $M(n,\R)$                                         \\
$\SL{n,\R}$ & $n^2$        & $\sl{n,\R}$ & $\{A \xeof M(n,\R)\:         \trace{A} \eq 0\}$   \\
$\OG{n}   $ & $n(n{-}1)/2$ & $\o {n}$    & $\{A \xeof M(n,\R)\:A \eq {-}A^T              \}$   \\
$\SOG{n}  $ & $n(n{-}1)/2$ & $\so{n}$    & $\{A \xeof M(n,\R)\:A \eq {-}A^T, \trace{A} \eq 0\}$\\
$\GL{n,\C}$ &              & $\gl{n,\C}$ & $M(n,\C)$                                         \\
$\SL{n,\C}$ &              & $\sl{n,\C}$ & $\{A \xeof M(n,\C)\:\trace{A} \eq 0\}$            \\
$\UG{n}   $ &              & $\u {n}$    & $\{A \xeof M(n,\C)\:A \eq {-}A^T            \}$     \\
$\SUG{n}  $ &              & $\su{n}$    & $\{A \xeof M(n,\C)\:A \eq {-}A^T, \trace{A} \eq 0\}$\\
\hline
\end{tabular}\\

\vspace{2mm}

\subparagraph{}
The Lie algebra of a Lie group is a useful tool
to study and classify Lie groups.
Restricting to \textit{local} isomorphy,
on obtains the following

\begin{theorem}
Two Lie groups $\lieG,\lieH$ are locally isomorphic
iff
their Lie algebras $\g,\h$ are isomorphic to each another
(in the sense of an isomorphism between Lie algebras).
\end{theorem}

\subparagraph{}
Hence Lie algebras are roughly classifying Lie groups.
The question occurs
whether for a given Lie algebra $\g$ there is a Lie group $\lieG$
such that
$\g$ is the Lie algebra of $\lieG$.
In the finite-dimensional case we obtain the following result:

\begin{theorem}
Let $\g$ be a \textit{finite}-dimensional Lie-Algebra and $[\g]$
be the class of all Lie groups,
whose Lie algebra is isomorphic to $\g$.
Then there is (up to Lie isomorphy) exactly one Lie group $\lieG_0 \xeof [\g]$,
which is \textit{simply connected}.
\end{theorem}

\subparagraph{}
This Lie group is called the
\textit{universal covering group} of the class $[\g]$ and
from the topological point of view is the simplest representative of this class.

\subparagraph{}
Each other Lie group $\lieG \xeof [\g]$ is associated
with the universal covering group $\lieG_0$ as follows:

\begin{theorem}
Let $\g$ be a finite-dimensional Lie-Algebra
and $\lieG_0$ a simply connected Lie group with Lie-Algebra $\g$.
Let $\lieG$ be a further connected Lie group with Lie-Algebra $\g$.
Then there is a mostly countable, discrete Lie normal subgroup $\difN$
of $\lieG_0$,
which lies in the center of $\lieG_0$,
such that $\lieG \= \lieG_0 / \difN$.
\end{theorem}

\subparagraph{}
Hence one obtains all connected Lie groups with given Lie algebra $\g$
by taking the quotient of the universal covering group $\lieG_0$
by its Lie normal subgroups.

\section{The Exponential Map}
%------------------------------------------------------------------------------%
The exponential map forms a connection between the Lie algebra
of a Lie group and the Lie group itself:

\begin{definition}
Let $\lieG$ be a Lie group.
A differentiable mapping $\Phi:\R \xto \lieG$ is called
\textit{Einparameter-subgroup} of $\lieG$,
if $\Phi(t+s) = \Phi(t) \circ \Phi(s)$ for all $s,t \xeof \R$.
\end{definition}

\subparagraph{}
Especially if $\Phi$ is a global flow on $\lieG$.
For a Lie group there is a interessanten connection between
their one-parameter subgroups
and the Lie-algebra of left-invariant vector fields.

\begin{theorem}
Let $\lieG$ be a Lie group, $\g$ the related Lie-Algebra
and $V \xeof \g$ be a left-invariant vector field on $\lieG$.
Then there is a unique one-parameter subgroup $\Phi_V$
such that $V$ the velocity field of $\Phi_V$ is.
On the other hand,
there is to each one-parameter subgroup $\Phi$ exactly one
left-invariant vector field $V$
such that the orbits of $\Phi$ are the integral curves of $V$.
\end{theorem}

\subparagraph{}
Left-invariant vector fields and one-parameter subgroups are
in a one-to-one relationship to each other.
This enables the definition of the \textit{exponential map}:

\begin{definition}
Let $V \xeof \g$ and $\exp:\g \xto \lieG$ be defined according to
$\exp(V) := \Phi_V(1,e)$.
Then $\exp: \g \xto \lieG$ is called the \textit{exponential map} of $\lieG$.
\end{definition}

\begin{example}
For the classical groups $\GL{d,\R}$ the exponential map
is given by the power series
\[
\exp(\matrA) \= \sum_{k=0}^{\infty} \frac{A^k}{k!},
\]
where $\matrA \xeof \sl{d,\R}$ is a $(d,d)$-matrix.
\end{example}

\section{Transformation Groups}
%------------------------------------------------------------------------------%
In many applications a Lie group is operating in a \textit{differentiable} way
on a manifold.

\begin{definition}
Let $\difM$ be a differentiable manifold.
Then a Lie group $\lieG$ is operating \textit{differentiable} on $\difM$,
if the operation $T: \cross{\lieG}{\difM} \xto \difM$ is differentiable.
The triple $(\lieG,T,\difM)$ is called a
\textit{continuous transformation group}.
\end{definition}

\subparagraph{}
Hence the transformation groups are the G-spaces $(\lieG,T,\difM)$,
known from the theory of topological groups,
for the differentiable case.

\subparagraph{}
Let $(\lieG,T,\difM)$ be a transformation group,
where $\lieG$ is operating \textit{transitively} on $\difM$,
that is,
there is exactly one orbit of the operation $T$.
Hence two points $p,q \xeof \difM$ own alsways
conjugated isotropy groups
such that
the homogenneous spaces $\lieG/\lieG_p$ and $\lieG/\lieG_q$
have identical elements.
In this case there is a continuous bijection
$\phi \: \lieG/\lieG_p \xto \difM$.

\subparagraph{}
If $\lieG$ is compact,
then $\phi$ is even a homeomorphism.
In this case the manifold is called a \textit{homogeneous space}.

\paragraph{}
Let us consider the following examples:

\begin{itemize}[leftmargin=15pt]
\item[1.]
Each Lie group $\lieG$ is operating transitively on its own
(for example by left translation).
If $\lieG$ is compact, then $\lieG$ forms a homogeneous space.

\item[2.]
The Lie group $\SOG{d+1}$ of orthogonal $(d+1,d+1)$-matrices with
$\det(A) \eq 1$ is operating transitively on $S^d$.
Since the groups $\SOG{d}$ are compact,
the spheres $S^d$ prove to be homogeneous spaces.
\end{itemize}

\subparagraph{}
If the Lie group $(\R,+)$ is operating on a differentiable manifold $\difM$,
then we are speaking of a \textit{global flow} on $\difM$.

\begin{definition}
Let $\difM$ be a differentiable manifold.
A \textit{global flow} on $\difM$ is a differentiable operation
of the group $(\R,+)$ on $\difM$,
that is,
a differentiable mapping
$\Phi \: \cross{\R}{\difM} \xto \difM$ having the property
$\Phi(0,p) \eq p$ and $\Phi(t+s,p) \eq \Phi(t,\Phi(s,p))$
for all $p \xeof \difM$ and all $t,s \xeof \R$.
\end{definition}

\subparagraph{}
Every orbit of the operation $\Phi$ is called a \textit{trajectory}.
The trajectories on a manifold can be characterized as follows:

\begin{theorem}
Let $\difM$ be a differentiable manifold and
$\Phi$ be a global flow on $\difM$.
Then the trajektories of $\Phi$ are either injektiv and smooth,
periodically smooth or punctate.
\end{theorem}

\begin{definition}
Let $\difM$ be a differentiable manifold and $V$ be a vector field on $\difM$.
Then each curve $\gamma:(a,b) \xto \difM$ mit
$(d\gamma/dt)(t_0) = V(\gamma(t_0))$ for $t_0 \xeof (a,b)$
is called a \textit{integral curve} of $V$ in the point $\gamma(t_0)$.
\end{definition}

\begin{example}
The group $S^1$ operiert on $\R^2$.
The trajectories of the operation are the point $(0,0)$
and all circles around the center point.
\end{example}

\subparagraph{}
A global flow $\Phi$ assignes to each point $p \xeof \difM$ of the orbit space
a tangent vector according to
$(\partial \Phi(t,p) / \partial t) \xeof T_p\,\difM$
In this way a vector field is defined,
called the \textit{velocity field} of the flow.
Hence each trajectory is an integral curve of the velocity field of the flow.

\newpage
\thispagestyle{empty}
\fancyhead[C] {\small{\textsc{V. Connections}}}
%------------------------------------------------------------------------------%
\part{Connections}                                                             %
%------------------------------------------------------------------------------%

\section{Lineare Connections}
%------------------------------------------------------------------------------%
Man fordert on a differentiable manifold the \"Anderung
eines vector field in direction of a tangent vectors \textit{axiomatisch}
with Hilfe of a linear connection,
that is,
a mapping $\nabla$,
die each vector $\vecY$ im tangential space of a point a \"Anderung
in a beliebige direction $\vecX$ zuweist.
Diese mapping soll sowohl for the Argument $\vecY$ as auch for den
direction vector $\vecX$ linear sein.

\begin{definition}
Let $\difM$ be a differentiable manifold
and $T\,\difM$ the tangent bundle on $\difM$.
A \textit{linear connection} on $\difM$ is a mapping
$\nabla: \cross{T\,\difM}{T\,\difM} \xto T\,\difM$
with the following properties:
\begin{itemize}[leftmargin=15pt]
\item[1.]
For each point $p \xeof \difM$ is $\nabla$ a mapping
$\cross{T_p \difM}{T_p \difM} \xto T_p \difM$, $\nabla_X Y \xeof T_p \difM$
for $X,Y \xeof T_p \difM$,
\item[2.]
$\nabla$ is linear {\wrt} the vector $Y$,
\[
\nabla_X (Y_1 + Y_2) = \nabla_X Y_1 + \nabla_X Y_2
\]
\item[3.]
$\nabla_X$ is linear {wrt} the direction vector $X$,
\[
\nabla_{X_1 + X_2} Y = \nabla_{X_1} Y + \nabla_{X_2} Y,
\]
\[
\nabla_{fX} Y = f \cdot \nabla_X Y,
\]
\item[1.]
for a scalar field $f$ and a vector field $\vecY$ holds the product rule
\[
\nabla_X (f \cdot \vecY) = f \cdot \nabla_X \vecY + X(f) \cdot \vecY
\]
\end{itemize}
\end{definition}

\subparagraph{}
If there is linear connection on $\difM$,
the the pair $(\difM,\nabla)$ is called a \textit{linear connected manifold}.

\subparagraph{}
The \textit{\"Ubertragungs forms} and \textit{\"Ubertragungskoeffizienten}
eines linear connection $\nabla$ im point $p \xeof \difM$ erkl\"art man
durch Anwendung of $\nabla$ on the (not necessarily natural)
base vectors $\{e_i\}$,
\[
\nabla_X(e_i) =: \gamma^j_i(X) e_j \quad \mbox{rsp. }
\nabla_{e_k} (e_i) =: \Gamma^j_{ik} e_j.
\]

\begin{example}
Im folgenden wird the parallel transport on a Fl\"ache im $\R^3$
untersucht.
Let for a open set $\openO \xsubset \R^2$ a mapping be defined by
\[
f\: \openO \xto \R, \quad (x,y) \xeof O \mapsto z = f(x,y).
\]
The points $\{(x,y,f(x,y)) \: (x,y) \xeof \openO\}$ are the elements of
two-dimensional manifold $F \xsubset \R^3$ defined by $f$.

\subparagraph{}
Verschiebt man a tangent vector $v \xeof T_p F$ in a benachbarten
point $q \xeof F$, so is $v$ im allgemeinen kein tangent vector mehr in $q$,
since the changes of the base vectors have Anteile in the normal direction
enthalten.
Man kann jedoch $v$ on den tangential space of $q$ projizieren and das
Projektionsergebnis as the \textit{parallel transport} of the Vektors $v$
vom point $p$ in den point $q$ auffassen.\\
Es gen\"ugt obviously, these parallel transport nur for base vectors
$\partial_i$ in $T_p F$ to betrachten.
For linear combinations of the base vectors the parallel transport
is the linear combination of the parallel verschobenen base vectors.

\subparagraph{}
In a point $p \xeof F$ is a nat\"urliche Basis of the tangential space durch
die vectors
\[
\partial_1 := (1,0,f_x)^T, \qquad \partial_2 := (1,0,f_y)^T
\]
gegeben, and for den normal vector in $p$ erh\"alt man
\[
n = \partial_1 \times \partial_2 = \frac {1}{\omega^2} (-f_x, -f_y, 1)^T,
\qquad
\omega^2 = 1 + f_x^2 + f_y^2.
\]
Let $e_z := (0,0,1)^T$, then holds the Identit\"at
\[
e_z = n + \frac{f_x}{\omega^2} \partial_1 + \frac{f_y}{\omega^2} \partial_2.
\]
Let $t = (t_1,t_2)^T \xeof \R^2$ be a direction,
then erh\"alt man for the directional derivatives of $\partial_1$ and $\partial_2$ {wrt} $t$
\[
\frac {d \partial_1}{dt} = f_{xx} \, t_1 \cdot e_z + f_{xy} \, t_2 \cdot e_z,
\]
\[
\frac {d \partial_2}{dt} = f_{yx} \, t_1 \cdot e_z + f_{yy} \, t_2 \cdot e_z.
\]
Durch Einsetzen of $e_z$ in the Gleichungen for $d \partial_1/dt$ and
$d \partial_2/dt$ ergeben sich schlie\"sslich the directional derivatives von
$\partial_1, \partial_2$ as linear combinations of $\partial_1, \partial_2$
and $d$.

\subparagraph{}
In the projection of the directional derivative onto the tangential space
vernachl\"assigt man the Anteile in direction of the normal vector $n$,
and hence the result of the projection is
\[
\frac {d \partial_1}{dt} =
\gamma^1_1(t) \, \partial_1 + \gamma^1_2(t) \, \partial_2, \qquad
\frac {d \partial_2}{dt} =
\gamma^2_1(t) \, \partial_1 + \gamma^2_2(t) \, \partial_2,
\]
where the coefficients $\gamma^i_j$ satisfy the following equations:
\[
\gamma^1_1(t) = \frac{1}{\omega^2}
(f_x \cdot f_{xx}, f_x \cdot f_{xy})^T (t_1,t_2),
\]
\[
\gamma^1_2(t) = \frac{1}{\omega^2}
(f_y \cdot f_{xx}, f_y \cdot f_{xy})^T (t_1,t_2),
\]
\[
\gamma^2_1(t) = \frac{1}{\omega^2}
(f_x \cdot f_{yx}, f_x \cdot f_{yy})^T (t_1,t_2),
\]
\[
\gamma^2_2(t) = \frac{1}{\omega^2}
(f_y \cdot f_{yx}, f_y \cdot f_{yy})^T (t_1,t_2).
\]
The so-called \textit{\"Ubertragungs forms} $\gamma^i_j$ legen somit fest,
wie sich the base vectors beim infinitesimalen \"Ubergang to a neighbourhood
point \"andern.
Obviously they are linear {\wrt} the directional vector $t$ and hence
are linear forms over $\R^2$.

\subparagraph{}
The coefficients $\Gamma^i_{jk}$ of these linear forms are functions of
the first and second partial derivatives of $f$ and are called the
\textit{\"Ubertragungskoeffizienten} {\wrt} the natural base $(dx_1, dx_2)$.

\subparagraph{}
In this example is the comparison of two tangent vectors in benachbarten
tangential spaces geometrically motivated and to the end
given by the partial derivatives of the function $f$.
\end{example}

\section{Covariant Derivatives}
%------------------------------------------------------------------------------%
Mit Hilfe of a linear connection kann a Differentiatonsbegriff f\"ur
vector fields erkl\"art werden.

\begin{definition}
Let $(M,\nabla)$ be a linear zusammenh\"angende manifold
and $\vecX,\vecY$ two zwei vector fields on $\difM$.
Then is called the vector field $\nabla_{\vecX} \vecY$
the \textit{covariant derivative}
of the vector field $\vecY$ {\wrt} the direction $\vecX$.
\end{definition}

\subparagraph{}
Since $\nabla_X Y$ in beiden Argumenten linear is,
stellt $\nabla_X Y$ a (1,1)-tensor field dar.
Let a nat\"urliche Basis in a neighbourhood of a point $p \xeof \difM$ gegeben,
then holds with $Y = Y^i \partial_i$ and $X = X^i \partial_i$ for the covariant
derivative aufgrund the computation rules {\wrt} linear connections
\[
\nabla_X Y =
(\frac{\partial Y^i}{\partial x_j} + \Gamma^i_{jk} \, Y^k) \, X^j \,
\partial_i \otimes dx^j.
\]

\subparagraph{}
For a vorgegebenes vector field $X$ kann the covariant derivative auf
die Gesamtheit of all tensor fields ausgedehnt werden.
For a scalar field $f$ we define in allen points $p \xeof \difM$
\[
(\nabla_X f)(p) := (X(f))(p) \xeof \R.
\]

\subparagraph{}
If two ($r$,$s$)-tensor fields $T_1, T_2$ are given,
then one requires the linearity of the covariant derivative
for the sum of the tensor fields,
\[
\nabla_X (\lambda_1 T_1 \, + \, \lambda_2 T_2) =
\lambda_1 \nabla_X T_1 \, + \, \lambda_2 \nabla_X T_2,
\]
the validness of the following product rule for the product $T_1 \otimes T_2$
\[
 \nabla_X (T_1 \otimes T_2) =
(\nabla_X T_1 \otimes T_2) + (T_1 \otimes \nabla_X T_2),
\]
the Vertauschbarkeit of the covariant derivative with the Spurbildung,
\[
\nabla_X \operatorname{Tr}(T) = \operatorname{Tr}(\nabla_X T),
\]
and the fact
that the covariant derivative
of a $(r,s)$-tensor field provides a tensor field of the type $(r,s+1)$.

\begin{theorem}
Let $(M,\nabla)$ be a linear zusammenh\"angende manifold.
Then l\"a\"sst sich the covariant derivative $\nabla$ for vector fields unter
Ber\"ucksichtigung der obigen Forderungen on genau a Weise to einer
covariant derivative for tensor fields erweitern.
\end{theorem}

Existenz.
The mapping $\nabla: \cross{T\,\difM}{T^r_s \difM} \xto T^r_s \difM$ mit
\[
\nabla(u,T)(\alpha_1,\ldots,v_n) := (\nabla_u T)(\alpha_1,\ldots,v_n) =
\]
\[
\nabla_u T(\alpha_1,\ldots, v_n) -
T(\nabla_u \alpha_1,\ldots,v_n) - \, \ldots \, -
T(\alpha_1,\ldots, \nabla_u v_n).
\]
erf\"ullt all oben genannten Forderungen, stellt also a Erweiterung der
covariant derivative on beliebige tensor fields dar.

\begin{example}
\begin{itemize}[leftmargin=15pt]
\item[1.]
\underline{(r,s) = (0,1)}.
The covariant derivative a linear form $\omega$ on $\difM$
ist the linear form
\[
(\nabla_X \omega)(v) = \nabla_X (\omega(v)) -  \omega(\nabla_X v).
\]
Let a nat\"urliche Basis gegeben mit
$X = X^i \partial_i$, $v = v^i \partial_i$ and $\omega = \omega_i dx^i$,
then holds for the coordinates of the covariant derivative of $\omega$
\[
(\nabla_X \omega)(v) =
(\frac {\partial \omega_i}{\partial x_k} - \Gamma^j_{ik} \omega_j) X^k dx^i
\]
such that
\[
\nabla_X( \omega_i dx^i ) =
(\frac {\partial \omega_i}{\partial x_k} - \Gamma^j_{ik} \omega_j) X^k \, dx^i.
\]
\item[2.]
\underline{(r,s) = (0,2)}.
The covariant derivative of a two-fold covariant tensor field $T$ ist
\[
(\nabla_X T)(v,w) =
\nabla_X T(v,w) - T(\nabla_X v,w) - T(v, \nabla_X w).
\]
Ist $T = t_{ij} dx^i \otimes dx^j$ the representation {wrt} a chart,
so holds
\[
\nabla_X T = ( \frac { \partial T_{ij}}{\partial X_k} -
\Gamma^l_{ik} T_{lj} - \Gamma^l_{jk} T_{il} ) X^k \, dx^i \otimes dx^j.
\]
\item[3.]
\underline{(r,s) = (1,1)}.
The covariant derivative of a gemischten tensor field $T$ zweiter Stufe ist
\[
(\nabla_X T)(\alpha,v) =
\nabla_X T(\alpha,v) - T(\nabla_X \alpha,v) - T(\alpha, \nabla_X v).
\]
Ist $T = t^i_j \partial_i \otimes dx^j$ the representation {wrt} a chart,
so holds
\[
\nabla_X T = ( \frac { \partial T^i_j}{\partial x_k} +
\Gamma^i_{lk} T^l_j - \Gamma^l_{jk} T^i_l ) X^k \, \partial_i \otimes dx^j.
\]
\end{itemize}
\end{example}

\section{Vector-Valued Differential Forms}
%------------------------------------------------------------------------------%
The covariant derivative is a Differentiationsproze\"ss,
der to a directional vector $X$ a tensor field $Y$ im point $p$
einen further tensor zuweist.
L\"a\"sst man den directional vector $X$ im tangential space $T_p \difM$ of a point
variieren,
so erscheint the covariant derivative as linear form on $T_p \difM$
with a \textit{tensor-valued resultat}.
Especially for vector fields erscheint the covariant derivative
as a linear form with vector-valued result.

\begin{definition}
A \textit{vector-valued differential form of degree one}
is a linear mapping
$\alpha: T\,\difM \xto T\,\difM$ with $\alpha(v) \xeof T_p \difM$
for $v \xeof T_p \difM$.
\end{definition}

\subparagraph{}
If a vector field $Y$ is given,
then the covariant derivative generates a vector-valued differential form
of first degree $\operatorname{D}Y$,
defined by
\[
\operatorname{D}Y(v) \eqdef \nabla_v Y \qquad \mbox{ for } v \xeof T_p \difM.
\]
It is called the \textit{absolute differential} of the vector field $Y$.

\subparagraph{}
The absolute differential is damit zun\"achst for vector fields on a
linear connected manifold $\difM$ erkl\"art.
For a scalar field $f$ definiert man the absolute differential
$\operatorname{D}f$ as the already bekannte Differential
$\operatorname{d}f$ of $f$: $\operatorname{D}f := \operatorname{d}f$.

\subparagraph{}
Um the concept of absolute differential on tensor fields auszudehnen,
erkl\"art man a \textit{tensor-valued differential form of first degree}
as follows:

\begin{definition}
A \textit{tensor-valued differential form of first degree} is a linear
mapping
$\alpha: T\,\difM \xto (T^r_s)\,\difM$ with $\alpha(v) \xeof (T^r_s)_p \difM$
for $v \xeof T_p \difM$.
\end{definition}

\subparagraph{}
If a tensor field $T$ on $\difM$ is given,
then the tensor-valued differential form of degree $1$ defined by
\[
\operatorname{D}T(v) := \nabla_v T \qquad \mbox{ for } v \xeof T_p \difM
\]
is called the \textit{absolute differential} of the tensor field $T$.

\subparagraph{}
The process of taking the exterior derivative can be extended to
\textit{vector-valued differential forms of degree $d$}.
The covariant derivative for a vector field hat sich already als
vector-valued differential form of first degree erwiesen,
a notation,
der on h\"oherstufige differential forms erweitert werden kann.

\begin{definition}
Let $\difM$ be a $d$-dimensional differentiable manifold
and $1 \leq s \leq n$.
A \textit{vector-valued differential form of degree $s$} is a
multilinear and anti-symmetric mapping
$\vec{\alpha}: T\,\difM^s \xto T\,\difM$ mit
\[
\vec{\alpha}(v_1,\ldots,v_n) \xeof T_p \difM \quad
\mbox{for } v_1,\ldots,v_s \xeof T_p \difM.
\]
\end{definition}

\begin{remark}
Als Erg\"anzung to theser definition erkl\"art man for $s=0$ eine
vector-valued differerential form of \textit{degree zero} as ein
gew\"ohnliche vector field $X:\difM \xto T\,\difM$.
\end{remark}

\subparagraph{}
If a natural base is given,
then the target vectors of a vector-valued differential form
of degree $s$ can be decomposed {\wrt} the base vectors $\partial_i$,
\[
\vec{\alpha}(v_1,\ldots,v_s) = \alpha^i(v_1,\ldots,v_s) \, \partial_i.
\]
Hierbei are the real-valued and in all arguments linear functions
$\alpha^i$ obviously ordinary $s$-forms.\\
This representation can be generalized,
indem man the natural base vectors $\partial_i$ by $d$ vector fields $X_i$ ersetzt,
which in each point $p \xeof \difM$ are linearly independent.

\subparagraph{}
The sum of two vector-valued $s$-forms $\vec{\alpha}$ and $\vec{\beta}$
and the multiplication a vector-valued $s$-form with a number
is defined according to
\[
(\vec{\alpha} + \vec{\beta}) (v_1,\ldots,v_s) :=
\vec{\alpha}(v_1,\ldots,v_s) + \vec{\beta}(v_1,\ldots,v_s),
\]
\[
(\lambda \cdot \vec{\alpha}) (v_1,\ldots,v_s) :=
\lambda \cdot \vec{\alpha}(v_1,\ldots,v_s),
\]
that is,
the set of vector-valued differential forms with degree $s$ form
a $\R$-vector space.
Ferner definiert man the \textit{outer product} a real-valued
differential form $\alpha$ of degree $s_1$
with a vector-valued differential form $\vec{\beta}$
of degree $s_2$
as vector-valued differential form $\alpha \wedge \vec{\beta}$
of degree $(s_1 + s_2)$,
\[
\alpha \wedge \vec{\beta} (v_1,\ldots,v_{s_1+s_2}) :=
\]
\[
\frac{1}{p! q!} \sum_{\pi} sign(\pi) \,
\alpha(v_{\pi(1)},\ldots,v_{\pi(s_1)} \cdot
\vec{\beta}(v_{\pi(s_1+1)},\ldots,v_{\pi(s_1+s_2)}.
\]
If $\beta^i$ are the coordinates of the $s$-form $\vec{\beta}$,
then $\alpha \wedge \beta^i$ are the coordinates of the products
$\alpha \wedge \vec{\beta}$,
\[
 \alpha \wedge \vec{\beta} = \alpha \wedge (\beta^i \partial_i) =
(\alpha \wedge \beta^i) \partial_i.
\]

\section{Exterior Derivatives of Vector-Valued Forms}
%------------------------------------------------------------------------------%
Um nun the exterior derivative $d$ of real-valued $s$-forms
auf a exterior derivative $d'$ for vector-valued $s$-forms
zu erweitern, fordert man for $d'$

\begin{itemize}[leftmargin=15pt]
\item[1.]
Linearity {wrt} the vector-valued differential forms,
\[
d'(\vec{\alpha_1} + \vec{\alpha_2}) =
d'(\vec{\alpha_1}) + d'(\vec{\alpha_2})
\]
\item[2.]
the validness of the product rule
\[
d'(\alpha \wedge \vec{\beta}) = d \alpha \wedge \vec{\beta} +
(-1)^s \alpha \wedge d'(\vec{\beta}),
\]
for the product of a vector-valued differential of degree $s$
with a real-valued differential form.
\item[3.]
$d'(\alpha) = d\alpha$ for real-valued differential forms $\alpha$.
\end{itemize}

\subparagraph{}
This extension is on differentiable manifolds nicht ohne
weiteres m\"oglich.
But if a manifold is linear connected,
then the extension of the exterior derivative
on vector-valued differential forms can be performed.

\begin{theorem}
Let $(M,\nabla$) be a linear connected manifold.
Then kann the exterior derivative $d$ for real-valued differential forms
auf genau a Weise to a exterior derivative $d_{\nabla}$
for vector-valued differential forms erweitert werden,
die den obigen Forderungen gen\"ugt.
\end{theorem}

\subparagraph{}
The existence of such a mapping $d_{\nabla}$ for a
vector-valued differential form $\vec{\alpha}$ is ensured 
by the Vorschrift
\[
d_{\nabla} \vec{\alpha} :=
(d \alpha_i + (-1)^n \alpha^j \wedge \gamma_j^i) \, \partial_i.
\]
Here the magnitudes $\gamma_j^i$ are denoted 
the \"Ubertragungs forms of the linear connection.

\subparagraph{}
Thus the mapping $d_{\nabla}$ depends on the linear connection $\nabla$
and is uniquely defined by this.
It is the generalization of the exterior derivative of real-valued $s$-forms
to vector-valued differential forms.

\subparagraph{}
The rule {\wrt} vanishing of the second exterior differentials
l\"a\"sst sich for vector-valued differential forms on manifolds mit
linearem connection nicht aufrechterhalten.
Es wird sich zeigen,
dass denjenigen linear connected manifolds,
auf denen the zweite exterior differential of a
vector-valued differential form ausnahmslos verschwindet,
a special state einzur\"aumen ist.

\newpage
\thispagestyle{empty}
\fancyhead[C] {\small{\textsc{VI. Riemannian Spaces}}}
%------------------------------------------------------------------------------%
\part{Riemannian Spaces}                                                       %
%------------------------------------------------------------------------------%

\section{Riemannian Metrics}
%------------------------------------------------------------------------------%
Riemannian manifolds distinguish by the fact
that a symmetric, regular and two-fold covariant
tensor field is defined,
called the \textit{Riemannian metric} rsp.
the \textit{metric fundamental tensor}.

\begin{definition}
Let $\difM$ be a differentiable manifold.
A \textit{Riemannian metric} on $\difM$ is
a two-fold covariant tensor field $g$ on $\difM$ having the properties
\begin{itemize}[leftmargin=15pt]
\item[1.]
$g(v,w) = g(w,v)$ for all $v,w \xeof T_p \difM , \, p \xeof \difM$ (symmetry)
\item[2.]
$g(v,v) > 0$ for all $v \neq 0$ (positive definiteness)
\end{itemize}

Falls on $\difM$ a Riemannian metric existiert,
so is called the Paar $(\difM,g)$ a \textit{Riemannian manifold}.
\end{definition}

\subparagraph{}
A Riemannian metric defines in each point $p \xeof \difM$
an inner product in the related tangential space $T_p \difM $.

\begin{example}
Let $\difM$ be a differentiable manifold and
$\iota:\difM \hookrightarrow \R^d$ be an embedding.
Then is each tangential space $T_p \difM$ a linear subspace von
$T_{\iota(p)} \R^d \cong \R^d$,
and the Einschr\"ankung of the kanonischen Skalarprodukts of the space $\Rd$
auf $T_p \difM$ for each $p \xeof \difM$
induziert a Riemannian metric on $\difM$.
\end{example}

\subparagraph{}
For differentiable manifolds special decompositions
of the unit are of interest.

\begin{definition}
Let $\difM$ be a differentiable manifold.
A decomposition of the unit on $\difM$ is called \textit{differentiable},
if sie an open covering of $\difM$ by charts untergeordnet ist
and the functions $f_i:\difM \xto \R$ of the decomposition are differentiable,
that is,
they are scalar fields on $\difM$.
\end{definition}

\subparagraph{}
Remember that each differentiable manifold is locally compact
and also para-compact.
Hence there is a Zerlegung der Eins on $\difM$,
which can even be chosen 'differentiable':

\begin{theorem}
On each differentiable manifold there is a differentiable decomposition
of the unit.
\end{theorem}

\subparagraph{}
This theorem is decisional for the existence of a Riemannian metric
on manifolds.
The topological properties, the in der definition a manifold
verankert are (Hausdorff-property and second countability axiom),
sind notwendige Voraussetzungen for the existence of differentiable
Zerlegungen der Eins.

\paragraph{}
There is the question,
whether there is a Riemannian metric on a differentiable manifold.
Hierbei spielt the existence of a differentiable Zerlegung der Eins
a decicional role.
Let $\{f_{\alpha}\}_{\alpha \xeof \indI}$ be a solche Zerlegung
and $\Off_{\alpha}$ the domain
of a chart $(\Off_{\alpha},\kappa_{\alpha})$ with
$\operatorname{supp} f_{\alpha} \xsubset \Off_{\alpha}$.
Man definiert for $p \xeof \operatorname{supp} f_{\alpha}$ a Teil-metric
gem\"a\"ss
\[
\bracket{u,v}_{\alpha} :=
\bracket{\operatorname{d}f_{\alpha}(u),\operatorname{d}f_{\alpha}(v)},
\]
where $u,v \xeof T_p \difM$ and
$\operatorname{d}f_{\alpha}(u),\operatorname{d}f_{\alpha}(v) \xeof \R^d$.
Summiert man the Teil-metricen auf,
\[
g(u,v) := \sum_{\alpha \xeof \indI} \bracket{u,v}_{\alpha},
\quad u,v \xeof T_p \difM
\]
so erh\"alt man on these Weise a metric on $\difM$.
Hieraus folgt

\begin{theorem}
Each differentiable manifold   owns a Riemannian metric.
\end{theorem}

\section{Riemannian Connections}
%------------------------------------------------------------------------------%
If on a Riemannian manifold a connection is given,
then the question arises,
whether the parallel transport is consistent
with the measurement of lengths and angles in tangential spaces,
as lengths and angles under parallel transport of vectors along a curve
are not changed.
To emphasize this fact,
a linear connection is denoted \textit{compatible} with the metric tensor,
if the lengths of vectors and the angles between vectors are unchanged
under parallel transport.

\begin{definition}
Let $(\difM,\nabla)$ be a linear connected manifold.
A metric $g$ is called compatible with the connection $\nabla$,
if for each differentiable curve $\gamma$
and for all parallel vector fields $X,Y$ we have
$\bracket{X,Y}(p) \eq \mbox{constant}$
for all points $p$ on $\gamma$.
\end{definition}

\subparagraph{}
This means that the isomporphism given by the parallel transport
\[
\tau_c \: T_p \difM \to T_q \difM
\]
between the tangential spaces of two points $p$ and $q$
connected by a curve $\gamma$ is even an \textit{isometry}.
The following theorem provides a necessary and sufficient criterion for the
compatibility of a linear connection with a Riemannian metric.

\begin{theorem}
Let $\difM$ be a Riemannian manifold.
A connection $\nabla$ is compatible with der metric of $\difM$
\iff for three arbitrary vector fields $X,Y,Z$:
\[
Z(g(X,Y)) \= g(\nabla_Z X,Y) + g(X,\nabla_Z Y).
\]
\end{theorem}

\paragraph{}
The following assertion is denoted the
\textit{main theorem of Riemannian geometry}:

\begin{theorem}
On a Riemannian manifold there is exactly one torsion-free
linear connection being compatible with the metric.
\end{theorem}

\begin{definition}
Let $\difM$ be a Riemannian manifold.
The uniquely defined linear connection $\nabla$ on $\difM$ is called
the \textit{Riemannian connection} rsp. \textit{Levi-Civita connection}.
\end{definition}

\subparagraph{}
The \"Ubertragungskoeffizienten $\Gamma^i_{jk}$ of the Riemannian connection
are uniquely defined by the metric tensor.
If $g_{ij}$ are the coordinates of the metric tensor $g$
{\wrt} a chart $(\Off,\kappa)$,
then the \"Uber\-tragungs\-koeffizien\-ten are given by
\[
\Gamma^i_{jk} \= \frac {1}{2} \, g^{il} \,
(\frac {\partial g_{lj}}{\partial x_k} +
 \frac {\partial g_{lk}}{\partial x_j} -
 \frac {\partial g_{jk}}{\partial x_l} ),
\]
where the magnitudes $g^{il}$ are the coordinates
of the two-fold contravariant tensor $g'$ associated with $g$.

\begin{definition}
The \"Ubertragungskoeffizienten $\Gamma^i_{jk}$ of a Riemannian connection
are called the \textit{Chri\-stof\-fel-Symbole 2.Art}.
The \textit{Christoffel-Symbole 1.Art} are definiert durch
\[
\Gamma_{ijk} \eqdef (\partial_i, \nabla_{\partial_k} \partial_j),
\]
\end{definition}

\subparagraph{}
The Christoffel-Symbole 1.Art are with den Christoffel-Symbolen 2.Art gem\"a\"ss
\[
\Gamma_{ijk} \= g_{il} \, \Gamma^l_{jk}
\]
verkn\"upft.
Aufgrund der Torsionsfreiheit of a Riemannian connection haben the
zugeh\"origen Christoffel-Symbole the following symmetry-properties:
\[
\Gamma_{ijk}  \= \Gamma_{ikj} \quad \mbox{ rsp. }
\Gamma^i_{jk} \= \Gamma^i_{kj}.
\]

\subparagraph{}
Riemannian connections can be characterized as follows:

\begin{theorem}
Let $\difM$ be a Riemannian manifold
and $\nabla$ der Riemannian connection on $\difM$.
Then metric tensor $g$ is \textit{constant},
that is,
for each vector field $\vecX$ on $\difM$ we have $\nabla_{\vecX} g \eq 0$.
\end{theorem}

\paragraph{}
Examples of Riemannian manifolds are:

\begin{itemize}[leftmargin=15pt]
\item[1.]
$\R^d$ is a Riemannian manifold.
\item[2.]
Each surface $(x,y,f(x,y))$ with $(x,y) \xeof 0 \xsubset \R^2$
is a Riemannian manifold.
\item[3.]
The product manifold of two Riemannian manifolds
is also a Riemannian manifold.
\item[4.]
Open subsets of Riemannian manifolds are Riemannian manifolds again.
\end{itemize}

\section{Parallel Vector Fields and Geodaesic Curves}
%------------------------------------------------------------------------------%
If a manifold is linear connected,
then concepts like
parallelity of vector fields, parallelity of vector fields
along curves and geodetic curves can be introduced.

\begin{definition}
Two vector fields $X,Y$ are called \textit{parallel}, if $\nabla_X Y \eq 0$.
A vector field $Y$ is called \textit{parallel} or \textit{constant} on $\difM$,
if $\nabla_X Y \eq 0$ for each vector field $X$.
\end{definition}

\subparagraph{}
For constant vector fields the change of their coordinate
in each direction and in each point is equal zero.
If $\kappa$ is a global chart for $\difM$,
then the coordinates functions of a constant vector field $Y$
are the solutions of the system of partial differential equations
\[
\frac {\partial Y^i}{\partial x_k} + \Gamma^i_{jk} Y^j \= 0,
\qquad (i,k=1,\ldots,n),
\]

\subparagraph{}
In generalization of this concept a \textit{constant tensor field}
is defined as follows:

\begin{definition}
A tensor field $T$ is called \textit{constant} on $\difM$,
if $\nabla_X T \eq 0$ for each vector field $X$.
\end{definition}

\begin{definition}
A \textit{regular curve} on $\difM$ is a injektive differentiable
mapping $\gamma:(a,b) \xsubset \R \xto \difM$.
A vector field $Y$ is called \textit{parallel along a regular curve} $\gamma$,
if $\nabla_{\gamma^{\cdot}} Y \eq 0$ for each point $\gamma(t)$ der curve.
\end{definition}

\subparagraph{}
Ist a vector field $Y$ parallel along a curve,
then in each point
\[
\frac {dY^i(t)}{dt} + \Gamma^i_{jk} Y^j(t) \frac {dx^k(t)}{dt} \= 0,
\qquad (i=1,\ldots,n)
\]
holds.
Here $Y^i(t) := Y^i \circ c(t)$ is the $i$-th coordinate of the function
$Y \circ c:[a,b] \xto T_p(M)$ {wrt} a natural base and
$x_i(t)$ the $i$-te coordinate of the function $(\kappa \circ c)(t)$.
The parallel transport of a vector along a curve h\"angt im
allgemeinen vom Verlauf der curve ab,
that is,
sie f\"uhrt to different results for two curves,
die on verschiedene Weise dieselben points verbinden.
The Aufgabe,
umgekehrt a l\"angs $c$ paralleles vector field zu konstruieren,
indem a vector $v(p_0)$ l\"angs der durch den point
$p_0$ hindurchgehenden curve $c$ parallel verschoben wird,
f\"uhrt on daselbe system of differential equations of first order,
das unter Vorgabe der Anfangswerte and the coordinates
of the starting vector $v(p_0)$ in a neighbourhood of $p$ uniquely solvable is.

\paragraph{}
A vector field besonderer Art l\"angs a curve $\gamma$ is dasjenige
that points in each point in direction of the tangent vector of the curve.

\begin{definition}
A curve $\gamma$ is called \textit{geodetic},
if there is a scalar-valued function $\lambda:[a,b] \xto \R$ with
\[
\nabla_{\gamma^{\cdot}} \gamma^{\cdot} = \lambda \gamma^{\cdot}.
\]
\end{definition}

\subparagraph{}
An equivalent formulation hierfor ist,
dass the curve $\gamma$ der Gleichung $\nabla_{\gamma'} \gamma' = 0$ gen\"ugt,
if $\gamma$ nach der Bogenl\"ange parametrisiert ist.

\subparagraph{}
Ist a chart $(U,\kappa)$ in a neighbourhood of a point gegeben,
so f\"uhrt the Ermittlung a geodetic curve to the ordinary
system of differential equations of second order
\[
\ddot{x}^i + \Gamma^i_{jk} \, \dot{x}^i \dot{x}^k = 0, \qquad (i=1,\ldots,n),
\]
where the functions $x_i(t)$ are the coordinates of $(\kappa \circ c)(t)$.
Durch each point $p \xeof \difM$ geht genau a Geod\"atische with vorgegebenem
tangent vector in $p$.

\begin{definition}
A Riemannian manifold $\difM$ is called (geodetic) \textit{complete},
if to each point $p \xeof \difM$ each geodetic curve durch $p$ for all
parameters  $t \xeof \R$ definiert ist.
\end{definition}

\subparagraph{}
Let $p,q$ be two points of a connected space $\difM$,
so there is a st\"uckweise differentiable Weg $c$ of $p$ nach $q$.
Summiert man the L\"angen der Teilwege auf,
so erh\"alt man the concept of distance between two points.

\begin{definition}
The Abstand $d(p,q)$ zweier points $p,q \xeof \difM$ is the Infimum der L\"angen
aller st\"uckweise differentiable curves between $p$ and $q$.
\end{definition}

\subparagraph{}
Mit theser Abstandsbegriff is $(M,d)$ a metric space.
The dadurch induzierte topology stimmt with der topology of $\difM$
as a differentiable manifold \"uberein.

\paragraph{}
The following Hopf-Rinow theorem macht the concept of completeness of a space interessant.

\begin{theorem}
Let $\difM$ be a connected Riemannian manifold and $p \xeof \difM$.
Then the following assertions are equivalent:
\begin{itemize}[leftmargin=15pt]
\item[1.]
$\exp_p$ is defined on $T_p \difM$.
\item[2.]
Closed and begrenzte subsets of $\difM$ are compact.
\item[3.]
$\difM$ is complete as a metric space.
\item[4.]
$\difM$ is geodetically complete.
\end{itemize}
\end{theorem}

\subparagraph{}
The validness of one of these statements imply
that to each two points $p,q \xeof \difM$ there is a geodetic curve
$\gamma$ with $L(\gamma) \eq d(p,q)$.
Moreover,
each compact Riemannian manifold is complete.

\section{Co-Derivatives and Harmonic Forms}
%------------------------------------------------------------------------------%
Let $\difM$ be a $d$-dimensional Riemannian manifold and for some
$k < n$
$\Omega^{n-k} \difM$ be the space of differential forms of degree $(n-k)$.
For each point $p \xeof \difM$ erh\"alt man den Hodge-operator as a mapping
between den vector spaces der anti-symmetric, purely covariant tensors.
Dies erlaubt es, den Hodge-operator on differential forms von
differentiable manifolds auszudehnen,
indem man punktweise for $p \xeof \difM$ definiert:
\[
\*: \Omega^s \difM \xto \Omega^{n-s} \difM, \quad
\*(\omega)(p) := \omega^\*(p).
\]

Mit Hilfe der mapping $\*$ wird aus der exterior derivative $d$
die \textit{Coableitung} a differential form definiert.

\begin{definition}
The mapping $\delta_k:\Omega^{n-k} \difM \xto \Omega^{n-k-1} \difM$ mit
\[
\delta_k(\omega) := (-1)^k (\* \circ d \circ \*^{-1})(\omega)
\]
is called the \textit{Coableitung}
einer differential form $\omega \xeof \Omega^{n-k} \difM$.
\end{definition}

The mapping $\*$ \"ubersetzt dabei den im Grad der differential forms
aufsteigenden \textit{de Rham-Komplex}
\[
\{0\}
  \stackrel{d}{\longrightarrow} \Omega^1 \difM
  \stackrel{d}{\longrightarrow} \ldots
  \stackrel{d}{\longrightarrow} \Omega^n \difM
  \stackrel{d}{\longrightarrow}
\{0\}
\]
in a dazu \"aquivalenten absteigenden Kettenkomplex
\[
\{0\}
  \stackrel{\delta}{\longrightarrow} \Omega^n \difM
  \stackrel{\delta}{\longrightarrow} \ldots
  \stackrel{\delta}{\longrightarrow} \Omega^1 \difM
  \stackrel{\delta}{\longrightarrow}
\{0\}
\]
with the coderivative as the chain homomorphism.

The Kombination der exterior derivative and der Coableitung liefert for each
$s=0,\ldots,n$ an endomorphism $\Delta$ on the space $\Omega^s \difM$,
der gem\"a\"ss
$\Delta := \delta \circ d + d \circ \delta : \Omega^s \difM \xto \Omega^s \difM$
definiert ist.

\begin{definition}
The operator
$\Delta := \delta \circ d + d \circ \delta :\Omega^s \difM \xto \Omega^s \difM$
is called der \textit{Laplace-Beltrami-operator} on the space $\Omega^s \difM$.
\end{definition}

Von besonderem Interesse are the differential forms $\omega$ mit
$\Delta \omega =0$.

\begin{definition}
A differential form $\omega$ is called \textit{harmonisch} on $\difM$,
if holds: $\Delta \omega =0$.
\end{definition}

The harmonic differential forms of degree $s$ obviously constitute
a vector space, denoted $\cal{H}^s \difM$,
and can be characterized as follows:

\begin{lemma}
A differential form $\omega \xeof \Omega^s \difM$ is harmonic
iff $d \omega \eq 0$ and $\delta \omega \eq 0$.
\end{lemma}

{\bf Satz 7} (Hodge-Theorem).
\begin{theorem}
Each de Rham-cohomology class is represented by exactly one
harmonic differential form.
\end{theorem}

{\bf Satz 8} (Zerlegunsssatz of Hodge).
\begin{theorem}
If $\omega \xeof \Omega^s \difM$ is a $s$-form with $0 < s < n$,
then there are differential forms $\phi \xeof \Omega^{s-1} \difM$ and
$\psi \xeof \Omega^{s+1} \difM$
as well as a harmonic differential form $\chi \xeof \Omega^s \difM$
such that
\[
\omega = d \phi + \delta \psi + \chi,
\]
where the forms $\phi, \psi, \chi$ are uniquely defined by $\omega$.
\end{theorem}

{\bf Satz 9} ((Poincare-Dualit\"at for the de Rham-Kohomologie).
\begin{theorem}
Let $\difM$ be an oriented, closed and $d$-dimensional Riemannian manifold.
Then by the Hodge operator there is an isomorphism
$\iota\:\cal{H}^s \difM \xto \cal{H}^{n-s} \difM$
between the vector spaces of harmonic forms $\cal{H}^s \difM$ and
$\cal{H}^{n-s} \difM$.
\end{theorem}

\section{The Riemannian Curvature Tensor}
%------------------------------------------------------------------------------%

\section{Relationships between Curvature and Topology}
%------------------------------------------------------------------------------%

\begin{comment}
\newpage
\thispagestyle{empty}
\fancyhead[C] {\small{\textsc{VI. Complex Manifolds}}}
%------------------------------------------------------------------------------%
\part{Complex Manifolds}
%------------------------------------------------------------------------------%

\section{{\KAEHLER} Manifolds}
%------------------------------------------------------------------------------%
\section{Riemannian Surfaces}
%------------------------------------------------------------------------------%
\end{comment}

\newpage
\thispagestyle{empty}
\fancyhead[C] {\small{\textsc{VII. Fiber Bundels}}}
%------------------------------------------------------------------------------%
\part{Fiber Bundles}                                                           %
%------------------------------------------------------------------------------%

\section{Generalities on Fiber Bundles}
%------------------------------------------------------------------------------%
Sind $E,B$ topological spaces and is $\pi:E \xto B$ a surjective
continuous mapping,
so is called to each point $p \xeof B$ the Urbild $\pi^{-1}(p)$ a \textit{Faser},
da $E$ sich zusammensetzt gem\"a\"ss
\[
E = \bigsqcup_{p \xeof B} \pi^{-1}(p).
\]
The set $E$ is also the Vereinigung of paarweise disjunkten fibers,
and es gibt a nat\"urliche projection mapping of elements in $E$ auf
den space $B$.

\subparagraph{}
Man spricht of a \textit{fiber bundle},
if all fibers gleich aussehen
sowie der space $E$ locally homeomorphic to $\cross{F}{B}$ ist.

\begin{definition}
Let $E,B,F$ be topological spaces sowie
$\pi:E \xto B$ a continuous surjective mapping.
The Tripel $(E,\pi,B)$ is called a \textit{fiber bundle} with typischer Faser $F$,
if es to each point $p \xeof B$ a neighbourhood $U_p \xsubset B$ and eine
mapping $h_p\: \cross{U_p}{F} \xto \pi^{-1}(U_p)$ gibt
such that $\pi^{-1}(U_p) \cong \cross{U_p}{F}$
(\textit{lokale Trivialisierung} of the bundles).
\end{definition}

Hence the sets $\pi^{-1}(U_p)$ and $\cross{U_p}{F}$ are homeomorphic
to each another.
Lokal sieht a fiber bundle also aus wie a product of topological spaces,
and the concept of fiber bundle forms a generalization of this product.

\subparagraph{}
The space $B$ is called der \textit{base space}
and $E$ is called der \textit{total space} rsp. \textit{bundle space}
of the fiber bundle.
Each local homeomorphism $h$ between $E$ and $\cross{B}{F}$ is called
a \textit{bundle chart} of the fiber bundle.

\paragraph{}
Examples of fiber bundles are:

\begin{itemize}[leftmargin=15pt]
\item[1.]
For each product $(\cross{B}{F},\pi,B)$ of topological spaces $B,F$
and the projection mapping $\pi$ we obtain a fiber bundle
$(\cross{B}{F}i,\pi,B)$ (\textit{product bundle}).
\item[2.]
The triple $(S^d,\pi,\R\,P^n)$ with $d \geq 1$ forms a fiber bundle
with the antipodic mapping $\pi \: x \xto -x$.
\end{itemize}

\paragraph{}
A \textit{cut} in a fiber bundle assigns to each point of the base space
a point of the asscociated fiber.

\begin{definition}
Let $(E,\pi,B)$ be a fiber bundle.
A \textit{cut} into the fiber bundle is a continuous mapping $s \:B \xto E$
with $\pi \circ s \eq id$.
\end{definition}

\begin{example}
Let $(\cross{B}{F},\pi,B)$ be a product bundle
and $f\:B \xto F$ be a continuous mapping.
Then $f$ induces a cut $s:B \xto \cross{B}{F}$ according to
$p \mapsto (p,f(p))$.
\end{example}

\subparagraph{}
Let $(E_1,\pi_1,B)$ and $(E_2,\pi_2,B)$ be two fiber bundle over dem
gemeinsamen base space $B$.
A \textit{bundle-homomorphism} is a fasertreue continuous mapping
$\phi:E_1 \xto E_2$ between den bundle spaces,
so the $\pi_1 = \pi_2 \circ \phi$.
bundle homomorphism f\"uhren also fibers of $E_1$ in fibers
von $E_2$ \"uber.

\begin{example}
Let $(S^n, \pi, \R\,P^n$) be the o.a. fiber bundle mit
$\pi(p) := \{p,-p\}$.
Then are the mappings $Id:S^n \xto S^n$ and $-Id:S^n \xto S^n$
(antipode mapping) the einzigen bundle homomorphism on $S^n$.
\end{example}

\section{Vector Space Bundles}
%------------------------------------------------------------------------------%
The tangent bundle of a differentiable manifold is the
most well-known example of a vector bundle.
However,
vector bundles cannot only be defined over differentiable manifolds,
but also over arbitrary topological spaces:

\begin{definition}
Let $E,B$ be topological spaces,
$\pi:E \xto B$ a continuous surjective mapping and $d \xeof \N$.
The triple $\chi := (E,\pi,B)$ is called a $d$-dimensional
rsp. complex vector bundle,
if the following items hold:

\begin{itemize}[leftmargin=15pt]
\item[1.]
the fiber $\pi^{-1}(p)$ is for each $p \xeof B$ a $d$-dimensional
rsp. complex vector bundle $V$.
\item[2.]
zu each $p \xeof B$ there is a neighbourhood $U_p \xsubset B$
such that $\pi^{-1}(U_p)$ is homeomorphic to $\cross{U}{V}$
\textit{(principle of local triviality)}.
\item[3.]
The vector bundle is called \textit{smooth} rsp. \textit{differentiable},
if $E,B$ are differentiable manifolds
and $\pi$ is a differentiable mapping.
\end{itemize}
\end{definition}

\subparagraph{}
The \textit{total space} $E$ is also locally homeomorphic to $\cross{U}{\R^d}$
for a certain neighbourhood $U \xsubset B$.
Each homeoomorphism $h:\pi^{-1}(U) \xto \cross{U}{\R^d}$ is called a
\textit{bundle chart} of the vector bundle.
The \textit{bundle atlas} of a vector bundle is the totality of all bundle charts.

\begin{examples}
\begin{itemize}[leftmargin=15pt]
\item[1.]
For each topological space $\topX$ and each $n \xeof \N$
is $\cross{\topX}{\Rd}$ a $d$-dimensional vector bundle.
Es is called the \textit{kanonische $d$-dimensional vector bundle} over $\topX$.
Especially for $d=1$ is called a $1$-dimensional vector bundle a
\textit{line bundle}.
\item[2.]
Let $\difM$ be a $d$-dimensional differentiable manifold,
$T\,\difM$ ihr tangential space and $\pi:T\,\difM \xto \difM$ the mapping,
die each tangent vector seinen Fu\"sspunkt zuordnet.
Then is $(T\,\difM,\pi,\difM)$ a es $d$-dimensional vector bundle.
\item[3.]
Let $\difM \xsubset \R^d$ be a differentiable manifold.
The \textit{normal bundle} of $\difM$ is the total space $E := \set{(p,v)}
{p \xeof \difM, v \xeof T_p \difM}$,
together with the projection mapping $\pi: E \xto \difM$.
\end{itemize}
\end{examples}

\begin{definition}
Let $\chi_1 \eq (E_1,\pi_1,B)$ and $\chi_2 \eq (E_2,\pi_2,B)$
be two vector bundles over the base space $B$.
A mapping $\phi:E_1 \xto E_2$ is called ein
\textit{vector bundle homomorphism (isomorphism)}
between $\chi_1$ and $\chi_2$,
if $\pi_2 \circ \phi \eq \pi_1$
and the restriction of $\phi$ on each fiber $\pi^{-1}(p)$
ein vector bundle-homomorphism (isomorphism) ist.
\end{definition}

\subparagraph{}
Homomorphisms between vector bundles f\"uhren also fibers wieder in
fibers \"uber.

\begin{example}
The tangent bundle and the cotangent bundle over of a differentiable
manifold are isomorphic to each another.
\end{example}

\subparagraph{}
Many constructions from the theory of vector spaces can be carried over
to vector bundles,
for example
taking the direct sum,
the tensor product and the constructing the dual space.

\subparagraph{}
Let $\chi_1 = (E_1,\pi_1,B)$ and $\chi_2 = (E_2,\pi_2,B)$ be two
vector bundle over the same base space $B$.
Erkl\"art man for each point $p \xeof B$ the vector space $E_p$ als
\textit{direct sum} der fibers $\pi_1^{-1}(p)$ and $\pi_2^{-1}(p)$,
so stellt the union $E$ of the vector spaces $E_p$
einen topological space $E$ dar
such that the Tripel $\chi_1 + \chi_2 := (E,\pi,B)$
wiederum a vector bundle ist.
The bundle $\chi_1 + \chi_2$ is called the \textit{Whitney-sum} of
the vector bundles $\chi_1$ and $\chi_2$.
Sind $E_1$ and $E_2$ $n_1$- rsp. $n_2$-dimensional,
so is $E$ a $n_1+n_2$-dimensional vector bundle over $B$.

\subparagraph{}
Let $\chi = (E,\pi,B)$ be a vector bundle and $B' \xsubset B$.
Mit $E' := \pi^{-1}(B')$ and $\pi' := \pi \mid_{B'}$ is the Tripel
$\chi' := (E',\pi',B')$ ebenfalls a vector bundle,
called the \textit{restriction} of $\chi$ on $B'$.

\begin{example}
Let $\difM$ be a differentiable manifold,
$(T\difM,\pi,\difM)$ the tangent bundle of $\difM$ and
$\Off \xsubset \difM$ a open subset.
Then is the tangent bundle of $\Off$
die Restriktion of the tangent bundle of $\difM$.
\end{example}

\subparagraph{}
Let $\chi = (E,\pi,B)$ be a vector bundle,
$\topX$ a topological space and $fiii \: X \xto B$ a continuous mapping.
Aufgrund der vorgegebenen mapping $f$ kann man over $\topX$ ein
vector bundle $f^{\*}\chi$ konstruieren,
which is said \textit{to be induced} by the vector bundle $\chi$.
The total space $E^{\*}$ of $f^{\*}\chi$ is the subset
\[
E^{\*} := \set{(x,e)}{x \xeof \topX, e \xeof E} \quad \mbox{with } f(x) = \pi(e).
\]
The projection mapping is defined according to $\pi^{\*}((x,e)) := x$.

If $\chi_1 \eq (E_1,\pi_1,B_1)$ and $\chi_2 \eq (E_2,\pi_2,B_2)$
are vector bundles,
then the cartesian product $\cross{\chi_1}{\chi_2}$ is the vector bundle
$(\cross{E_1}{E_2}, \pi, \cross{B_1}{B_2})$ with $\pi(e_1,e_2) := (\pi_1(e_1),
\pi_2(e_2)$ for all $e_1 \xeof E_1$ and all $e_2 \xeof E_2$.

\begin{example}
For two manifolds $\difM_1,\difM_2$ the cartesian product
of the tangent bundles of $\difM_1$ and $\difM_2$ is isomorphic
to the tangent bundle of the product manifold $\cross{\difM_1}{\difM_2}$.
\end{example}

\subparagraph{}
Let $\chi_1 \eq (E_1,\pi_1,B)$ and $\chi_2 \eq (E_2,\pi_2,B)$
be vector bundles over one and the same base space $B$.
Let for $p \xeof B$ the vector space
$E_p := \tensor{\pi_1^{-1}(p)}{\pi_2^{-1}(p)}$
be the tensor product of the fibers in $p$.
Then the $\tensor{\chi_1}{\chi_2}$ with the total space $\cup E_p$ is also
a vector bundle,
called the \textit{tensor product} of $\chi_1$ and $\chi_2$.

\subparagraph{}
A vector bundle $(E,\pi,B)$ is called \textit{trivial},
if $E$ is homeomorphic to $\cross{B}{\R^d}$.

\subparagraph{}
If especially $B$ is a differentiable manifold
and $E$ the tangential bundle of $B$,
then the concept of Parallelisierbarkeit of a manifold can be derived:

\begin{definition}
A differentiable manifold $\difM$ is called \textit{parallelisierbar},
if the tangent bundle $(T\,\difM,\pi,\difM)$ is trivial.
\end{definition}

\begin{example}
$S^1$ is parallelisierbar, but $S^2$ is not so.
\end{example}

\subparagraph{}
Let $(E,\pi,B)$ be a vector bundle
and $\Gamma(B,E)$ the set of all continuous cuts.
The mapping $s \xeof \Gamma(B,E)$ with $p \mapsto 0 \xeof \pi^{-1}(p)$
is called der \textit{Nullschnitt} of the vector bundle.
Dieser bettet den base space in den total space ein.

\subparagraph{}
If $\chi$ is a $d$-dimensional vector bundle,
then the $d$ cuts $s_1,\ldots,s_n$ are called \textit{nirgends abh\"angig},
if for each $p \xeof B$
the vectors $s_1(b),\ldots,s_n(b)$ are linearly independent.

\begin{theorem}
A $d$-dimensional vector bundle is trivial
iff there are $d$ cuts being nirgends abh\"angig.
\end{theorem}

\begin{definition}
A vector bundle $\chi$ is called \textit{Euclidean},
if there is a continuous mapping $b:E \xto \R$
such that
the restriction of $b$ on each Faser $\pi^{-1}(p), p \xeof B$
forms an inner product.
\end{definition}

\subparagraph{}
If there is an inner product on a vector bundle,
then it is called the \textit{Euclidean metric} of the bundle.
If the base space is a differentiable manifold
and $E$ its tangent bundle,
the the inner product $b$ is called a \textit{Riemannian metric}
and $B$ a \textit{Riemannian manifold}.

\subparagraph{}
If the base space is \textit{paracompact},
then to each vector bundle over $B$ an Euclidean metric can be constructed.

\subparagraph{}
If a vector bundle $\chi$ is Euclidean and trivial,
then there are $d$ cuts $s_1,\ldots,s_d$,
which forms an orthonormal system in each fiber of $\chi$.

\subparagraph{}
If a subvector bundle $\chi$ of the vector bundle $\eta$ is given,
then there is a further vector bundle $\chi'$
such that
$\eta$ is the Whitney sum of $\chi$ and $\chi'$.
In this case the bundle $\chi'$ is called the \textit{orthogonal complement}
of $\chi$.

\section{Principal Fiber Bundles}
%------------------------------------------------------------------------------%
A principal fiber bundle in its most general form is a fiber bundle
$(E,\pi,B)$,
where a topological group $G$ on each fiber transitiv operiert.
Von besonderem Interesse are thejenigen principal fiber bundle,
where $E$ and $B$ are differentiable manifolds and $G$ is a Lie group.

\begin{definition}
Let $(E,\pi,B)$ be a differentiable fiber bundle and $G$ a Lie group.
The triple $(E,\pi,B)$ is called a \textit{principal fiber bundle}
with structure group $G$,
if the following conditions hold:

\begin{itemize}[leftmargin=15pt]
\item[1.]
Let $\set{U_i}{ \xeof \indI}$ be a covering of the base space $B$.
Then there is to each $i \xeof \indI$ a diffeomorphism
$\phi_i: \pi^{-1}(U_i) \xto \cross{U_i}{G}$.
The mapping $\phi$ is called a \textit{bundle chart} of the principal fiber bundle.
\item[2.]
Ist $p \xeof U_i \cap U_j$,
then there is to each $x \xeof \pi^{-1}(p)$
two bundle coordinates $(p,g_i)$ and $(p,g_j)$,
and there is an element $g_{ij}(p) \xeof G$
such that $g_i \eq g_{ij}\cdot g_j$.
\end{itemize}
\end{definition}

\begin{examples}
\begin{itemize}[leftmargin=15pt]
\item[1.]
For each $d$-dimensional manifold $B$ geh\"ort in a natural way
a principal fiber bundle with structure group $\GL(n,\R)$.
Let hierzu $E$ be the set of all $(n+1)$-Tupel $(p,e_1,\ldots,e_n)$,
where $p \xeof B$ a point and the vectors $e_1,\ldots,e_n$ a base of
$T_p B$ darstellen.
Then is called $(E,\pi,B)$ the \textit{Rahmenb\"undel} rsp. \textit{Reper-bundle} der
manifold $B$.
\item[2.]
Each Lie group $G$ with a mapping $\pi:G \xto \{p\}$ in einen
ein-punktigen space $\{p\}$ is a principal fiber bundle with
structure group $G$.
\item[3.]
Let $(E,\pi,B)$ be a smooth $d$-dimensional vector bundle
and $\{U_i\}$ a covering of the base manifold $B$.
If we chose products of the form $\cross{U_i}{\GL(n,\R)}$, $(i \xeof \xeof)$,
then the Verklebung of these products is a principal fiber bundle
with structure group $\GL(n,\R)$.
\end{itemize}
\end{examples}

\subparagraph{}
The structure group is acting from the right on $E$.
\[
\ldots
\]

\section{Connections and Curvature on Principal Fiber Bundles}
%------------------------------------------------------------------------------%

\section{Characteristic Classes}
%------------------------------------------------------------------------------%
The characteristic classes of a vector bundle
over a differentiable manifold $\difM$ are elements of the cohomology algebra
of $\difM$.
Sie messen the Nichttrivialit\"at of the bundles.
Ihre existence stellt Hindernisse for the Vorhandensein of mehreren linear
independent cuts in the vector bundle dar.
The most important characteristic classes are \textit{Chern} classes,
\textit{Pontrjagin} classes, \textit{Euler} classes
and \textit{Stiefel-Whitney} classes.
For product bundles of the form $\cross{B}{\R^d}$
the characteristic classes vanish.

\subparagraph{}
Let $\difM$ be a $d$-dimensional, compact and oriented manifold.
Since $H^n(\difM,\R) \neq 0$,
existiert a $d$-form $\omega \neq 0$ with $d \, \omega \eq 0$ and
\[
\int_M \omega \= 1.
\]
The related cohomology class $[\omega] \xeof H^n(\difM,\R)$ is called the
\textit{orientation class} of $\difM$.

\subparagraph{}
Let $\difM$ be a $d$-dimensional, compact and connected real manifold.
Then there is a $d$-form $\omega$ with $d\,\omega \eq 0$.
The related class $[\omega] \xeof H^n(\difM,\R)$ is called
the \textit{Euler class} $e(T\,M)$ of the tangent bundle of $\difM$.

\paragraph{}
The Euler class owns the following properties.

\begin{theorem}
(Satz of Gau\"ss-Bonnet-Chern).
For each $\omega \xeof e(T\difM)$ holds
\[
\chi(\difM) \= \int_{\difM} \omega.
\]
If the dimension of the manifold is odd, then $e(T\,\difM) \eq 0$.
\end{theorem}

\subparagraph{}
If a manifold owns a non-vanishing Euler class,
then it is not parallelisierbar,
that is,
each continuous vector field on $\difM$  owns at least a root.
Especially $\chi(\difM) \neq 0$ always implies $e(T\,\difM) \eq 0$.

\begin{example}
For $\difM \eq S^d$ we have $\chi(S^n) \eq 1 + (-1)^d$.
\end{example}

\subparagraph{}
By the \textit{Euler class} of an oriented real manifold is meant
the Euler class of its tangential bundle.

\paragraph{}
These characteristic classes have been design nearly at the same time
by \textsc{E.\,Stiefel} and \textsc{H.\,Whitney}.

\begin{definition}
Let $\chi\:E \xto B$ be a real $d$-dimensional vector bundle.
A sequence of cohomology classes $w_i(\chi) \xeof H^i(B,\Z_2)$, $i=0,\ldots$
are called \textit{Stiefel-Whitney classes} of $\chi$ if

\begin{itemize}[leftmargin=15pt]
\item[1.]
If $\eta$ is a further vector bundle over $B$ and $f\:E(\chi) \xto E(\eta)$
a vector bundle homomorphism,
so holds $w_i(\chi) = f^{\*} w_i(\eta)$ for all $i \eq 0,\ldots$.
\item[2.]
If $\chi,\eta$ are vector bundles over $B$,
so holds for the Whitney-sum
\[
w_k(\chi \oplus \eta) = \sum_{i=0}^k w_i(\eta) \cdot w_{k-i}(\eta).
\]
\item[3.]
For the line bundle $\gamma_1^1$ over the circle $S^1$
we have $w_1(\gamma_1^1) \eq 0$.
\end{itemize}
\end{definition}

Especially for $i \eq 0$ is $w_0(\chi) \eq 1 \xeof H^0(B,\Z_2)$.

\subparagraph{}
The concept of Stiefel-Whitney class was introduced in 1935
by \textsc{E.\,Stiefel} and \textsc{H.\,Whitney}.
Stiefel betrachtete characteristic homology classes for the
tangential bundle of a differentiable manifold,
w\"ahrend Whitney the sphere bundle over a simplicial complex
untersuchte.

\subparagraph{}
Aus den properties der Stiefel-Whitney classes folgt,
dass zueinander isomorphe vector bundle identische Stiefel-Whitney classes
besitzen.
Moreover,
if $\epsilon$ is a trivial bundle over $B$,
then $w_i(\chi) \eq 0$ for $i > 0$ holds.
In this case we have for each further bundle $\eta$
\[
w_i(\epsilon \oplus \eta) = w_i(\eta).
\]
If $\chi$ is a $\Rd$-bundle with Euclidean metric having a non-vanishing cut,
then $w_n(\chi) \eq 0$.

\begin{definition}
The \textit{total Stiefel-Whitney class} of a $d$-dimensional real
vector bundle $\chi$ over $B$ is the element
\[
w(\chi) = 1 + w_1(\chi) + w_2(\chi) + \ldots + w_n(\chi)
\]
of the cohomology algebra $H^{\*}(B)$.
\end{definition}

\subparagraph{}
For the total cohomology classes of vector bundles over the base space $B$
the \textit{product theorem} of Whitney holds:
\[
w(\chi \oplus \eta) = w(\chi) \cdot w(\eta).
\]
Hierbei is $\cdot \: \cross{H^{\*}}{H^{\*}} \xto H^{\*}$ the cup product
of the cohomology algebra.

%------------------------------------------------------------------------------%
%                                 BIBLIOGRAPHY                                 %
%------------------------------------------------------------------------------%
\newpage
\thispagestyle{empty}

\begin{thebibliography}{99}
\begin{small}

\begin{comment}
\bibitem{Yo}  K.\,Yosida:
\textit{Functional Analysis}.
Springer, 1980.
\end{comment}

\end{small}
\end{thebibliography}

%------------------------------------------------------------------------------%
%                               END OF DOCUMENT                                %
%------------------------------------------------------------------------------%
\end{document}
